\chapter{Curves, toric varieties, and abelian varieties}
\index{curve!CLD calculus}\index{toric variety!CLD calculus}\index{abelian variety!canonical metric}
\label{chap:cld-examples}

This chapter is a test suite for the local differential theory.  Curves turn
curvature into a graph Laplacian, toric varieties turn it into a real
Monge--Amp\`ere operator, and abelian varieties turn Raynaud uniformization
into a tropical real torus.  In each case one can see precisely which local
and global notions agree.

We retain the conventions
\[
           v=-\log|\cdot|,\qquad \phi_e=-\log\|e\| ,
\]
and we write the Chambert--Loir--Ducros operator as $\dprimeddprime$
throughout.  The glyph $\ddc$ is typed by its setting
(\Cref{sec:appK-collisions}): in \Cref{sec:ch19-thuillier-currents} it is
Thuillier's Laplacian $\mathrm A^0\to\mathrm A^1$, and in
\Cref{sec:ch19-cld-bfj-setup} it is the BFJ model operator of
\Cref{sec:ch20-ddc}.  It is never $\dprimeddprime$.
In particular, $\dprimeddprime[\log|f|]=\delta_{\operatorname{div}(f)}$.
The normalization of Lebesgue measure on $N_{\mathbf R}$ is the one for
which a fundamental parallelepiped of the lattice $N$ has volume one.

\section{Orientation}
\label{sec:ch19-orientation}

\subsection{What is being tested}
\label{sec:ch19-what-tested}

\Crefrange{chap:superforms}{chap:delta-forms} construct a calculus.  A
calculus is only as good as the answers it returns on objects that were
already understood by other means, and this chapter runs that check on the
three families of varieties for which the answers were known before the
calculus existed.

\begin{itemize}
\item On a \emph{smooth proper curve} the classical answer is potential
theory on a metric graph: Laplacians, Green functions, and the slope formula
of \Cref{thm:ch08-slope-formula}.  We verify that
$\dprimeddprime$ of a model function is the graph Laplacian
\eqref{eq:graph-laplacian}, that the Poincar\'e--Lelong formula
\Cref{thm:cld-poincare-lelong} retracts to the slope formula, and that the
two competing local definitions of subharmonicity---Thuillier's
\cite[Def.~3.1.5]{Thu05}, by comparison, and Wanner's current-theoretic
condition $\dprimeddprime[u]\ge0$---define the same class of continuous
functions.  We first keep that current-theoretic class distinct from the
approximation-generated CLD hierarchy of \Cref{def:ch15-hierarchy}, and then
prove their local first-stage equality in
\Cref{thm:p063-current-psh-ordinary-regularization}.
\item On a \emph{toric variety} the classical answer is convex analysis:
semipositive invariant metrics have convex weights on $N_{\mathbf R}$ and their
Monge--Amp\`ere measures are Alexandrov measures, with the factor $n!$
recorded in \eqref{eq:appJ-toric-degree}.  We verify that CLD positivity,
delta-form positivity, and functorial positivity all collapse to convexity,
and that the tropical pushforward of the CLD product is the real
Monge--Amp\`ere operator with no hidden normalization.
\item On an \emph{abelian variety} the classical answer is Raynaud
uniformization together with the theorem of the cube: the canonical metric
is rigid, its curvature current is translation invariant, and its top power
is a Haar measure on the skeleton.  We verify that the canonical tropicalization
is a harmonic tropicalization in the sense of \Cref{sec:harmonic-trop}, so
that the special products of formal curvature currents with harmonic-chart
$G$-forms fall within \Cref{thm:bt-harmonic-compatibility}.  We then compute
the resulting measures; harmonicity alone is not a product theorem for
arbitrary currents.
\end{itemize}

Two further items are not verifications but delimitations, and they are the
reason this chapter closes the part rather than opening it.  The first is
\Cref{sec:ch19-cld-bfj}: the earliest point at which the local calculus of
this part and the global model theory of \Cref{part:bfj} can be compared,
together with a precise statement of the locus on which they agree.  The
second is \Cref{sec:ch19-limitations}: a table of assertions that are true,
or at least meaningful, for model and continuous psh objects and that are
\emph{not} known, or not even meaningful, for arbitrary singular psh
functions.  \Cref{part:bfj,part:comparisons} refer back to that table
whenever a hypothesis looks unnecessarily strong.

\subsection{Prerequisites}
\label{sec:ch19-prerequisites}

The chapter is deliberately thin on new machinery and thick on
cross-reference.  The following are used constantly and are not restated.

\begin{center}
\small
\begin{tabular}{p{.32\textwidth}|p{.25\textwidth}|p{.28\textwidth}}
\textbf{Input}&\textbf{Statement}&\textbf{Used in}\\ \hline
metric graphs, graph Laplacian, Green functions
 &\Cref{sec:metric-graphs}, \eqref{eq:graph-laplacian},
  \eqref{eq:graph-green}
 &\Cref{sec:ch19-curves}\\
slope formula on a semistable model
 &\Cref{thm:ch08-slope-formula}
 &\Cref{thm:ch19-curve-slope-formula}\\
Tate curve and its circle skeleton
 &\Cref{ex:ch08-tate-curve,ex:tate-degeneration}
 &\Cref{ex:ch19-tate-slope,ex:ch19-tate-canonical}\\
toric skeleton, Gauss section, Alexandrov measure
 &\Cref{lem:ch08-toric-gauss-section},
  \eqref{eq:alexandrov-ma}, \Cref{lem:toric-real-total-mass}
 &\Cref{sec:ch19-toric}\\
model metrics and Chambert--Loir measures
 &\Cref{sec:metrics-model,sec:metrics-cl-measure},
  \Cref{thm:cl-model-properties}
 &\Cref{sec:ch19-models,sec:ch19-cld-bfj}\\
Stokes and Green formulae
 &\Cref{thm:cld-global-stokes,thm:ch11-graph-green},
  \Cref{cor:analytic-stokes-green}
 &\Crefrange{sec:ch19-curves}{sec:ch19-abelian}\\
Poincar\'e--Lelong and the vertex atom
 &\Cref{thm:cld-poincare-lelong,lem:ch15-vertex-atom}
 &\Cref{sec:ch19-curves}\\
Bedford--Taylor products and harmonic compatibility
 &\Cref{thm:bt-symmetry-locality,thm:bt-harmonic-compatibility}
 &\Cref{sec:ch19-abelian,sec:ch19-canonical}\\
delta-forms, formal positivity, harmonic tropicalizations
 &\Cref{thm:gk-tropical-frame-positivity,%
        thm:gk-formal-positivity-comparison,thm:harmonic-balancing}
 &\Crefrange{sec:ch19-toric}{sec:ch19-canonical}
\end{tabular}
\end{center}

\subsection{Four normalizations in force}
\label{sec:ch19-normalizations}

Every numerical claim below depends on the following four choices, all fixed
in \Cref{app:normalizations} and none of them optional.

\begin{enumerate}
\item\label{itm:ch19-norm:val} \emph{Valuation versus logarithm.}
$v=-\log|\cdot|$; a metric weight is $\phi_e=-\log\|e\|$; and
$\dprimeddprime[\log|f|]=+\delta_{\divis(f)}$, hence
$\dprimeddprime[v(f)]=-\delta_{\divis(f)}$.  See
\eqref{eq:appJ-cld-pl} and \eqref{eq:appJ-valuation-pl}.
\item\label{itm:ch19-norm:lap} \emph{The graph Laplacian is the sum of
outgoing slopes}, \eqref{eq:graph-laplacian}; a convex kink therefore has
positive mass, and a harmonic function has vanishing outgoing-slope sums.
\item\label{itm:ch19-norm:lattice} \emph{Lebesgue measure on
$N_{\mathbf R}$ is lattice-normalized}: a fundamental parallelepiped
of $N$ has volume one (\Cref{def:ch11-lattice-measure}).
On a strictly semistable curve edge $xy=a$, the length is
$-\log|a|$ and the slope parameter is $-\log|x|$.
The residue-degree weights and the nonreduced-model length convention
are specified in \Cref{sec:ch19-graph-laplacian}; the distinction
from primitive lattice length is recorded in \Cref{rem:lattice-length}.
\item\label{itm:ch19-norm:factorial} \emph{No probability normalization.}
The top wedge of $n$ Chern currents of an ample $L$ has mass $(L^n)$,
which in the toric case is $n!\vol_M(P)$, not $\vol_M(P)$; the factor
$n!$ is \eqref{eq:appJ-toric-degree} and is displayed in every mass check
below.
\end{enumerate}

A reader who reproduces a computation of this chapter and obtains a mass off
by $n!$, by a sign, or by a lattice index has almost certainly departed
from one of \Cref{sec:ch19-normalizations}\,\cref{itm:ch19-norm:val}--\cref{itm:ch19-norm:factorial}; the
mass ledger of \Cref{sec:ch19-mass-ledger} lists every check in the chapter
together with the degree it is checked against.

\subsection{Plan}
\label{sec:ch19-plan}

\Cref{sec:ch19-curves} treats curves: the graph Laplacian, the slope
formula, and a completely worked elliptic curve with multiplicative
reduction.  \Cref{sec:ch19-wanner} proves Wanner's comparison theorem and
then proves unconditional ordinary-smooth monotone regularization on smooth
proper curves.  \Cref{sec:ch19-toric} treats toric
metrics, tropical Monge--Amp\`ere measures, and three worked surfaces.
\Cref{sec:ch19-models} records the model and good-reduction computations in
arbitrary dimension.  \Cref{sec:ch19-abelian} sets up Raynaud
uniformization and proves that the canonical tropicalization is harmonic;
\Cref{sec:ch19-canonical} constructs canonical metrics, computes their
currents and Monge--Amp\`ere measures, and carries the Tate curve and an
abelian surface to explicit numbers.  \Cref{sec:ch19-cld-bfj} proves the
first CLD--BFJ comparison theorem.  \Cref{sec:ch19-limitations} is the
limitations table, \Cref{sec:ch19-ledger} the comparison ledger, and
\Cref{sec:ch19-dictionary} the chapter dictionary.  Four long comparison
proofs are collected in \Cref{app:ch19-comparison-proofs}; their statements
remain here at the points where they enter the examples.

\section{Curves and graph Laplacians}
\label{sec:ch19-curves}

Let $C$ be a smooth proper curve over a complete nontrivially valued
field $K$, and fix a strictly semistable formal model $\mathfrak C$ over
$K^\circ$.  Its skeleton and retraction are denoted by
$\Gamma$ and $\tau:C^{\An}\to\Gamma$.
If a model exists only after extending $K$, these objects initially
belong to the base-changed curve.  We specify below how the formulas
descend when the model is defined over $K$.

\subsection{The geometry behind the calculation}
\label{sec:ch19-curve-geometry}

The skeleton is attached to the chosen model.  Over an algebraically
closed field, its complement is a union of open discs; over a general
field this description is made after scalar extension.
A model function $u=F\circ\tau$ is constant along the branches outside
$\Gamma$.  Its curvature is supported on $\Gamma$, but computing it
against a smooth test function may require a larger skeleton adapted
to that test function.  This distinction is used in the proof below.

The resulting curvature current is a finite signed measure with atoms
given by slopes.  This conclusion concerns the specified model
functions: a general $(1,1)$-current on a curve is a distribution,
and need not be a measure.  Likewise, a continuous subharmonic function
need not factor through one fixed finite skeleton.

\subsection{The graph Laplacian}
\label{sec:ch19-graph-laplacian}

Give an edge coming from a node chart $xy=a$ its annular length
$\ell_e=-\log|a|$.  A coordinate along it is $t=-\log|x|$;
the opposite orientation replaces $t$ by $\ell_e-t$.
Let $k=\widetilde K$.  For a node $q_e$ of a strictly semistable
model, the calibration weight is
\[
                         m_e=[\kappa(q_e):k].
\]
The extension here is separable: in a semistable chart the node
stratum is \'etale over $k$.  After extension to a completed algebraic
closure $\Omega$ of $K$, the edge has $m_e$ geometric lifts, each of
weight one and the same annular length.  Thus $m_e=1$ over an
algebraically closed residue field, or for a rational node; strict
semistability alone does not force that value.
This is the curve case of
\Cref{prop:trop-calibration-base-change}.

For a continuous piecewise affine function $F$ on $\Gamma$, define
\begin{equation}
\label{eq:cld-curve-graph-laplacian}
 \Delta_\Gamma F
   =\sum_{v\in\Gamma}
      \left(\sum_{e\ni v}m_e\,d_eF(v)\right)\delta_v ,
\end{equation}
where $d_eF(v)$ is the outgoing derivative in annular length.
The inner sum is over edge germs, so a loop contributes twice.
After subdivision at all break points, the sum is finite.
A convex kink has positive mass, and the two endpoint contributions
of each affine edge cancel.  Consequently
\begin{equation}
\label{eq:ch19-lap-mass-zero}
                          \Delta_\Gamma F(\Gamma)=0 .
\end{equation}
This is the weighted version of \eqref{eq:graph-laplacian}.
Graph measures will also be viewed as measures on $C^{\An}$ by the
inclusion $\Gamma\hookrightarrow C^{\An}$.

For later use, write $p:C_\Omega^{\An}\to C^{\An}$ and
$q:\Gamma_\Omega\to\Gamma$ for the maps under scalar extension.
The semistable coordinates give
$\tau\circ p=q\circ\tau_\Omega$.
Each lifted edge maps isometrically to its original edge, since
the absolute value on $K$ has not been rescaled.  Counting the
geometric edge germs over a vertex gives
\begin{equation}
\label{eq:ch19-graph-scalar-descent}
                 q_*\Delta_{\Gamma_\Omega}(F\circ q)
                              =\Delta_\Gamma F .
\end{equation}
Indeed, the derivative along a lift of $e$ is $d_eF(v)$, and summing
over all vertices above $v$ counts exactly $m_e$ such lifts.
The assertion also holds at a subdivision point.

\begin{proposition}
\label{prop:skeletal-curvature-curve}
If $u=F\circ\tau$ is a model function, then
\[
                           \dprimeddprime[u]=\Delta_\Gamma F .
\]
If $u=-\log\|1\|_{\mathfrak M}$ is represented by a vertical line
bundle $\mathfrak M$ on $\mathfrak C$, the coefficient at the vertex
$x_Y$ of a component $Y$ is $\deg_k(\mathfrak M|_Y)$.
For a rational model metric on a tensor power, this coefficient is
divided by that tensor power.
\end{proposition}

\begin{proof}
First suppose that $K$ is algebraically closed, so all edge weights
are one.  Fix a smooth test function $\psi$ on $C^{\An}$.
Choose a compact skeleton $\Sigma\supseteq\Gamma$ adapted to $\psi$:
$\psi$ is constant on every component of the complement, and
$\Sigma$ supports its differential forms.
The existence of such a skeleton, and a finite subdivision on which
the restrictions of $\psi$ are smooth in annular coordinates, are
\cite[Lemma~18.2.14(a) and the calculation preceding
Lemma~18.2.17]{CLD25}.  The function $u$ is affine on its subdivided
edges and is constant along every added branch.

The restrictions of an analytic smooth test function satisfy
\[
                         \sum_{e\ni v}d_e\psi(v)=0
\]
at every vertex of this adapted skeleton.  To see the source of the
identity, take a smooth cutoff $\chi$ equal to one near $v$ and
supported away from the other vertices.  Apply analytic Stokes to
$\chi\ddprime\psi$.  Its differential is supported on the same
edge germs; the integral along each edge is the endpoint difference
of $\chi\,d_e\psi$.  The outer endpoints contribute zero, and the
inner endpoints give the negative of the displayed sum.
This is \cite[Lemma~18.2.17]{CLD25}.  It is a restriction on
analytic smooth test functions, not an assertion that an arbitrary
abstract graph is a balanced affine tropical space.

The integral of the smooth top form $\dprimeddprime\psi$ is computed
on $\Sigma$ in the annular coordinates.  Two integrations by parts
on each edge give
\[
\begin{aligned}
 \langle\dprimeddprime[u],\psi\rangle
   &=\int_{C^{\An}}u\,\dprimeddprime\psi
     =\sum_e\int_e u_e(t)\psi_e''(t)\,\rd t\\
   &=\sum_v\left(
        \psi(v)\sum_{e\ni v}d_eu(v)
        -u(v)\sum_{e\ni v}d_e\psi(v)\right)\\
   &=\sum_v\psi(v)\sum_{e\ni v}d_eu(v).
\end{aligned}
\]
The added branches have zero $u$-slope.  Removing them leaves
exactly $\langle\Delta_\Gamma F,\psi|_\Gamma\rangle$.
In particular, only the differentiated current has been shown to
be an atomic graph measure; the argument never replaces $[u]$ by a
current supported on $\Gamma$.

For general $K$, apply this calculation on $C_\Omega$ to
$u_\Omega=p^*u$.  Compatibility of integration with scalar
extension gives
\[
                   p_*[u_\Omega]=[u],\qquad
                   p_*\dprimeddprime[u_\Omega]=\dprimeddprime[u].
\]
For the first identity, test against a smooth top form and use its
associated measure with the continuous coefficient $u$; the second
follows because the current differentials commute with pushforward.
These are the weighted form of
\cite[Propositions~15.3.1 and 16.1.1]{CLD25}.
Together with \eqref{eq:ch19-graph-scalar-descent}, they prove the
claimed measure identity over $K$.

For the coefficient assertion, apply
\Cref{thm:cld-formal-intersection}, that is,
\cite[Theorem~18.5.5]{CLD25}, to $\mathfrak M$.
On a reduced strictly semistable model over $\Omega^\circ$ there
is one point over each generic component, with multiplicity,
value-group index and defect all equal to one.  Its coefficient is
therefore the residual degree.  Under descent, the degrees on the
geometric components above $Y$ sum to $\deg_k(\mathfrak M|_Y)$.
This proves the coefficient formula over $K$ as well.
Linearity gives the tensor-power denominator.
\end{proof}

\begin{remark}
\label{rem:ch19-nonreduced-model-lengths}
For a proper regular snc curve model over a discretely valued field,
write $\mathcal C_s=\sum_Yb_YY$ and $D=\sum_Ya_YY$.
The model-curvature coefficient at $x_Y$ is
$b_Y(D\cdot Y)$ by \Cref{thm:cld-recovers-cl}.
Its expression in the intersection graph keeps the multiplicities.
Put $\lambda=v_K(\pi)$; the vertex value is
$u_Y=\lambda a_Y/b_Y$.  With the intersection lengths
$\ell_e=\lambda/(b_Yb_Z)$ at a node between $Y$ and $Z$ and
weights $m_e=[\kappa(q_e):k]$, one obtains
\[
\begin{aligned}
 \sum_{e\ni Y}m_e\frac{u_Z-u_Y}{\ell_e}
   &=b_Y\sum_{Z\ne Y}
         \left(a_Z-\frac{a_Yb_Z}{b_Y}\right)(Z\cdot Y)\\
   &=b_Y(D\cdot Y).
\end{aligned}
\]
The last equality uses $(\mathcal C_s\cdot Y)=0$.
These are the lengths of \Cref{rem:lattice-length}; replacing them
by a different lattice normalization requires the corresponding
change of weights.  In the regular strictly semistable case all
$b_Y$ are one and the node equation is $xy=\pi$.
\end{remark}

For a formal metric represented by $\mathcal L$ on $L^{\otimes m}$,
the curve formula is
\begin{equation}
\label{eq:ch19-curve-formal-measure}
 c_1(\overline L)
    =\frac1m\sum_Yb_Y\deg_k(\mathcal L|_Y)\delta_{x_Y}.
\end{equation}
Here $b_Y=1$ on the strictly semistable model of the proposition;
the displayed multiplicities apply to relatively normal models
over a discrete valuation ring.  For other formal models the
factors are those of \Cref{thm:cld-formal-intersection}.
The total mass is $\deg_C L$, and positivity is equivalent to
nonnegativity of each displayed residual degree, hence to relative
nefness in these curve-model cases.  This recovers
\eqref{eq:ch08-curve-cl-measure} in its discretely valued setting.

\begin{proposition}
\label{prop:ch19-curve-model-curvature}
Let $\overline L_0$ be a formal metric on $L$ with Chern measure
$\mu_0$ as in \eqref{eq:ch19-curve-formal-measure}, and let
$\varphi=F\circ\tau$ be a model function, $F$ piecewise affine on
$\Gamma$.  Assume that the formal metric and the model function are
represented on a common strictly semistable formal model.
Use its skeleton $\Gamma$, subdivided at the break points of $F$;
in particular, all atoms of $\mu_0$ lie at vertices.  Write $\overline L_\varphi$ for the metric
$\|\cdot\|_\varphi=\rme^{-\varphi}\|\cdot\|_0$.  Then
\[
 c_1(\overline L_\varphi)=\mu_0+\Delta_\Gamma F ,
\]
a signed measure of total mass $\deg_CL$, independent of $\varphi$.  It
is positive if and only if
\[
 \mu_0(\{v\})+\sum_{e\ni v}m_e\,d_eF(v)\ \geq\ 0
 \qquad\text{for every }v\in\Gamma ,
\]
a finite list of linear inequalities in the vertex values of $F$.
\end{proposition}

\begin{proof}
By \eqref{eq:cld-pl-section-sign} applied to a local frame,
$c_1(\overline L_\varphi)=c_1(\overline L_0)+\dprimeddprime[\varphi]$, and
$\dprimeddprime[\varphi]=\Delta_\Gamma F$ by
\Cref{prop:skeletal-curvature-curve}.  The mass statement is
\eqref{eq:ch19-lap-mass-zero}.  Positivity of an atomic measure on $\Gamma$
is nonnegativity of each atom, and inserting vertices at the break points of
$F$ makes the list finite; at a two-valent point interior to an edge of the
model, $\mu_0$ has no mass, and the condition reads that the slope of $F$ is
nondecreasing there, that is, that $F$ is convex across that point.
\end{proof}

\begin{remark}
\label{rem:ch19-curve-finite-test}
For the model metrics above, relative nefness and positivity of the
curvature measure are equivalent to the same residual-degree
inequalities.  The comparison with formal positivity and with the
Chern delta-form, after extension to $\widehat{K^a}$, uses
\Cref{thm:ch17-psh-residual-nef,thm:gk-formal-positivity-comparison};
it is not a consequence of the abstract graph definition alone.
The passage to continuous potentials is considered in
\Cref{sec:ch19-wanner}, with its separate hypotheses.
\end{remark}

\begin{example}
\label{ex:ch19-blowup-residue-degree}
Assume here that $K$ is discretely valued, with $v_K(\pi)=1$.
Blow up a closed point $p$ of the special fibre of
$\mathcal C=\mathbf P^1_{K^\circ}$, where $\kappa(p)/k$ is
separable of degree $d$.  Let $Y$ be the strict transform and
$E\simeq\mathbf P^1_{\kappa(p)}$ the exceptional component.
Both have multiplicity one, and
\[
                          E\cdot Y=d,\qquad E^2=-d .
\]
The model metric of the generically trivial bundle $\mathcal O(E)$
has weight $\varphi_E$, with values $0$ at $x_Y$ and $1$ at $x_E$.
Its skeleton is the unit interval, but the edge has weight $d$.
Thus
\begin{equation}
\label{eq:exceptional-model-laplacian}
                   \dprimeddprime\varphi_E
                           =d\,\delta_{x_Y}-d\,\delta_{x_E}.
\end{equation}
This agrees with the residual intersection degrees and has total
mass zero.  After an unramified extension splitting $\kappa(p)/k$,
there are $d$ exceptional components and $d$ unit-weight pendant
edges; pushing their Laplacians down gives the same two coefficients.
For a rational point, $d=1$, and the calculation is
\Cref{ex:metrics-p1-blowup} and
\eqref{eq:metrics-p1-blowup-laplacian}.
It shows explicitly why a nonrational node cannot simply be given
weight one.
\end{example}

\subsection{The slope formula as a Poincar\'e--Lelong formula}
\label{sec:ch19-slope}

The intersection proof of \Cref{thm:ch08-slope-formula} treated a
regular semistable model over a discrete valuation ring.  We now
compare the analytic divisor measure with the weighted graph
Laplacian.  Identifying that graph Laplacian requires the annular
and residue-slope calculations in addition to Poincar\'e--Lelong.

\begin{theorem}
\label{thm:ch19-curve-slope-formula}
Let $K$ be complete and nontrivially valued, let $C$ be a smooth proper
$K$-curve, and let $\mathfrak C$ be a strictly semistable model with
skeleton $\Gamma$ and retraction $\tau$.  Let $f\in K(C)^\times$,
and assume that every point of $\Supp\divis(f)$ retracts to a vertex
of $\Gamma$.  Put $F=v_\bullet(f)|_\Gamma$.
Regard the divisor as the measure
$\sum_P\ord_P(f)[\kappa(P):K]\delta_P$.
With the edge weights and annular lengths of
\Cref{sec:ch19-graph-laplacian}, one has
\begin{equation}
\label{eq:ch19-slope-from-pl}
\begin{gathered}
 \tau_*\bigl(\dprimeddprime[\log|f|]\bigr)
     =\tau_*\delta_{\divis(f)}
     =\tau_*\divis(f),\\
                    \Delta_\Gamma F=-\tau_*\divis(f).
\end{gathered}
\end{equation}
The function $\log|f|$ is harmonic on
$C^{\An}\setminus\Supp\divis(f)$, and $F$ is affine with integral
slope on every edge of $\Gamma$.
\end{theorem}

\begin{proof}
\emph{The analytic measure.}
By \Cref{thm:cld-poincare-lelong}, the current
$\dprimeddprime[\log|f|]$ is the finite signed divisor measure.
The factors $[\kappa(P):K]$ are those of integration over a
zero-cycle (\Cref{prop:ch14-delta-Y-current}).
Thus pushforward along the continuous retraction is defined by
\[
             \int_\Gamma a\,\rd(\tau_*\mu)
                  =\int_{C^{\An}}(a\circ\tau)\,\rd\mu,
             \qquad a\in C^0(\Gamma),
\]
and gives the first line of \eqref{eq:ch19-slope-from-pl}.
This operation is pushforward of Radon measures; it does not require
$\tau$ to be an analytic morphism.
Off the divisor, $f$ is an analytic unit and $\log|f|$ is
pluriharmonic.

\emph{The graph calculation over an algebraically closed field.}
On an edge annulus there is no zero or pole of $f$.
The Newton-polygon calculation
\cite[Proposition~2.10]{BPR13Structure} gives
$F(t)=c+mt$ with $m\in\mathbf Z$ in the annular coordinate $t$.
This proves affineness and integrality before any vertex mass is
computed.

At a vertex $v$, normalize $f$ by a scalar so that its value at $v$
has absolute value one.  Its reduction $\widetilde f_v$ is a
nonzero rational function on the smooth projective residue curve
at $v$.  The local residue-slope formula
\cite[Theorem~5.15(3)]{BPR13Structure} identifies the outgoing
$F$-slope along a branch with the order of $\widetilde f_v$ at
the corresponding residue point.  The degree of this principal
divisor is zero, so the sum over all branches is zero.

The branches outside $\Gamma$ are open discs.  If such a disc $D$
is attached at $v$, choose a coordinate $z$ with $|z|<R$.
There are only finitely many zeros and poles of $f$ in $D$.
Writing them as $a_P$, with signed orders $n_P$, gives
\[
              f(z)=h(z)\prod_{P\in D}(z-a_P)^{n_P},
\]
where $h$ is an analytic unit on $D$ and $|h|$ is constant.
Near its outer end, $|z-a_P|=|z|$ for all these points.
The outward-from-$v$ parameter is $-\log|z|+\log R$, so the
$F$-slope into $D$ is $\sum_{P\in D}n_P$.
Separating these disc branches from the graph branches in the
degree-zero relation yields
\[
             \sum_{e\ni v}d_eF(v)
                  =-\sum_{\tau(P)=v}\ord_P(f).
\]
There is no mass in an edge interior, by affineness and the
hypothesis on the divisor.  This proves the second line of
\eqref{eq:ch19-slope-from-pl} over an algebraically closed field.
The residue-slope identity is the analytic input in this argument;
the disc calculation accounts explicitly for the divisor mass
that is lost when the discs are retracted.

\emph{Descent.}
For general $K$, pass to the completed algebraic closure $\Omega$,
and use the maps $p$ and $q$ of
\eqref{eq:ch19-graph-scalar-descent}.
The base-changed model is still strictly semistable, and
$F_\Omega=F\circ q$.  The preceding argument applies on every
geometric component.  On one side,
$q_*\Delta_{\Gamma_\Omega}F_\Omega=\Delta_\Gamma F$.
On the other, a closed point $P$ contributes total multiplicity
$[\kappa(P):K]$ after scalar extension: these multiplicities are
the lengths of the finite algebra $\kappa(P)\otimes_K\Omega$.
They include the inseparable multiplicities.  Hence
$p_*\divis(f_\Omega)=\divis(f)$ with the measure convention in the
statement.  The identity
$\tau\circ p=q\circ\tau_\Omega$ proves the required descent.
Since lifted edges have the same lengths, their integral slopes
also descend unchanged.
\end{proof}

\begin{remark}
\label{rem:ch19-slope-two-proofs}
For an oriented edge from $v$ to $w$, the slope is
$\bigl(F(w)-F(v)\bigr)/\ell_e$.
It is a difference of the two normalized vertex values only when
$\ell_e=1$.  For example, on $xy=\pi^r$ with $v_K(\pi)=1$,
the function $v(x)$ has endpoint values $0$ and $r$ and slope
$1$.  The unit-length case in \Cref{thm:ch08-slope-formula} is
$xy=\pi$.  Its intersection proof and the analytic calculation
above agree in their common scope.  The latter uses both
Poincar\'e--Lelong and the local residue-slope formula; neither
affineness nor the vertex identity follows just by naming the
retraction.
\end{remark}

\subsection{An elliptic curve with multiplicative reduction}
\label{sec:ch19-tate-slope}

We now carry a complete example to numbers.  It is the Tate curve of
\Cref{ex:ch08-tate-curve}, with the same numerical data, so that every
intermediate quantity can be compared with \Cref{chap:model-calculations}.

\begin{example}[The slope formula on a Tate curve]
\label{ex:ch19-tate-slope}
\index{Tate curve!slope formula}
Let $K$ be algebraically closed, complete and nontrivially valued, let
$q\in K^\times$ with
\[
                     \ell:=v(q)=3,
\]
and let $E=E_q=\mathbf G_{m,K}^{\An}/q^{\mathbf Z}$ be the Tate curve.  By
\Cref{ex:ch08-tate-curve} its skeleton is the circle
$\Gamma=\mathbf R/3\mathbf Z$ with coordinate $t=v(z)$, the semistable
model has special fibre a cycle of three projective lines meeting in three
nodes, and the vertices are $v_1,v_2,v_3$ at $t=0,1,2$.

\emph{The theta function.}  Put
\[
 \theta(z)=\prod_{n\geq0}\bigl(1-q^nz\bigr)
           \prod_{n\geq1}\bigl(1-q^nz^{-1}\bigr),
\]
a convergent product on $\mathbf G_{m,K}^{\An}$ with
$\divis(\theta)=\sum_{n\in\mathbf Z}(q^n)$, simple zeros, and the
functional equation
\begin{equation}
\label{eq:ch19-theta-functional}
 \theta(qz)=-z^{-1}\theta(z).
\end{equation}
Write $\eta_t\in\mathbf G_{m,K}^{\An}$ for the point of the apartment with
$v(z)=t$, that is, the Gauss point of the circle $|z|=\rme^{-t}$.  Since
the valuation of a convergent Laurent series at $\eta_t$ is the minimum of
the valuations of its monomials,
\[
 v_{\eta_t}(1-q^nz)=\min\{0,\,n\ell+t\},
 \qquad
 v_{\eta_t}(1-q^nz^{-1})=\min\{0,\,n\ell-t\},
\]
whence
\begin{equation}
\label{eq:ch19-theta-valuation}
 V(t):=v_{\eta_t}(\theta)
      =\sum_{n\geq0}\min\{0,n\ell+t\}
       +\sum_{n\geq1}\min\{0,n\ell-t\}.
\end{equation}
Thus $V$ is the concave piecewise affine function on $\mathbf R$ with
\[
 V\equiv0\ \text{on}\ [0,\ell],\qquad
 V(t)=t\ \text{on}\ [-\ell,0],\qquad
 V(t)=\ell-t\ \text{on}\ [\ell,2\ell],
\]
and in general slope $-k$ on $[k\ell,(k+1)\ell]$ for $k\geq0$ and
slope $k+1$ on $[-(k+1)\ell,-k\ell]$.  Its outgoing-slope sum at
$t=k\ell$ is $-1$ for every $k\in\mathbf Z$, so
$\Delta_{\mathbf R}V=-\sum_k\delta_{k\ell}$; with
$\log|\theta|=-V$ this is \eqref{eq:ch19-slope-from-pl} for $\theta$ on
$\mathbf G_m^{\An}$, whose divisor retracts to
$\sum_k\delta_{k\ell}$.  This is the first check.

\emph{A rational function on $E$.}  Fix $a\in K^\times$ with
$\alpha:=v(a)=1$ and set
\[
 f(z)=\frac{\theta(z/a)\,\theta(za)}{\theta(z)^2}.
\]
By \eqref{eq:ch19-theta-functional},
\[
 \frac{\theta(qz/a)\theta(qza)}{\theta(qz)^2}
 =\frac{\bigl(-(z/a)^{-1}\bigr)\bigl(-(za)^{-1}\bigr)}{(-z^{-1})^2}
   \cdot\frac{\theta(z/a)\theta(za)}{\theta(z)^2}
 =f(z),
\]
so $f$ descends to a rational function on $E$, with
\[
 \divis_E(f)=(P_a)+(P_{a^{-1}})-2(O),
\]
where $P_b$ denotes the class of $b\in K^\times$ and $O=P_1$.  The
three points are distinct because $a^{\pm1}\notin q^{\mathbf Z}$.

\emph{The retraction of the divisor.}  We have $\tau(P_b)=v(b)\bmod3$, so
\[
 \tau_*\divis_E(f)=\delta_{t=1}+\delta_{t=2}-2\delta_{t=0}.
\]
Every one of the three points lies over a vertex of $\Gamma$, so the
hypothesis of \Cref{thm:ch19-curve-slope-formula} holds.

\emph{The potential.}  The automorphism $z\mapsto z/b$ of
$\mathbf G_m^{\An}$ carries $\eta_t$ to $\eta_{t-v(b)}$, so
$v_{\eta_t}(\theta(z/a))=V(t-1)$ and
$v_{\eta_t}(\theta(za))=V(t+1)$.  Hence
\[
 F(t)=v_{\eta_t}(f)=V(t-1)+V(t+1)-2V(t).
\]
On the fundamental domain $t\in[0,3]$ we have $V(t)=0$; moreover
$V(t-1)=t-1$ for $t\in[0,1]$ and $V(t-1)=0$ for $t\in[1,3]$, while
$V(t+1)=0$ for $t\in[0,2]$ and $V(t+1)=2-t$ for $t\in[2,3]$.
Therefore
\begin{equation}
\label{eq:ch19-tate-slope-potential}
 F(t)=
 \begin{cases}
  t-1, & 0\leq t\leq1,\\
  0,   & 1\leq t\leq2,\\
  2-t, & 2\leq t\leq3,
 \end{cases}
\end{equation}
with $F(0)=F(3)=-1$, consistent with periodicity.

\emph{The check.}  The slopes of $F$ are $+1$ on $(0,1)$, $0$ on
$(1,2)$, and $-1$ on $(2,3)$.  Outgoing-slope sums:
\[
\begin{array}{lcl}
 t=0: & \ (+1)+(+1)=+2, &\text{(right slope $+1$; left slope $-1$, outgoing $+1$)}\\
 t=1: & \ (-1)+(0)=-1,  &\\
 t=2: & \ (0)+(-1)=-1.  &
\end{array}
\]
Hence
\[
 \Delta_\Gamma F=2\delta_{t=0}-\delta_{t=1}-\delta_{t=2}
   =-\bigl(\delta_{t=1}+\delta_{t=2}-2\delta_{t=0}\bigr)
   =-\tau_*\divis_E(f),
\]
which is exactly \eqref{eq:ch19-slope-from-pl}.  Total mass on both sides is
$0$, as \eqref{eq:ch19-lap-mass-zero} and
\eqref{eq:ch08-principal-degree-zero} require.

\emph{What has been checked.}  Three independent normalizations:
the sign of normalization~\cref{itm:ch19-norm:val}, since the same computation with
$\log|f|$ instead of $v(f)$ reverses both sides; the Laplacian
convention~\cref{itm:ch19-norm:lap}, since a maximum of $F$ would otherwise
produce a positive atom; and the lattice length of the edges, since the
slopes $\pm1$ are computed in the coordinate $t=v(z)$, which is the
lattice arclength of $\Gamma$ by \Cref{prop:ch08-annulus-skeleton}.
\end{example}

\begin{example}[The same computation seen on a model]
\label{ex:ch19-tate-slope-model}
It is instructive to redo \Cref{ex:ch19-tate-slope} by intersection theory,
which is the route of \Cref{thm:ch08-slope-formula}.  On the $3$-cycle
model $\mathfrak C$, write $Y_0,Y_1,Y_2$ for the components at
$t=0,1,2$, so that $(Y_i\cdot Y_{i\pm1})=1$ and $(Y_i^2)=-2$.  The
divisor of $f$ on $\mathfrak C$ is $H+V$ with
\[
 H=\overline{(P_a)}+\overline{(P_{a^{-1}})}-2\overline{(O)},
 \qquad
 V=\sum_ic_iY_i,\quad c_i=F(i)=(-1,0,0),
\]
by \eqref{eq:ch19-tate-slope-potential}.  Then
\[
 (H\cdot Y_0)=-2,\quad (H\cdot Y_1)=1,\quad (H\cdot Y_2)=1,
\]
because the closures of the three rational points meet the components over
which they retract, transversally and once each, while
\[
 (V\cdot Y_0)=c_0(Y_0^2)+c_1+c_2=2,
 \qquad
 (V\cdot Y_1)=c_0+c_1(Y_1^2)+c_2=-1,
\]
and $(V\cdot Y_2)=-1$ by symmetry.  Each of the three sums
$(H\cdot Y_i)+(V\cdot Y_i)$ vanishes, which is
\eqref{eq:ch08-principal-degree-zero}, and the vertical intersection numbers
$(2,-1,-1)$ are precisely the atoms of $\Delta_\Gamma F$ computed above.
The two routes agree atom by atom.
\end{example}

\section{Wanner's comparison theorem}
\label{sec:ch19-wanner}

There are two a priori different local notions on a curve.  Thuillier calls
an upper semicontinuous function subharmonic if it satisfies comparison with
harmonic functions on affinoid domains \cite[Def.~3.1.5]{Thu05}.  In the
original CLD setting, a continuous function is psh if its current
$\dprimeddprime[u]$ is positive.

\subsection{The two definitions}
\label{sec:ch19-wanner-definitions}

The contrast is the one drawn in \Cref{sec:ch15-classical} between
definitions from below and definitions from above, in the single case where
both can be pushed through to a theorem.

\emph{From above.}  Fix a smooth strictly $K$-analytic curve $C$ and an
open $W\subseteq C^{\An}$.  A continuous function $h$ on an affinoid
domain $V\subseteq W$ is \emph{harmonic} if locally it factors through
an adapted semistable skeleton with a piecewise affine profile whose
weighted outgoing-slope sums vanish at interior points; this is the notion of \Cref{sec:ch08-harmonic} and
\Cref{prop:ch08-dirichlet}, transported to $V$ by retraction.  An upper
semicontinuous $u:W\to[-\infty,\infty)$ is
\emph{Thuillier-subharmonic} if for every strictly affinoid $V\Subset W$
and every harmonic $h$ on $V$ with $u\leq h$ on $\partial V$, one
has $u\leq h$ on $V$.  Membership is proved by verifying comparisons;
nothing is approximated, and no extension of $K$ is used.

Both halves of that definition are \cite{Thu05}, and it is worth recording
how they are set up there, because the construction is what makes the class
usable.  The sheaf $\mathcal H_C$ of harmonic germs is built from the
skeleta $S(\mathcal X)$ of simply semistable formal models, with their
retractions $\tau_{\mathcal X}$, by
$\mathrm H(\mathcal X)=\tau_{\mathcal X}^{*}\mathrm H(S(\mathcal
X),\partial\mathcal X_\eta)$
\cite[Def.~2.3.1, Def.~2.3.12, Def.~2.3.16]{Thu05}; that this defines a sheaf
on every smooth strictly $K$-analytic curve, and not merely on the
generic fibre of one model, uses the rigid-analytic semistable reduction
theorem of Bosch and L\"utkebohmert in the form \cite[Thm.~2.3.8]{Thu05}.
The condition displayed above is his \cite[Def.~3.1.5]{Thu05}, and the cone
of subharmonic germs it defines is a sheaf $\mathcal{SH}_C$ by
\cite[Cor.~3.1.13]{Thu05}.  The class is nonempty for the expected reason:
$\log|f|$ is harmonic when $f$ is invertible \cite[Prop.~2.3.20]{Thu05}
and subharmonic when $f$ is invertible on a dense open
\cite[Prop.~3.1.6]{Thu05}.  \Cref{rem:ch08-harmonic-vs-log} records that
$\mathcal H_C$ is strictly larger than the sheaf generated by the
$\log|f|$, and identifies the obstruction.

\emph{From the current.}  Wanner calls a continuous function $u$ on $W$
\emph{psh} when $\dprimeddprime[u]$ is a positive current.  This is not the
definition of CLD-psh adopted in \Cref{def:ch15-hierarchy}: membership in
that hierarchy implies positivity of the current by
\Cref{prop:cld-psh-positive-current}, but the converse is a separate
regularization question.  The terminology must therefore be kept explicit
in the comparison theorem.

\subsection{The theorem}
\label{sec:ch19-wanner-theorem}

\begin{theorem}\label{thm:wanner-curve-comparison}
Let $K$ be algebraically closed, complete, and nontrivially valued, let
$C$ be a smooth proper algebraic curve, and let $W\subset C^{\mathrm{an}}$
be open.  For a continuous function $u:W\to\mathbf R$,
\[
                  \dprimeddprime[u]\geq0
        \quad\Longleftrightarrow\quad
                  u\text{ is Thuillier-subharmonic}.
\]
Consequently every continuous CLD-psh function on $W$ is
Thuillier-subharmonic.
\end{theorem}
\begin{proof}
This is imported from \cite[Theorem~4.7 and Corollary~4.8]{Wanner19}; no
proof is claimed here.  In that source, ``psh'' is defined by positivity of
the current $\dprimeddprime[u]$.  The forward implication in the final
sentence is \Cref{prop:cld-psh-positive-current} followed by the imported
equivalence.  It is deliberately not reversed: \cite{Wanner19} does not
identify current positivity with the approximation-generated hierarchy of
\Cref{def:ch15-hierarchy} without a separate regularization input.  That
input is supplied later by
\Cref{thm:p063-current-psh-ordinary-regularization}.
\end{proof}

\begin{corollary}
\label{cor:ch19-continuous-zero-current-harmonic}
Under the hypotheses of \Cref{thm:wanner-curve-comparison}, let
$h:W\to\mathbf R$ be continuous.  If
\[
                         \dprimeddprime[h]=0,
\]
then $h$ is Thuillier harmonic, hence harmonic in the sense of
Gubler--Jell--Rabinoff and classically pluriharmonic in the sense of
Gubler--Rabinoff.
\end{corollary}

\begin{proof}
The current equality makes both $\dprimeddprime[h]$ and
$\dprimeddprime[-h]$ positive.  By
\Cref{thm:wanner-curve-comparison}, both $h$ and $-h$ are
Thuillier-subharmonic, so $h$ is Thuillier harmonic.  On a rig-smooth curve,
Thuillier harmonicity agrees with Gubler--Jell--Rabinoff harmonicity by
\cite[Proposition~15.7]{GJR26}.  In particular $h$ is PL.  Both $h$ and
$-h$ are semipositive, and \Cref{thm:gr-classical-pl} therefore makes both
classically psh.  Thus $h$ is classically pluriharmonic.
\end{proof}

% P21.3 residual-Picard curve regularity, 11 September 2026.
\begin{proposition}[Residual Picard criterion for ordinary smooth harmonic germs]
\label{prop:p213-residual-picard-smooth-curve}
Let \(K\) be algebraically closed, complete, and nontrivially valued, let
\(C\) be a smooth proper algebraic \(K\)-curve, and let
\(W\subset C^{\mathrm{an}}\) be open.  Suppose that
\(h\in C^0(W,\mathbf R)\) satisfies
\[
                         \dprimeddprime[h]=0.
\]
For a type~\(2\) point \(x\in W\), let \(C_x\) be its smooth projective
residue curve and let
\[
 \operatorname{ob}_x(h)\in
 \operatorname{Pic}^0(C_x)\otimes_{\mathbf Z}\mathbf R
\]
be the residual class of the harmonic germ of \(h\) at \(x\).  Then the
following are equivalent:
\begin{enumerate}
\item \(h\) is ordinary CLD-smooth near \(x\) and
      \(\dprimeddprime h=0\) there as a smooth form;
\item \(\operatorname{ob}_x(h)=0\).
\end{enumerate}
At points of types \(1\), \(3\), and \(4\), the germ of \(h\) is
automatically ordinary CLD-smooth with zero smooth Hessian.

Consequently, if
\begin{equation}
\label{eq:p213-residual-picard-vanishing}
       \operatorname{Pic}^0(C_x)\otimes_{\mathbf Z}\mathbf R=0
       \qquad\text{for every type~\(2\) point }x\in W,
\end{equation}
then \(h\) is ordinary CLD-smooth on \(W\) and
\(\dprimeddprime h=0\) as a smooth form.  The hypothesis
\eqref{eq:p213-residual-picard-vanishing} holds, in particular, when the
residue field \(\widetilde K\) is algebraic over a finite field, or when
\(C\) is a Mumford curve.
\end{proposition}


The proof is given in \Cref{sec:appN-residual-picard}.

\subsection{What the theorem does and does not give}
\label{sec:ch19-wanner-scope}

\begin{corollary}
\label{cor:ch19-curve-three-classes}
Under the hypotheses of \Cref{thm:wanner-curve-comparison}, every continuous
CLD-psh function $u$ on $W\subseteq C^{\An}$ is
Thuillier-subharmonic and classically psh in the sense of
\Cref{sec:ch15-classical}.  In particular, the maximum principle of
\Cref{cor:ch15-cld-curve-max} is available for every continuous CLD-psh
function on a curve.  The converse implication from either of the two larger
classes requires qualification: the Thuillier/current-positive converse is
\Cref{thm:p063-current-psh-ordinary-regularization}, whereas no converse from
the full classically psh class is asserted.
\end{corollary}

\begin{proof}
CLD-psh implies positivity of $\dprimeddprime[u]$ by
\Cref{prop:cld-psh-positive-current}, hence Thuillier-subharmonic by
\Cref{thm:wanner-curve-comparison}; it is classically psh independently by
\Cref{thm:cld-to-classical-psh}.  The maximum principle is the one available
for the curve-test class.
\end{proof}

Wanner's published ordinary-smooth regularization theorem has a smaller
scope: it applies when the residue field is algebraic over a finite field,
or when $C=\mathbf P^1$, or when $C$ is a Mumford curve
\cite[Corollary~5.4]{Wanner19}.  The unrestricted ordinary-smooth statement
is proved in \Cref{thm:p063-current-psh-ordinary-regularization} below.  The
weakly smooth theory gives a different regularization on rig-smooth curves
by \cite[Proposition~15.11]{GJR26}.

\begin{warning}{Three statements that look alike}
On a curve, the following are three different assertions.
\begin{enumerate}
\item A continuous function has positive $\dprimeddprime$ if and only if
 it is Thuillier-subharmonic (\Cref{thm:wanner-curve-comparison}).
\item A continuous CLD-psh function has positive $\dprimeddprime$
 (\Cref{prop:cld-psh-positive-current}).
\item A continuous current-psh function belongs locally to the first
 decreasing-limit stage of the ordinary CLD hierarchy.  This is
 \Cref{thm:p063-current-psh-ordinary-regularization}; Wanner had previously
 proved it under the three hypotheses of \cite[Corollary~5.4]{Wanner19}.
\end{enumerate}
The first statement compares two definitions from above; the third is the
additional approximation statement and must not be folded into it.
\end{warning}

\begin{proposition}
\label{prop:ch19-cld-smooth-harmonic-nonregularization}
There exist a complete algebraically closed nontrivially valued
non-Archimedean field $K$, a smooth proper connected $K$-curve $C$, a
type~$2$ point $P\in C^{\mathrm{an}}$, a relatively compact connected open
neighbourhood $U$ of $P$, and an integral PL harmonic function
$h:U\to\mathbf R$ with the following properties.
\begin{enumerate}
\item The function $h$ is weakly smooth and
\[
                         \dprimeddprime h=0
\]
as a weakly smooth form.
\item There is no net $(h_i)$ of ordinary CLD-smooth harmonic functions on
$U$---that is, $\dprimeddprime h_i=0$ as an ordinary smooth form---which
converges locally uniformly to $h$.  In particular, no decreasing net of
ordinary CLD-smooth harmonic functions on $U$ converges pointwise to $h$.
\end{enumerate}
This does not disprove decreasing regularization by ordinary CLD-smooth psh
functions: such approximants need not have zero curvature.
\end{proposition}

\begin{proof}
We use the curve constructed in \cite[\S17]{GJR26}.  The ground field $K$ is
complete, algebraically closed, and nontrivially valued, with algebraically
closed residue field not algebraic over a finite field.  The metrized complex
in that construction has a theta graph with a vertex $P$.  The residual
curve at $P$ is an elliptic curve $E$; the three graph directions $a,b,c$ at
$P$ correspond to distinct points $x,y,z\in E(\widetilde K)$, where $z$ is
the identity and $x$ is non-torsion.  The realization theorem used there
gives a smooth proper connected genus-$3$ curve $C/K$ with this metrized
complex.  Write $\tau:C^{\mathrm{an}}\to\Theta$ for the retraction onto its
theta skeleton.

The proof of \cite[Proposition~17.4]{GJR26} subdivides the theta skeleton to
$\Theta'$ and constructs a piecewise $\mathbf Z$-linear function
$F:\Theta'\to\mathbf R$, affine on every edge, whose three outgoing slopes at
$P$ are a permutation of $1,-1,0$.  It also constructs a
vertical Cartier divisor $D$ on the corresponding strictly semistable model
such that
\[
                         H:=F\circ\tau=-\log\lVert1\rVert_{\mathcal O(D)}
\]
on $C^{\mathrm{an}}$.  Thus $H$ is a globally defined integral PL function;
we do not assert that it is harmonic away from the star chosen below.  The
same proof computes the restriction of $\mathcal O(D)$ to the residual
elliptic curve as the class of $[z]-[x]$.

Choose type~$2$ points $Q_a,Q_b,Q_c$ at value-group-rational positions,
sufficiently close to $P$ that each $[P,Q_e]$ contains no vertex of
$\Theta'$ other than $P$, and put
\[
              \Sigma_0=\bigcup_{e\in\{a,b,c\}}[P,Q_e],
              \qquad Y_0=\tau^{-1}(\Sigma_0).
\]
The restriction $F|_{\Sigma_0}$ is affine on each arm, and its three outgoing
slopes sum to $1-1+0=0$ at the sole interior vertex $P$; no balancing
condition is imposed at the boundary points $Q_e$.  Therefore
\cite[Proposition~15.9]{GJR26} makes $H|_{Y_0}$ integral harmonic.  The
domain $Y_0$ is compact and strictly affinoid, with analytic boundary
$\{Q_a,Q_b,Q_c\}$.  Put
\begin{equation}
\label{eq:ch19-harmonic-reg-open-star}
                         U=Y_0\setminus\partial Y_0,
\end{equation}
and set $h=H|_U$.  Then $U$ is a relatively compact connected open
neighbourhood of $P$, and $h$ is integral PL and harmonic.

Its residual slope divisor $D_P(h)$ is supported on
\begin{equation}
\label{eq:ch19-harmonic-reg-support}
                         S=\{x,y,z\}\subset E(\widetilde K),
\end{equation}
and its class is
\begin{equation}
\label{eq:ch19-harmonic-reg-nonzero-class}
 [D_P(h)]=[z]-[x]\neq0
       \quad\text{in}\quad
 \operatorname{Pic}^0(E)\otimes_{\mathbf Z}\mathbf R.
\end{equation}
Indeed, $[z]-[x]$ has infinite order by the choice of $x$.  Its infinite
cyclic subgroup remains injective after tensoring with $\mathbf R$, because
$\mathbf R$ is torsion-free, hence flat, over $\mathbf Z$.  By
\cite[Proposition~10.17]{GJR26}, $h$ is weakly smooth and
$\dprimeddprime h=0$.

Choose type~$2$ points $R_e\in(P,Q_e)$ at value-group-rational positions, put
\[
              \Sigma_Y=\bigcup_{e\in\{a,b,c\}}[P,R_e],
              \qquad Y=\tau^{-1}(\Sigma_Y)\Subset U,
\]
and subdivide the stable model at $R_a,R_b,R_c$.  The corresponding formal
open gives a strictly semistable formal model $\mathfrak Y$ of the strict
affinoid domain $Y$, whose skeleton is $\Sigma_Y$ and whose boundary is
$\{R_a,R_b,R_c\}$.  Let $g$ be harmonic on $U$ in the sense of
Gubler--Jell--Rabinoff.  By \cite[Proposition~15.7]{GJR26}, $g|_Y$ is
harmonic in Thuillier's sense.  Thuillier's strict-affinoid skeletal
description \cite[Definitions~2.3.1, 2.3.12, and 2.3.16]{Thu05} places
$g|_Y$ in $H(Y)$.  The independence of this space from the chosen strictly
semistable model \cite[Proposition~2.3.3]{Thu05} gives
$g|_Y=G\circ\tau_{\mathfrak Y}$, where $G$ is piecewise affine on
$S(\mathfrak Y)$ and has zero sum of outgoing slopes away from $\partial Y$.
At every subdivision point in the interior of an arm, the two one-sided
slopes therefore agree.  Thus $G$ is affine with constant slope on each
$[P,R_e]$.  Since $g$ factors through this three-armed retraction, all other
tangent directions at $P$ have slope zero.  Consequently
\begin{equation}
\label{eq:ch19-harmonic-reg-fixed-support}
       D_P(g)\in
       V_S:=\operatorname{Div}^0_S(E)\otimes_{\mathbf Z}\mathbf R,
\end{equation}
and the three coefficients of $D_P(g)$ are the outgoing slopes of $g$ on the
fixed arms of $Y$.  Here $D_P(g)$ is the residual slope divisor; its class is
the residue class of the germ of $g$ at $P$.

Suppose in addition that $g$ is ordinary CLD-smooth.  By
\cite[Remark~15.10]{GJR26}, its germ lies in the ordinary smooth harmonic
subsheaf $\mathcal F_P$.  The description of
$(\mathcal H/\mathcal F)_P$ by \cite[Lemma~2.3.22]{Thu05} then says that its
residual divisor is a real linear combination of principal divisors, or
equivalently that its class in
$\operatorname{Pic}^0(E)\otimes_{\mathbf Z}\mathbf R$ vanishes.  Together
with \eqref{eq:ch19-harmonic-reg-fixed-support}, this gives
\begin{equation}
\label{eq:ch19-harmonic-reg-picard-kernel}
 D_P(g)\in
 P_S:=\ker\!\left(
 V_S\longrightarrow\operatorname{Pic}^0(E)\otimes_{\mathbf Z}\mathbf R
 \right).
\end{equation}
The space $V_S$ is finite-dimensional and $P_S$ is a linear subspace, hence
$P_S$ is closed in the coefficient topology of $V_S$.

Suppose that a net $(h_i)$ of ordinary CLD-smooth harmonic functions on $U$
converges locally uniformly to $h$.  The inclusion of ordinary smooth forms
in weakly smooth forms, together with \cite[Propositions~10.14 and
10.17]{GJR26}, shows that every $h_i$ is harmonic in the
Gubler--Jell--Rabinoff sense.  The fixed-support argument makes $h_i$ and $h$
affine on each $[P,R_e]$.  If $\ell_e$ is the length of this segment, the
outgoing slope is
\begin{equation}
\label{eq:ch19-harmonic-reg-slope-values}
                         \frac{h_i(R_e)-h_i(P)}{\ell_e},
\end{equation}
and the analogous formula holds for $h$.  Uniform convergence on $Y$
therefore gives coefficientwise convergence
\begin{equation}
\label{eq:ch19-harmonic-reg-divisor-limit}
                         D_P(h_i)\longrightarrow D_P(h)
                         \quad\text{in }V_S.
\end{equation}
Every term on the left belongs to the closed subspace $P_S$ by
\eqref{eq:ch19-harmonic-reg-picard-kernel}.  Hence $D_P(h)\in P_S$, contrary
to \eqref{eq:ch19-harmonic-reg-nonzero-class}.

Finally, suppose that $(h_i)$ is a decreasing net of ordinary CLD-smooth
harmonic functions on $U$ with pointwise limit $h$.  On the compact set $Y$,
the net form of Dini's theorem makes the convergence uniform.  Indeed, for
$\varepsilon>0$ the increasing family of open sets
\[
              \{p\in Y:h_i(p)-h(p)<\varepsilon\}
\]
covers $Y$; compactness supplies finitely many indices and directedness one
index dominating them all.  The preceding contradiction applies.
\end{proof}

The finite-dimensional closure argument in
\Cref{prop:ch19-cld-smooth-harmonic-nonregularization} uses the vanishing
curvature of every approximant.  It does not apply to ordinary CLD-smooth psh
approximants: their positive curvature can move and concentrate, and uniform
convergence does not force endpoint derivatives of convex functions to
converge.  Thus the proposition contradicts neither
\cite[Proposition~15.11]{GJR26} nor
\Cref{thm:p063-current-psh-ordinary-regularization}.  The latter theorem
resolves the ordinary CLD-smooth psh regularization question left open in
\cite[Remark~5.5]{Wanner19}: positive curvature is what lets the residual
slope defect of an approximant escape into fresh tangent directions.

\subsection{Currents on a curve, and Thuillier's approximation theorem}
\label{sec:ch19-thuillier-currents}

The three assertions just separated are separated for a reason, and the
reason is visible already in \cite{Thu05}, where a complete
current-theoretic potential theory on curves is built and an unconditional
approximation theorem is proved---for \emph{his} notion of smoothness, which
is not the CLD one.  This subsection records the two statements of that
theory that sharpen what is said above, and then says precisely where the
two notions of smoothness part company.

We first fix Thuillier's vocabulary.  Let $k$ be a complete
non-Archimedean field with nontrivial absolute value and let $X$ be a
smooth strictly $k$-analytic curve.  Write $\mathrm I(X)$ for the set of
points of $X$ of type~$\mathrm{II}$ or~$\mathrm{III}$, equivalently the
set of Shilov points of the $k$-affinoid domains of $X$.  For an open
$\Omega\subseteq X$:
\begin{itemize}
\item $\mathrm A^0(\Omega)$ is the space of \emph{smooth} functions in the
sense of \cite[Def.~3.2.1(i)]{Thu05}: continuous real functions $\varphi$
on $\Omega$ admitting a locally finite cover of $\Omega$ by
$k$-affinoid domains $Y$ with $\varphi|_Y$ harmonic on the interior of
$Y$;
\item $\mathrm A^1(\Omega)$ is the space of \emph{smooth forms}: real
measures on $\Omega$ whose support is a locally finite subset of
$\mathrm I(\Omega)$ \cite[Def.~3.2.1(ii)]{Thu05};
\item $\mathrm D^0(\Omega)=\mathrm A^1_c(\Omega)'$ and
$\mathrm D^1(\Omega)=\mathrm A^0_c(\Omega)'$ are the currents of degree
$0$ and $1$, with the weak topology \cite[Def.~3.3.1]{Thu05}.
\end{itemize}
There is exactly one operator $\ddc:\mathrm A^0_X\to\mathrm A^1_X$
compatible with base change and with the discrete Laplacian
\eqref{eq:graph-laplacian} on every skeleton, that is with
$\ddc(\Phi\circ\tau_{\mathcal Y})=\Delta_{S(\mathcal Y)}\Phi$, and its
kernel is the sheaf of harmonic germs
\cite[Thm.~3.2.10 and Cor.~3.2.11(i)]{Thu05}.  It is symmetric,
$\int\varphi_1\,\ddc\varphi_2=\int\varphi_2\,\ddc\varphi_1$
\cite[Prop.~3.2.12]{Thu05}, hence transposes to a continuous operator
$\ddc:\mathrm D^0\to\mathrm D^1$.  Two structural facts make it behave like
an elliptic operator: it is \emph{regular}, in that
$\mathrm A^0(X)\subset\mathrm D^0(X)$ is exactly the set of degree-zero
currents $T$ with $\ddc T\in\mathrm A^1(X)$
\cite[Prop.~3.3.11]{Thu05}, and for $X$ connected and proper its image is
the hyperplane $\{\rho\in\mathrm D^1(X):\langle\rho,1\rangle=0\}$
of degree-one currents orthogonal to the constants
\cite[Thm.~3.3.13]{Thu05}.

\refstepcounter{theorem}\label{thm:ch19-thuillier-sh-currents}
\begin{quote}
\textbf{Foundational input \thetheorem\ (Thuillier's current-theoretic
characterization; source boundary).}
Let $k$ be a complete non-Archimedean field with nontrivial absolute value
and let $X$ be a smooth strictly $k$-analytic curve.
\begin{enumerate}
\item\label{itm:ch19-thu:inject}
Restriction to $\mathrm I(X)$ identifies a subharmonic function on $X$
with a degree-zero current, and the map
$\mathcal{SH}(X)\hookrightarrow\mathrm D^0(X)$ so obtained is injective.
\item\label{itm:ch19-thu:measures}
A degree-one current on $X$ is positive if and only if it is a positive
Radon measure on the locally compact space $|X|$.
\item\label{itm:ch19-thu:cone}
Under \cref{itm:ch19-thu:inject}, the real cone $\mathcal{SH}(X)$ is exactly
the set of $T\in\mathrm D^0(X)$ such that $\ddc T$ is positive, that is
such that $\langle T,\ddc\varphi\rangle\geq0$ for every
$\varphi\in\mathrm A^0_c(X)$ with $\varphi\geq0$.
\item\label{itm:ch19-thu:bijection}
If moreover $X$ is connected and is either the analytification of an affine
algebraic curve or a connected component of the interior of a $k$-affinoid
space, then $T\mapsto\ddc T$ induces a bijection from
$\mathcal{SH}(X)/\Gamma(X,\mathcal H_X)$ onto the positive Radon measures
on $|X|$.
\item\label{itm:ch19-thu:X0}
A subharmonic function on $X$ is determined by its restriction to
$X_0=\{x\in X:[\kappa(x):k]<\infty\}$.
\end{enumerate}
Items \cref{itm:ch19-thu:inject}--\cref{itm:ch19-thu:X0} are imported,
respectively, from \cite[Lem.~3.4.1, Thm.~3.4.12, Cor.~3.4.13,
Thm.~3.6.26]{Thu05}; no proof is claimed here.
\end{quote}

Item~\cref{itm:ch19-thu:cone} is the curve prototype of the whole CLD
picture: a potential is a current, and semipositivity of the potential is
positivity of its curvature current, with no regularity hypothesis
whatsoever.  It is the exact analogue of the classical statement that a
subharmonic function on a domain in $\mathbf C$ is a distribution with
positive Laplacian, and it is what the CLD calculus reproduces in higher
dimension only for the continuous---or at best the locally bounded---part of
the theory; see \Cref{sec:ch19-limitations}.  Item~\cref{itm:ch19-thu:X0} has
no counterpart asserted here for higher-dimensional singular CLD potentials:
it says that the rigid points already determine a subharmonic function.
Rigid points are dense on a strictly analytic curve, but density alone does
not determine an arbitrary upper semicontinuous function.  The additional
subharmonicity assertion is proved through the capacity theory of
\cite[\S3.6]{Thu05}.

\refstepcounter{theorem}\label{thm:ch19-thuillier-approximation}
\begin{quote}
\textbf{Foundational input \thetheorem\ (Thuillier's approximation
theorem; source boundary).}
Let $k$ be a complete non-Archimedean field with nontrivial absolute value,
let $X$ be a smooth strictly $k$-analytic curve, let
$\Omega\subseteq X$ be open and let $u\in\mathcal{SH}(\Omega)$.  Then for
every relatively compact open $\Omega'\Subset\Omega$ the restriction
$u|_{\Omega'}$ is the lower envelope of a decreasing family of
\emph{smooth} subharmonic functions on $\Omega'$, smooth being understood
in the sense of $\mathrm A^0$ above.  Moreover a smooth function
$\varphi\in\mathrm A^0(X)$ is subharmonic if and only if
$\ddc\varphi\geq0$.

No hypothesis is placed on $k$, on its residue field, or on $X$ beyond
those displayed, and the approximating family may be taken to be indexed by
the finite subsets $S\subseteq\mathrm I(Y)$ containing $\partial Y$, for
a compact $k$-analytic $Y$ with $\Omega'=Y\smallsetminus\partial Y$,
by taking $u_S$ harmonic on $Y\smallsetminus S$ and equal to $u$ on
$S$.
The approximation and the smooth characterization are imported from
\cite[Thm.~3.4.2, Prop.~3.4.4]{Thu05}; no proof is claimed here.
\end{quote}

The statement for metrics rather than functions is
\cite[Thm.~3.4.19]{Thu05}: on a proper smooth curve a subharmonic metric on
a line bundle is the \emph{increasing} upper envelope of a family of smooth
subharmonic metrics, the direction being reversed because the weight
$-\log\|s\|$ reverses it.  Both directions are used in
\cite[\S4.3.3]{Thu05} to deduce an arithmetic Hilbert--Samuel statement.

\begin{remark}[Thuillier-smooth is not CLD-smooth]
\label{rem:ch19-thuillier-smooth-vs-cld}
\Cref{thm:ch19-thuillier-approximation} and item~(3) of the warning above are
not in conflict, and the reason is a genuine difference of vocabulary that it
is worth stating once and for all.

A Thuillier-smooth function is one that is locally harmonic on affinoid
pieces; equivalently, by \cite[Prop.~3.2.4]{Thu05}, one that is locally of
the form $\Phi\circ\tau_{\mathcal Y}$ with $\Phi$ continuous and
piecewise affine on the skeleton $S(\mathcal Y)$.  In the language of this
book they are skeletal piecewise affine functions, and their $\ddc$ is
an atomic measure, by construction.  Break points of type~$3$ are allowed
in Thuillier's class; this description should not be identified without
further hypotheses with strict $\mathbf R$-PL model functions.  A
CLD-smooth function is locally $F\circ\trop$ with $F$ a $C^\infty$
function on $N_{\mathbf R}$; its $\dprimeddprime$ is absolutely
continuous along the edges of a skeleton and, as \Cref{rem:ch19-thuillier-balanced}
records, carries no atom at all.  The two classes are therefore
incomparable in an essential way, and neither contains the other:
$\max\{0,-u\}$ of \Cref{ex:ch19-p1-standard} is Thuillier-smooth and not
CLD-smooth, and $F_\varepsilon\circ\trop$ of \Cref{ex:ch19-p1-nonpl} is
CLD-smooth on the dense torus and not Thuillier-smooth there.

Consequently \Cref{thm:ch19-thuillier-approximation} gives a decreasing net
of Thuillier-smooth, hence skeletal piecewise affine, subharmonic functions
for arbitrary upper semicontinuous subharmonic $u$.  It does not by itself
give ordinary CLD-smooth approximants, because ordinary smoothness detects a
principal residual-divisor condition at type~$2$ points.  Wanner overcame
that obstruction under the three hypotheses of
\cite[Cor.~5.4]{Wanner19}; those are the cases in which the relevant
residual Picard obstruction is torsion, as in
\cite[Thm.~2.3.21]{Thu05}.

The proof below removes those hypotheses for \emph{psh} approximants.  It
first introduces an arbitrarily small strictly positive graph Laplacian.
At each type~$2$ vertex, Riemann--Roch moves the compensating effective
residual divisor away from the skeletal directions.  The remaining
degree-zero divisor is principal, hence comes from an ordinary smooth
harmonic germ; smooth convex caps then turn its negative fresh-direction
slopes into positive curvature.  This mechanism is unavailable for harmonic
approximants, which explains the preceding counterexample.
\end{remark}

\begin{remark}[The vanishing vertex atom of a smooth function on a balanced
space]
\label{rem:ch19-thuillier-balanced}
Here is the one-line computation behind the previous remark, in the tropical
language of \Cref{chap:tropical-spaces,chap:cld-poincare-lelong}.  Let
$X$ be a one-dimensional balanced weighted tropical space, let $p$ be a
vertex with tangent directions $\sigma\in T_pX$, weights $m_\sigma$ and
primitive generators $u_{\sigma/p}$, and let $F$ be smooth near the image
of $p$, so that $u=F\circ\trop$ is CLD-smooth.  The atom of
$\dprimeddprime[u]$ at $p$ is, by \Cref{lem:ch15-vertex-atom},
\[
 \sum_{\sigma\in T_pX}m_\sigma\,d_\sigma u(p)
 =\sum_{\sigma\in T_pX}m_\sigma\bigl\langle\rd F,u_{\sigma/p}\bigr\rangle
 =\Bigl\langle\rd F,\sum_{\sigma\in T_pX}m_\sigma u_{\sigma/p}\Bigr\rangle
 =0 ,
\]
the last equality being the balancing condition
\Cref{def:ch10-weighted-balanced}.  So on a balanced one-dimensional space
every CLD-smooth function has vanishing vertex atoms, and
$\dprimeddprime[F\circ\trop]$ is the edgewise measure $F''\rd s$.  A
function with a strictly positive vertex atom---for instance the model
function of \Cref{ex:ch19-p1-standard}---is therefore never CLD-smooth, yet
on every curve in the scope of
\Cref{thm:p063-current-psh-ordinary-regularization} it is a decreasing limit
of CLD-smooth psh functions, the curvature of the approximants concentrating
at the vertex as in \Cref{ex:ch19-p1-nonpl}.
\Cref{thm:ch19-thuillier-approximation}
is the analytic counterpart of exactly this phenomenon, with the roles of the
two smooth classes exchanged: Thuillier's smooth functions are the ones that
carry the atoms, so his approximation is unconditional, and the difficulty is
displaced onto the comparison of the two notions of smoothness rather than
onto the approximation itself.
\end{remark}

\subsection{Ordinary smooth monotone regularization on curves}
\label{sec:ch19-ordinary-smooth-curve-regularization}

We now prove the unrestricted continuous result.  The base field is kept
algebraically closed, complete, and nontrivially valued, exactly as in
Wanner's comparison theorem.  Properness is used to stay in the algebraic
semistable framework; the conclusion itself is local on the analytification.
The proof treats ordinary Chambert--Loir--Ducros smoothness, not the enlarged
weakly smooth sheaf.

% Exact candidate statement for hostile audit and production Rethlas.
% Date: 27 August 2026.
% This packet is isolated from the live book.  It makes no novelty or
% priority claim and must not be promoted before a final correct receipt.

\begin{lemma}
\label{lem:p063-local-tree-affinoid}
Let $K$ be algebraically closed, complete, and nontrivially valued, let
$C$ be a smooth algebraic $K$-curve, let $W\subset C^{\mathrm{an}}$ be
open, and let $x\in W$.  There is a connected strict affinoid domain
$Y\subset W$, with $x\in Y^\circ$, which admits a strictly semistable
skeleton that is a finite metric tree.  Moreover, if $f$ is lisse on a
neighbourhood of $Y$, then, after refining that skeleton and making a
finite metric subdivision, one can write
\[
                         f|_Y=F\circ\tau
\]
for the retraction $\tau:Y\to\Gamma$ onto a finite metric tree $\Gamma$
and a continuous function $F$ that is affine on every edge.  The added
vertices away from the boundary are of type $2$ or type $3$.
\end{lemma}

\begin{lemma}
\label{lem:p063-tree-rational-strictification}
Let $\Gamma$ be a finite metric tree, let $\partial\Gamma$ be a nonempty
set of vertices containing every endpoint designated as boundary, and let
$F:\Gamma\to\mathbf R$ be continuous and affine on every edge.  Suppose
that
\begin{equation}
\label{eq:p063-graph-subharmonicity}
 \Delta_\Gamma F(v):=\sum_{e\ni v}d_eF(v)\geq0
 \qquad
 \bigl(v\in V(\Gamma)\setminus\partial\Gamma\bigr),
\end{equation}
where $d_eF(v)$ denotes the outgoing slope.  For every $\varepsilon>0$
there is a continuous function $G:\Gamma\to\mathbf R$, affine with
rational slope on every edge, such that
\begin{equation}
\label{eq:p063-graph-rational-error}
             \lVert G-F\rVert_{L^\infty(\Gamma)}<\varepsilon
\end{equation}
and
\begin{equation}
\label{eq:p063-graph-strict-mass}
 \Delta_\Gamma G(v)>0
 \qquad
 \bigl(v\in V(\Gamma)\setminus\partial\Gamma\bigr).
\end{equation}
If there are no interior vertices, the last condition is empty.
\end{lemma}

\begin{lemma}
\label{lem:p063-residual-divisor-moving}
Let $C_0$ be a smooth projective connected curve of genus $g$ over an
infinite algebraically closed field $k_0$, let $M\subset C_0$ be finite, and let
$B\in\operatorname{Div}_{\mathbf Q}(C_0)$ have positive degree.  There are
an integer $N>0$ and an effective integral divisor $A$ such that
\begin{equation}
\label{eq:p063-moving-divisor}
          A\sim NB,
          \qquad
          \operatorname{Supp}(A)\cap M=\varnothing.
\end{equation}
Consequently $B-N^{-1}A$ is a rational multiple of a principal divisor.
\end{lemma}

\begin{lemma}
\label{lem:p063-principal-residual-germ}
Let $K$ be as in \cref{lem:p063-local-tree-affinoid}, let $C$ be a smooth
$K$-analytic curve, and let $v\in C$ be an interior point of type $2$.
Let $C_v$ be its residual smooth projective curve.  Choose finitely many
tangent directions $\xi$ at $v$, let $p_\xi\in C_v(\widetilde K)$ be the
corresponding residual points, and let $r_\xi\in\mathbf Q$ satisfy
\begin{equation}
\label{eq:p063-residual-degree-zero}
                         \sum_\xi r_\xi=0.
\end{equation}
Assume that
\[
                         D=\sum_\xi r_\xi[p_\xi]
\]
is a rational multiple of a principal divisor.  For every $c\in\mathbf R$
there is an ordinary Chambert--Loir--Ducros smooth harmonic germ $h$ at
$v$ such that
\begin{equation}
\label{eq:p063-prescribed-residual-slopes}
                    h(v)=c,
                    \qquad
                    d_\xi h(v)=r_\xi
\end{equation}
in the displayed directions and $d_\xi h(v)=0$ in every other tangent
direction.  Here ordinary smoothness means smoothness for analytic moment
maps, not weak smoothness for harmonic tropicalizations.
\end{lemma}

\begin{proposition}
\label{prop:p063-tree-lisse-ordinary-regularization}
Let $K$ be as above, let $C$ be a smooth proper algebraic $K$-curve, and
let $Y\subset C^{\mathrm{an}}$ be a connected strict affinoid domain
admitting a strictly semistable skeleton that is a finite tree.  Let
$f:Y\to\mathbf R$ be lisse and suppose that $f|_{Y^\circ}$ is
Thuillier-subharmonic.  For every $\varepsilon>0$ there is $g\in C(Y)$
such that
\begin{equation}
\label{eq:p063-lisse-uniform-regularization}
             \lVert g-f\rVert_{L^\infty(Y)}<\varepsilon
\end{equation}
and $g|_{Y^\circ}$ is ordinary Chambert--Loir--Ducros smooth and
plurisubharmonic.
\end{proposition}

\begin{theorem}
\label{thm:p063-current-psh-ordinary-regularization}
Let $K$ be an algebraically closed, complete, nontrivially valued
non-Archimedean field, let $C$ be a smooth proper algebraic $K$-curve,
and let $W\subset C^{\mathrm{an}}$ be open.  For a continuous function
$u:W\to\mathbf R$, the following conditions are equivalent.
\begin{enumerate}
\item
\begin{equation}
\label{eq:p063-current-positivity}
                         \mathrm d'\mathrm d''[u]\geq0.
\end{equation}
\item The function $u$ is Thuillier-subharmonic.
\item For every $x\in W$, there are a connected strict affinoid domain
$Y\subset W$, with $x\in Y^\circ$, and functions $u_n\in C(Y)$ such that
$u_n|_{Y^\circ}$ is ordinary Chambert--Loir--Ducros smooth and
plurisubharmonic,
\begin{equation}
\label{eq:p063-monotone-ordinary-approximation}
                         u_{n+1}\leq u_n\quad\text{on }Y,
                         \qquad
                         u_n\longrightarrow u|_Y
\end{equation}
uniformly on $Y$.
\end{enumerate}
Consequently, in this curve scope, current positivity, Thuillier
subharmonicity, and local membership of a continuous function in the first
ordinary Chambert--Loir--Ducros decreasing-limit stage agree.
\end{theorem}


The proof, including the local tree-affinoid and convex-cap constructions,
is given in \Cref{sec:appN-curve-regularization}.

\begin{remark}[The role of positive curvature]
\label{rem:ch19-positive-curvature-regularization}
The harmonic nonregularization result
\Cref{prop:ch19-cld-smooth-harmonic-nonregularization} fixes all slope mass
in a finite-dimensional residual divisor space.  In the proof above, the
small strict graph mass supplies positive degree at each interior type~$2$
vertex.  The moving-divisor argument transfers the compensating negative
slopes to fresh open-disc directions, while the convex cap replaces each
such negative slope by a nonnegative smooth curvature density.  Thus the
approximants need not remain in the closed harmonic subspace, and there is no
contradiction.
\end{remark}

\begin{example}[The standard metric on $\mathbf P^1$]
\label{ex:ch19-p1-standard}
On the chart $[1:T]$, the standard model metric on $\mathcal O(1)$
satisfies
\[
              -\log\|X_0\|=\log\max\{1,|T|\}.
\]
Put $u=-\log|T|$.  On the toric skeleton this weight is
$\max\{0,-u\}$, with one unit kink at $u=0$.  Hence
\[
                   c_1(\mathcal O(1),\|\cdot\|_{\mathrm{st}})
                          =\delta_{\zeta_G},
\]
where $\zeta_G$ is the Gauss point.  Its mass is one, the degree of
$\mathcal O(1)$.

The same statement read through \Cref{thm:wanner-curve-comparison}: the
function $\max\{0,-u\}$ is convex, hence Thuillier-subharmonic on the
apartment, hence has positive curvature current, and its Laplacian is the
unit atom at the kink.  Its CLD-psh membership here comes instead from the
finite continuous toric approximation statement in
\Cref{sec:ch19-toric-continuous}, namely
\cite[Theorem~5.3.11]{BGJKK25}.
Read through \Cref{sec:toric-convex}, it is the case $n=1$, $P=[0,1]$ of
the recession identity $F=\check h_P$, whose Alexandrov measure is
$\vol_M(P)\delta_0=\delta_0$ and whose non-Archimedean mass is
$1!\cdot1=1$.  Three languages, one atom.
\end{example}

\begin{example}[A continuous psh weight that is not PL]
\label{ex:ch19-p1-nonpl}
Keep $X=\mathbf P^1$, $L=\cO(1)$, and let $u=v(T)$ be the coordinate on
the apartment $\mathbf R$.  For $\varepsilon>0$ put
\[
 F_\varepsilon(u)=\varepsilon\log\bigl(1+\rme^{-u/\varepsilon}\bigr),
\]
as in \Cref{sec:toric-convex}.  Then $F_\varepsilon$ is smooth, strictly
convex, and $0\leq F_\varepsilon-\max\{0,-u\}\leq\varepsilon\log2$, so
the dense-torus weight $F_\varepsilon\circ\trop$ extends to a continuous
semipositive metric on $\cO(1)$, and
\[
 c_1\bigl(\cO(1),\|\cdot\|_{F_\varepsilon}\bigr)
 =F_\varepsilon''(u)\,\rd u
 =\frac{1}{\varepsilon}
   \frac{\rme^{-u/\varepsilon}}{\bigl(1+\rme^{-u/\varepsilon}\bigr)^2}\,\rd u,
\]
a measure on the toric skeleton, with smooth positive density on the
apartment and total mass
$\int_{\mathbf R}F_\varepsilon''=F_\varepsilon'(+\infty)-F_\varepsilon'(-\infty)
=0-(-1)=1=\deg\cO(1)$.  The metric is CLD-smooth on the dense torus;
finite continuous extension across the toric boundary does not make it
ordinary CLD-smooth there.  As $\varepsilon\downarrow0$ these measures
converge weakly to $\delta_{\zeta_G}$, recovering
\Cref{ex:ch19-p1-standard}.  The example shows that even on
$\mathbf P^1$ the curvature of a continuous semipositive metric need not
be atomic, so \eqref{eq:ch19-curve-formal-measure} is genuinely a statement
about \emph{model} metrics.
\end{example}

\section{Toric metrics and convex functions}
\label{sec:ch19-toric}

\subsection{Setup and the two dictionaries}
\label{sec:ch19-toric-setup}

Let $T=\operatorname{Spec}K[M]$ be a split torus, with cocharacter lattice
$N$, and let $X_\Sigma$ be the toric variety of a fan $\Sigma$.  The
extended tropicalization is
\[
                  \operatorname{trop}:X_\Sigma^{\mathrm{an}}\longrightarrow N_\Sigma,
\]
where $N_\Sigma$ is the tropical toric partial compactification.

A toric line bundle $L$, together with a toric meromorphic section $s$,
is encoded by an integral piecewise linear support function on $\Sigma$.
A torus-invariant metric is encoded on the dense torus by a function
$\phi:N_{\mathbf R}\to\mathbf R$ satisfying
\begin{equation}
                       -\log\|s(x)\|=\phi(\operatorname{trop}(x)).
\label{eq:toric-metric-potential}
\end{equation}
The boundary growth of $\phi$ is prescribed by $L$; it is not an
arbitrary convex function on $N_{\mathbf R}$.

Two dictionaries are in play and must not be confused.  The first is between
\emph{toric geometry} and \emph{convex geometry}: for a nef line bundle,
and in particular the ample bundles used below, $L$ with
its toric section corresponds to a polytope $P\subset M_{\mathbf R}$, and
by \Cref{sec:toric-convex} the correct recession function of the weight
$\phi$ is the reflected support function
\[
 \check h_P(u)=\sup_{m\in P}\langle-m,u\rangle,
\]
the reflection being forced by \eqref{eq:toric-trop}.  A toric metric on
$L$ has $\phi-\check h_P$ bounded.  For an arbitrary nonconvex weight this
boundedness alone does not ensure continuous extension across the toric
boundary: the metric ratio must extend continuously to $N_\Sigma$.  We
start with metrics already defined on $X_\Sigma^{\An}$.  With this convention,
different toric sections change $\phi$ and $\check h_P$ by the same
affine function, so nothing below depends on the choice.  The second
dictionary is between \emph{convex geometry} and \emph{tropical differential
calculus}: $\dprimeddprime\phi$ is the Hessian superform, its top power is
$n!$ times the Monge--Amp\`ere determinant by
\eqref{eq:appJ-real-hessian}, and the factor $n!$ is the one recorded in
normalization~\cref{itm:ch19-norm:factorial}.

The following lemma is the elementary reason why the two dictionaries are
compatible; it is the toric case of \Cref{prop:cld-local-normal-form}.

\begin{lemma}
\label{lem:ch19-toric-chart}
Let $\iota:N_{\mathbf R}\to T^{\An}$ be the Gauss section
\eqref{eq:toric-gauss-section}.  Then:
\begin{enumerate}
\item $\trop\circ\iota=\operatorname{id}$, $\iota$ is a homeomorphism onto the toric
 skeleton $\Sk(T)\subseteq T^{\An}$, and $\trop$ restricted to
 $\Sk(T)$ is its inverse;
\item a choice of basis of $M$ makes $\trop$ a smooth tropicalization
 in the sense of \Cref{chap:analytic-tropicalization}, with all coordinates
 of the form $-\log|\chi^{m}|$ for invertible $\chi^m$;
\item if $\phi$ is smooth on $N_{\mathbf R}$, its pullback
 $\phi\circ\trop$ is CLD-smooth on $T^{\An}$ and
 $\dprimeddprime[\phi\circ\trop]=\trop^*\dprimeddprime\phi$.
\end{enumerate}
\end{lemma}

\begin{proof}
Part~(1) is \Cref{lem:ch08-toric-gauss-section}.  For (2), the coordinates
$u_i=-\log|\chi^{m_i}|$ for a basis $(m_i)$ of $M$ are moment coordinates of
invertible functions on $T$, hence a smooth tropicalization by
\Cref{def:trop-local-skeleton}; the image is all of $N_{\mathbf R}$ and the
canonical calibration is the constant weight $1$ by
\Cref{thm:trop-calibration-canonical}, since $\trop$ is the standard torus
tropicalization of \Cref{prop:trop-torus-skeleton}.  Part (3) follows from the
definition of a CLD-smooth function on a tropical chart
(\Cref{def:ch13-cld-form}) together with the compatibility of $\dprime$ and
$\ddprime$ with chart pullback, \Cref{prop:cld-intrinsic-calculus}.  Because
the calibration is trivial, no lattice index intervenes; this is what makes
the toric case free of the correction factors of
\Cref{thm:trop-multiplicity-formula}.
There is no converse to (3) for intrinsic CLD smoothness.  For example,
$\min(0,u)^2\circ\trop$ on $\mathbf G_m^{\An}$ is CLD-smooth although
$\min(0,u)^2$ is not $C^\infty$ on $\mathbf R$
\cite[Example~2.3.5]{BGJKK25}.  A different moment chart can present an
intrinsically smooth function that is not smooth in this fixed chart.
\end{proof}

\begin{lemma}
\label{lem:ch19-toric-recession}
Let $\phi,\phi'$ be convex functions on $N_{\mathbf R}$ whose difference
is bounded.  Then their recession functions agree, their subgradient images
have the same closure, and
\[
 \MA_{\mathbf R}(\phi)(N_{\mathbf R})=\MA_{\mathbf R}(\phi')(N_{\mathbf R}).
\]
If moreover $\phi-\check h_P$ is bounded for a lattice polytope $P$, then
that common mass is $\vol_M(P)$.
\end{lemma}

\begin{proof}
The recession function of a convex $\phi$ is
$\lim_{t\to\infty}t^{-1}\phi(u_0+tw)$, which is unchanged by a bounded
perturbation.  By convex duality
(\Cref{thm:appD-fenchel-moreau} and \eqref{eq:appD-support}), the closure of
$\bigcup_u\partial\phi(u)$ is the closure of the domain of the Legendre transform
$\phi^*$, and the recession function of $\phi$ is the support function of
that domain; hence equal recession functions give the same closed subgradient
image.  The total Alexandrov mass is the Lebesgue measure of that image by
\eqref{eq:alexandrov-ma}.  The last assertion is
\Cref{lem:toric-real-total-mass}.
\end{proof}

\begin{theorem}\label{thm:toric-positivity}
Work over an algebraically closed complete nontrivially valued field.
For a piecewise smooth toric metric, psh, functorial psh, and functorial
delta-psh are equivalent to convexity of $\phi$.  For a formal toric metric,
this is also equivalent to semipositivity of the formal model.
Here psh and functorial psh have the positive-current meaning of
\cite{GK19}.  Membership in the book's approximation-generated CLD class
for finite continuous invariant metrics is supplied separately by
\Cref{sec:ch19-toric-continuous}.
\end{theorem}
\begin{proof}
Choose a toric polyhedral subdivision on which the metric weight is smooth
on every cell.  In the toric tropical frame, the Chern delta-current is the
sum of the facewise Hessian of $\phi$ and its corner cycle.  By
\Cref{thm:gk-tropical-frame-positivity}, it is positive exactly when the
Hessian is positive semidefinite on every cell and every corner weight is
nonnegative.  These two conditions are precisely convexity of the
piecewise-smooth function $\phi$.  The implication from convexity to
functorial delta-psh, for arbitrary algebraic pullbacks, is the additional
source input \cite[Proposition~8.13]{GK19}: its proof applies the tropical
positivity criterion of Proposition~8.7.  Checking only integer-affine
\emph{toric} pullbacks would not establish this functorial assertion.
The field hypothesis is the one used for ordinary delta-forms in that source.

For a formal toric metric, $\phi$ is rational piecewise affine.  Its corner
weights are the degrees of the residual toric line bundle on invariant
curves.  Nonnegativity of all those weights is equivalent to relative
nefness of the formal toric model, since the cone of curves of a complete
toric stratum is generated by invariant curves.  Thus convexity is also
equivalent to formal semipositivity, as in
\cite[Corollary~8.14]{GK19}.  With the alternative potential
$\psi=\log\|s\|=-\phi$, the same condition is concavity of $\psi$.
\end{proof}

\begin{corollary}
\label{cor:ch19-toric-convexity-test}
Let $\phi$ be a piecewise smooth toric metric weight for $L$ on the
dense torus.  Then $\overline L$ is psh over $T^{\An}$ if and only if
$\phi$ is convex on $N_{\mathbf R}$; and if $\Sigma$ is complete, so
that $\phi-\check h_P$ is bounded, this holds if and only if
$\phi=\check h_P+\psi$ with $\check h_P+\psi$ convex, in which case
$\phi$ is automatically the supremum of the affine functions it dominates.
\end{corollary}

\begin{proof}
The first assertion is \Cref{thm:toric-positivity}.  For the second, a
finite convex function on $N_{\mathbf R}$ is the supremum of its affine
minorants by \Cref{thm:appD-fenchel-moreau}, and the boundedness of
$\phi-\check h_P$ restricts the slopes of those minorants to $-P$ by
\Cref{lem:ch19-toric-recession}.
\end{proof}

\begin{remark}[Why toric positivity has no hidden content]
\label{rem:ch19-toric-no-content}
The equivalences of \Cref{thm:toric-positivity} make the invariant toric
family a useful normalization check: the stated positivity notions reduce
to convexity, and the geometric mass is $n!\vol_M(P)$.  The elementary
Hessian and corner calculation checks the sign; the imported functorial
comparison supplies the assertion about arbitrary algebraic pullbacks.
\Cref{sec:ch19-mass-ledger} records the independent degree checks.
\end{remark}

\subsection{Continuous and singular invariant metrics}
\label{sec:ch19-toric-continuous}

The continuous and singular toric comparison is more extensive.  Fix a
continuous reference Green function and let $\theta$ be its tropical Chern
current.  A function $\varphi$ on an open subset $U\subset N_\Sigma$ is
$\theta$-psh precisely when $\varphi\circ\operatorname{trop}$ defines a
continuous \emph{singular} invariant psh metric on
$\operatorname{trop}^{-1}(U)$.  For finite continuous potentials, this is
also equivalent to local uniform approximation by smooth psh metrics.
Here ``continuous singular'' permits extended-real weights with value
$-\infty$; it does not mean a finite continuous metric.  If
$\Sigma$ is complete, $U=N_\Sigma$, and the valuation is nontrivial,
the finite continuous conditions are
equivalent to global Zhang semipositivity.  This is the exact statement of
\cite[Theorem~5.3.11]{BGJKK25}.

Three features of this statement deserve emphasis, because they are the
features that the general theory does \emph{not} have.

\begin{enumerate}
\item\label{itm:ch19-toric-inv} \emph{Invariance is a hypothesis, not a
conclusion.}  The equivalence is asserted for torus-invariant metrics.  A
noninvariant continuous metric on a toric variety is not covered, and no
averaging argument reduces the general case to the invariant one, because
averaging over the compact torus does not preserve the class of psh metrics
in an obvious way; see \Cref{sec:ch19-limitations}.
\item\label{itm:ch19-toric-sing} \emph{Singular potentials are allowed on
the tropical side.}  A $\theta$-psh $\varphi$ may take the value
$-\infty$, and the correspondence with invariant metrics persists.  What
does not persist without further hypotheses is the existence of
Bedford--Taylor products; those require local boundedness, exactly as in
\Cref{chap:cld-bedford-taylor}.
\item\label{itm:ch19-toric-global} \emph{The last equivalence is global.}
Zhang semipositivity is a statement about uniform approximation by
semipositive \emph{model} metrics on a proper variety; that it agrees with a
local convexity condition is special to the invariant toric setting, and it
is the toric shadow of the comparison proved in
\Cref{sec:ch19-cld-bfj}.
\end{enumerate}

\subsection{The Monge--Amp\`ere comparison}
\label{sec:ch19-toric-ma}

Let $\overline L_1,\ldots,\overline L_q$ be locally psh-approximable toric
metrics corresponding to tropical potentials $\varphi_i$ and reference
currents $\theta_i$.  Then
\begin{equation}
 \operatorname{trop}_*\left(
      c_1(\overline L_1)\wedge\cdots\wedge c_1(\overline L_q)
                         \right)
       =\bigwedge_{i=1}^q\bigl(\dprimeddprime[\varphi_i]+\theta_i\bigr).
\label{eq:cld-toric-ma-comparison}
\end{equation}
This is \cite[Theorem~6.1.1]{BGJKK25}.  The pushforward is defined on
$U\subset N_\Sigma$ by duality with pullback of compactly supported
Lagerberg test forms under the proper toric tropicalization.  The
right-hand side is the tropical
Bedford--Taylor product, so \eqref{eq:cld-toric-ma-comparison} is an equality of
currents, and of measures when $q=\dim X_\Sigma$.

On an open subset of $N_{\mathbf R}$, if $\phi$ is $C^2$ and convex,
then
\begin{equation}
                  (\dprimeddprime\phi)^n
                    =n!\det\bigl(\partial_i\partial_j\phi\bigr)\,\rd x.
\label{eq:real-ma-density}
\end{equation}
For a nonsmooth convex $\phi$, the right-hand side extends to
$n!$ times its Alexandrov Monge--Amp\`ere measure.  Hence the
non-Archimedean toric measure is literally a real Monge--Amp\`ere measure
after tropical pushforward, with no hidden normalization.

We record the top-degree case in the form used in every computation below.

\begin{theorem}
\label{thm:ch19-toric-ma}
Let $X_\Sigma$ be a proper toric variety of dimension $n$ over a
complete nontrivially valued field, let $L$ be an ample toric line bundle
with polytope $P\subset M_{\mathbf R}$, and let $\|\cdot\|_\phi$ be a
continuous invariant semipositive metric with weight $\phi$, so that
$\phi$ is convex and $\phi-\check h_P$ is bounded.  Then
$c_1(L,\|\cdot\|_\phi)^n$ is supported on the toric skeleton,
\begin{equation}
\label{eq:ch19-toric-ma}
 \trop_*\,c_1\bigl(L,\|\cdot\|_\phi\bigr)^n
 =n!\,\MA_{\mathbf R}(\phi),
\end{equation}
and both sides have total mass
\[
                     n!\vol_M(P)=(L^n).
\]
\end{theorem}

\begin{proof}
\emph{Step 1: the model case.}  If $\phi$ is a convex rational piecewise
affine weight defining a toric rational model metric (including the
value-group conditions on its constants), \eqref{eq:ch19-toric-ma} is
\Cref{thm:toric-ma-comparison}, whose proof identifies the coefficient of
the Chambert--Loir measure at a vertex $u$ of the domain-of-affinity
decomposition with $n!\vol_M(\partial\phi(u))$.  That the CLD product agrees
with the Chambert--Loir measure for formal metrics uses
\Cref{thm:cld-formal-intersection} in its general field scope (and
\Cref{thm:cld-recovers-cl} over a DVR); agreement with the delta-form product
is \Cref{thm:gk-delta-cl-measure}, after completed algebraic closure for
that delta-form interpretation.

\emph{Step 2: support.}  Model measures in Step~1 are carried by the
Gauss section.  Their weak limits remain supported on the closed extended
skeleton $\iota(N_\Sigma)$.  The boundary noncharging in the toric
comparison removes $\iota(N_\Sigma\smallsetminus N_{\mathbf R})$.
For smooth tropical weights, the corresponding local support statement is
\cite[Proposition~2.3.6]{BGJKK25}: a pulled-back top-degree Lagerberg form
is supported on the Gauss section.  Fibrewise constancy of a potential,
without this support theorem or the model-limit argument, is not a proof
of the support of its curvature product.

\emph{Step 3: continuity in $\phi$.}  Global semipositive toric model
approximation follows from the toric metric dictionary
\cite[\S5.3.7 and Theorem~5.3.11]{BGJKK25}, equivalently the approximation
used in \Cref{cor:toric-ma-continuous}.  Its uniformity is on the compact
extended tropicalization after subtracting the reference weight.  Agreement
on a large ball and equality of recession functions alone would not prove
that global uniform statement.  On the non-Archimedean side, uniform
convergence of weights is uniform convergence of metrics, so
\Cref{thm:cld-bt-continuity} gives weak convergence of the products; on the
convex side, Alexandrov measures converge weakly under locally uniform
convergence of finite convex functions.  Passing to the limit in Step~1
proves \eqref{eq:ch19-toric-ma} in general.  This is
\Cref{cor:toric-ma-continuous}.

\emph{Step 4: the mass.}  By \Cref{lem:ch19-toric-recession} the total
Alexandrov mass is $\vol_M(P)$, so the left-hand side has mass
$n!\vol_M(P)$; and $n!\vol_M(P)=(L^n)$ is the toric degree formula
\eqref{eq:appJ-toric-degree}.  Independently, the mass of the top wedge of
Chern currents of any continuous semipositive metric on an ample $L$ is
$(L^n)$ by \Cref{thm:ch17-degree}; the two computations agree, which
is the check.
\end{proof}

\begin{example}[A polyhedral toric measure]
\label{ex:ch19-toric-polyhedral}
Let
\[
                  \phi(u)=\max_{1\leq j\leq r}
                              \{\langle m_j,u\rangle+c_j\}
\]
be a convex polyhedral function whose boundary behavior defines a metric on
a proper toric $n$-fold.  At a vertex $u_0$, let
$P_{u_0}\subset M_{\mathbf R}$ be the convex hull of the active slopes
$m_j$.  Then the Alexandrov mass at $u_0$ is
\[
                            \operatorname{vol}_M(P_{u_0}).
\]
Consequently the normalized superform, and hence non-Archimedean,
Monge--Amp\`ere measure has an atom of mass
$n!\operatorname{vol}_M(P_{u_0})$ at the canonical skeletal point over
$u_0$.  When the slopes and constants define a rational model metric,
this is the normalized lattice volume that
appears as a special-fibre intersection number on the associated toric
model.
\end{example}

\subsection{Three toric surfaces, computed}
\label{sec:ch19-toric-worked}

The next three examples exhibit, on surfaces, the three types of measure
that a tropical Monge--Amp\`ere operator can produce: purely atomic, carried
by a segment, and absolutely continuous.  In each case the total mass is
verified against the top self-intersection number.  All three use
\Cref{thm:ch19-toric-ma} and
normalization~\cref{itm:ch19-norm:factorial}; the reader who reproduces them
has verified the factor $n!=2$ three times over.

\begin{example}[Two atoms on $\mathbf P^1\times\mathbf P^1$]
\label{ex:ch19-toric-two-atoms}
\index{toric variety!Monge--Ampere measure computed}
Let $X=\mathbf P^1\times\mathbf P^1$, $N=\mathbf Z^2$,
$L=\cO(1,1)$, with polytope $P=[0,1]^2$, so that
\[
 \vol_M(P)=1,
 \qquad
 (L^2)=2!\cdot1=2=(H_1+H_2)^2 .
\]
The reflected support function is
\[
 \check h_P(u)=\max\{0,-u_1\}+\max\{0,-u_2\}
              =\max\{0,-u_1,-u_2,-u_1-u_2\}.
\]

Consider instead
\[
 \phi=\max\{\ell_0,\ell_1,\ell_2,\ell_3\},
 \qquad
 \begin{aligned}
  \ell_0&=0, &\ell_1&=-u_1,\\
  \ell_2&=-u_2, &\ell_3&=-u_1-u_2-1 .
 \end{aligned}
\]
The slopes $(0,0),(-1,0),(0,-1),(-1,-1)$ are the four vertices of
$-P$, so $\phi-\check h_P$ is bounded (indeed $-1\leq\phi-\check
h_P\leq0$) and $\phi$ is a model metric weight for $L$ as soon as
$1\in v(K^\times)$.  Note that $\phi$ is \emph{not} a sum of functions of
$u_1$ and of $u_2$ separately, because of the constant $-1$.

\emph{Locating the atoms.}  By
\Cref{lem:ch08-alexandrov-polyhedral}\,(3), an atom sits exactly at a point
where three of the four affine functions agree and are maximal.  The four
triples give
\[
\begin{array}{lll}
 \ell_0=\ell_1=\ell_2:& u=(0,0),
  &\ell_3=-1<0:\ \text{maximal};\\
 \ell_1=\ell_2=\ell_3:& u=(-1,-1),
  &\ell_0=0<1=\ell_1:\ \text{maximal};\\
 \ell_0=\ell_1=\ell_3:& u=(0,-1),
  &\ell_2=1>0:\ \text{not maximal};\\
 \ell_0=\ell_2=\ell_3:& u=(-1,0),
  &\ell_1=1>0:\ \text{not maximal},
\end{array}
\]
each value in the third column being taken at the point in the second.
\emph{The masses.}  By \Cref{lem:ch08-alexandrov-polyhedral}\,(1),
\[
 \partial\phi(0,0)=\operatorname{conv}\{(0,0),(-1,0),(0,-1)\},
\]
\[
 \partial\phi(-1,-1)=\operatorname{conv}\{(-1,0),(0,-1),(-1,-1)\},
\]
two triangles, each half of the unit square $-P$, of lattice area
$\tfrac12$.  Hence
\[
 \MA_{\mathbf R}(\phi)
 =\tfrac12\delta_{(0,0)}+\tfrac12\delta_{(-1,-1)},
 \qquad\text{mass }1=\vol_M(P),
\]
and by \eqref{eq:ch19-toric-ma}
\[
 \trop_*c_1\bigl(L,\|\cdot\|_\phi\bigr)^2
 =2!\left(\tfrac12\delta_{(0,0)}+\tfrac12\delta_{(-1,-1)}\right)
 =\delta_{(0,0)}+\delta_{(-1,-1)},
\]
of total mass $2=(L^2)$.

The two subgradient triangles meet along the segment from $(-1,0)$ to
$(0,-1)$, which has empty interior, so the disjointification in
\eqref{eq:alexandrov-ma} costs nothing; this is
\Cref{lem:ch08-alexandrov-polyhedral}\,(2) in action.

\emph{Reading the answer geometrically.}  The corner locus of $\phi$ is
the tropical curve consisting of the segment $\{u_1=u_2\in[-1,0]\}$
joining the two atoms, together with four unbounded rays.  It is the
wall complex of the two-dimensional domain-of-affinity subdivision, not
the full skeleton of a surface model.  Under the value-group hypothesis
already imposed, the associated toric model has two special-fibre surfaces
of residual self-intersection $1$, matching the two unit atoms.
\end{example}

\begin{example}[A measure carried by a segment]
\label{ex:ch19-toric-segment}
Keep $X=\mathbf P^1\times\mathbf P^1$ and $L=\cO(1,1)$.  Define
\[
 B(t)=\begin{cases}
        0, & t\geq0,\\
        \tfrac12t^2, & -1\leq t\leq0,\\
        -t-\tfrac12, & t\leq-1,
      \end{cases}
\]
a $C^1$ convex function with $B''=\one_{[-1,0]}$ in the sense of
measures, and put
\[
 \phi(u)=\max\{0,-u_1\}+B(u_2).
\]
Then $B(t)-\max\{0,-t\}$ is bounded by $\tfrac12$, so
$\phi-\check h_P$ is bounded and $\phi$ is a continuous semipositive
toric weight for $L$.  Since $\phi$ splits as a sum of a function of
$u_1$ and a function of $u_2$, its subgradient is the product
$\partial\phi(u)=\partial\max\{0,-\cdot\}(u_1)\times\{B'(u_2)\}$ and its
Alexandrov measure is the product measure
\[
 \MA_{\mathbf R}(\phi)
 =\delta_{u_1=0}\otimes\one_{[-1,0]}(u_2)\,\rd u_2 ,
\]
the uniform measure of mass $1$ on the segment
$S=\{0\}\times[-1,0]\subset N_{\mathbf R}$.  Its total mass is
$1\cdot1=1=\vol_M(P)$, and
\[
 \trop_*c_1\bigl(L,\|\cdot\|_\phi\bigr)^2
 =2\,\one_S\cdot(\text{lattice length on }S),
 \qquad\text{mass }2=(L^2).
\]
Thus the Monge--Amp\`ere measure of a perfectly ordinary continuous
semipositive toric metric on a surface can be carried by a
one-dimensional polyhedron and be neither atomic nor absolutely continuous.
This is the toric prototype of the phenomenon that
\Cref{sec:ch19-canonical} meets again on abelian varieties, where the
canonical measure decomposes into pieces of several dimensions.
\end{example}

\begin{example}[An absolutely continuous measure on $\mathbf P^2$]
\label{ex:ch19-toric-p2-smooth}
Let $X=\mathbf P^2$, $L=\cO(1)$, with polytope
$P=\operatorname{conv}\{0,e_1,e_2\}$, so that
\[
 \vol_M(P)=\tfrac12,
 \qquad
 (L^2)=2!\cdot\tfrac12=1 .
\]
The reflected support function is
$\check h_P(u)=\max\{0,-u_1,-u_2\}$.  For $\varepsilon>0$ put
\[
 \phi_\varepsilon(u)
 =\varepsilon\log\bigl(1+\rme^{-u_1/\varepsilon}+\rme^{-u_2/\varepsilon}\bigr),
\]
so that $0\leq\phi_\varepsilon-\check h_P\leq\varepsilon\log3$.  Write
\[
 s_0=1,\quad s_1=\rme^{-u_1/\varepsilon},\quad s_2=\rme^{-u_2/\varepsilon},
 \qquad S=s_0+s_1+s_2,\qquad p_i=s_i/S ,
\]
so that $p_0+p_1+p_2=1$ and $p_i>0$.  Then
\[
 \nabla\phi_\varepsilon=(-p_1,-p_2),
\]
which is a diffeomorphism from $N_{\mathbf R}$ onto the interior of
$-P=\{(x,y):x\leq0,\ y\leq0,\ x+y\geq-1\}$; this already gives the total
mass $\vol_M(-P)=\tfrac12$ by \eqref{eq:alexandrov-ma}.  Differentiating,
$\partial p_i/\partial u_i=-\varepsilon^{-1}p_i(1-p_i)$ and
$\partial p_i/\partial u_j=\varepsilon^{-1}p_ip_j$ for $i\neq j$, so the
Hessian is
\[
 D^2\phi_\varepsilon
 =\frac1\varepsilon
   \begin{pmatrix}
     p_1(1-p_1) & -p_1p_2\\
     -p_1p_2    & p_2(1-p_2)
   \end{pmatrix},
\]
a positive definite matrix, and
\[
 \det D^2\phi_\varepsilon
 =\frac{p_1p_2}{\varepsilon^2}
   \bigl[(1-p_1)(1-p_2)-p_1p_2\bigr]
 =\frac{p_1p_2}{\varepsilon^2}\bigl[1-p_1-p_2\bigr]
 =\frac{p_0p_1p_2}{\varepsilon^2}.
\]
Therefore, by \eqref{eq:smooth-real-ma} and \eqref{eq:ch19-toric-ma},
\begin{equation}
\label{eq:ch19-p2-density}
 \MA_{\mathbf R}(\phi_\varepsilon)
 =\frac{p_0p_1p_2}{\varepsilon^2}\,\rd u,
 \qquad
 \trop_*c_1\bigl(\cO(1),\|\cdot\|_{\phi_\varepsilon}\bigr)^2
 =\frac{2\,p_0p_1p_2}{\varepsilon^2}\,\rd u .
\end{equation}
\emph{Mass check, twice.}  Substituting $p=(p_1,p_2)$ as coordinates,
$\rd p=\det D^2\phi_\varepsilon\,\rd u$, so
\[
 \int_{N_{\mathbf R}}\det D^2\phi_\varepsilon\,\rd u
 =\int_{\{p_1,p_2>0,\ p_1+p_2<1\}}\rd p=\tfrac12 ,
\]
which is $\vol_M(P)$; the non-Archimedean mass is
$2\cdot\tfrac12=1=(\cO(1)^2)$.  Independently, the total mass is
$\vol_M$ of the closure of the gradient image, which is $-P$, by
\Cref{lem:ch19-toric-recession}.

\emph{Degeneration.}  As $\varepsilon\downarrow0$, the functions
$\phi_\varepsilon$ decrease uniformly to $\check h_P$.
The Alexandrov measures in \eqref{eq:ch19-p2-density} converge weakly to
$\tfrac12\delta_0$, and the pushed-forward curvature measures to $\delta_0$
on $N_{\mathbf R}$.  The analytic curvature measures converge to
$\delta_{\iota(0)}$, of mass $1=(\cO(1)^2)$, the measure of the standard
model metric in \Cref{ex:ch08-toric-pn}.  The example therefore also
illustrates \Cref{thm:cld-bt-continuity}: continuous semipositive metrics,
smooth on the dense torus,
whose curvature measures concentrate on a point of the skeleton.
\end{example}

\begin{example}[The Hirzebruch surface, revisited as a mass check]
\label{ex:ch19-toric-hirzebruch}
\Cref{ex:ch08-toric-hirzebruch} computed, for
$\mathbf F_1=\operatorname{Bl}_p\mathbf P^2$ and $L=2H-E$, the Alexandrov measure
$\MA_{\mathbf R}(F)=\delta_{(1,1)}+\tfrac12\delta_{(0,-1)}$ of a certain
model metric function $F$, with total mass
$\tfrac32=\vol_M(P)$.  By \eqref{eq:ch19-toric-ma} the associated
non-Archimedean measure is
\[
 \trop_*c_1(L,\|\cdot\|_F)^2
 =2\delta_{(1,1)}+\delta_{(0,-1)},
 \qquad\text{mass }3=(2H-E)^2=4-1 .
\]
The point of repeating this here is that in \Cref{chap:model-calculations}
the identity was derived from the toric intersection formula on a model,
whereas \Cref{thm:ch19-toric-ma} derives it from the Bedford--Taylor product
of Chern currents.  The two atoms, $2$ and $1$, are the residual
self-intersection numbers of the two special-fibre surfaces of that model.
Agreement of these two computations is one of the three transfers listed in
\Cref{sec:ch16-transfer}.
\end{example}

\subsection{Invariant Monge--Amp\`ere equations}
\label{sec:ch19-toric-equation}

The theory in \cite{BGJKK25} also solves invariant Monge--Amp\`ere equations.
For an ample toric line bundle on a proper toric variety and a positive
measure of the correct mass supported on the canonical skeleton, with no
mass on its tropical boundary, there is a toric continuous singular psh
metric solving the equation, unique up to scaling
\cite[Theorem~6.3.13]{BGJKK25}.  This statement goes beyond bounded
potentials; its invariant hypotheses should not be dropped.

\begin{example}[Solving the equation in dimension one]
\label{ex:ch19-toric-solve-p1}
For $X=\mathbf P^1$, $L=\cO(d)$, and $P=[0,d]$, the equation reads
$1!\,\MA_{\mathbf R}(\phi)=\mu$ for a positive measure $\mu$ of mass
$d$ on $\mathbf R$.  A convex $\phi$ with $\phi''=\mu$ exists and is
unique up to an affine function.  The slope normalization
$\phi'(-\infty)=-d$ and $\phi'(+\infty)=0$ selects it up to a constant:
one can integrate the bounded monotone function
$\phi'(t)=-\mu((t,+\infty))$ from $0$ to $u$.
These slope limits give the prescribed recession function, but do
\emph{not} imply that $\phi-\check h_{[0,d]}$ is bounded.
Finite continuous extension of the metric at both toric endpoints is
equivalent here to $\int_{\mathbf R}|t|\,\rd\mu(t)<\infty$.
For example $d(1+t)^{-2}\mathbf1_{t\geq0}\,\rd t$ has the required mass
and slope limits but infinite first moment, so its solution has an unbounded
metric ratio at one endpoint.  This is why the general existence theorem
produces a continuous \emph{singular} metric and uses the extended
Monge--Amp\`ere construction of \cite[\S6.3]{BGJKK25}.
``No mass at the tropical boundary'' means that $\mu$ is a measure on
$\mathbf R$ with no endpoint atoms.  Compare \Cref{prop:graph-ma-solution} and
\eqref{eq:curve-poisson}, which solve the same problem on a compact metric
graph.
\end{example}

\section{Model and good-reduction examples}
\label{sec:ch19-models}

Curves and toric varieties are the two families in which the local calculus
can be reduced entirely to convex analysis.  In arbitrary dimension no such
reduction is available, but for \emph{model} metrics the answer is still
completely explicit: the measure is atomic, carried by the finitely many
divisorial points over the generic points of the special fibre, and its
coefficients are residual intersection numbers.  This section records that
computation and its immediate consequences, in the exact form used in
\Cref{sec:ch19-cld-bfj}.

\begin{example}[Good reduction in every dimension]
\label{ex:ch19-good-reduction}
Let $\mathcal X$ be a proper smooth model of a smooth projective
$n$-fold, with geometrically integral special fibre, and let $\mathcal L$
extend $L$.  The formal metric has
\[
                       c_1(\overline L)^n=(L^n)\delta_{x_s},
\]
where $x_s$ is the divisorial point reducing to the generic point of
$\mathcal X_s$.  For
$(\mathbf P^n,\mathcal O(1))$ with its standard model, this is
$\delta_{\zeta_G}$.  This is simultaneously the CLD product, the
Gubler--K\"unnemann delta-form product after completed algebraic closure,
and the Chambert--Loir measure.

The general CLD formal-intersection calculation is
\Cref{thm:cld-formal-intersection}; the delta-form identification with the
Chambert--Loir measure is \Cref{thm:gk-delta-cl-measure}, applied after
completed algebraic closure.  In the discretely valued setting,
\Cref{thm:cld-recovers-cl} gives the explicit component formula for the
same measure.  The mass
is $(L^n)$ because $\mathcal L|_{\mathcal X_s}$ has the same top
self-intersection as $L$, by flatness and specialization invariance
(\Cref{thm:model-projection-formula}).
\end{example}

\begin{example}[Several special-fibre components]
\label{ex:ch19-several-components}
Let $K$ be discretely valued and $\mathcal X$ a normal proper flat model,
with $\mathcal X_s=\sum_Yb_YY$, and let
$\mathcal L_i$ extend $L_i$.  Then
\begin{equation}
\label{eq:ch19-model-mixed-measure}
 \bigwedge_{i=1}^n c_1(\overline L_i)
    =\sum_Y b_Y
       \bigl(c_1(\mathcal L_1|_Y)\cdots
             c_1(\mathcal L_n|_Y)\bigr)\delta_{x_Y}.
\end{equation}
If all $\mathcal L_i$ are relatively nef, every coefficient is
nonnegative.  If one metric is changed by a vertical divisor, multilinearity
turns the change of measure into the corresponding special-fibre
intersection with that divisor.  This is the higher-dimensional analogue of
\eqref{eq:exceptional-model-laplacian}.

Formula \eqref{eq:ch19-model-mixed-measure} is
\eqref{eq:cld-chambert-loir-formula}, and it is also the BFJ mixed model
measure \eqref{eq:model-mixed-measure}; the coincidence of these three
expressions is the algebraic half of \Cref{thm:ch19-cld-bfj-comparison}.
The points $x_Y$ are normalized by $v_{x_Y}(\pi)=v_K(\pi)$; their
multiplicities occur in the coefficients, not a second time in the points.
Over a general valuation ring, the formal-intersection factors are those
of \Cref{thm:cld-formal-intersection}, rather than this DVR formula.
\end{example}

\begin{proposition}
\label{prop:ch19-model-basic}
With the notation of \Cref{ex:ch19-several-components}, assume
$\mathcal X$ proper over $K^\circ$.  Then:
\begin{enumerate}
\item\label{itm:ch19-model:mass} the total mass of
 \eqref{eq:ch19-model-mixed-measure} is $(L_1\cdots L_n)_X$;
\item\label{itm:ch19-model:pos} if every $\mathcal L_i$ is relatively nef,
 the measure is positive;
\item\label{itm:ch19-model:vert} if $\mathcal L_1$ is replaced by
 $\mathcal L_1(D)$ for a vertical Cartier divisor $D$, the measure
 changes by
 $\sum_Yb_Y\bigl(\cO(D)|_Y\cdot c_1(\mathcal L_2|_Y)\cdots
   c_1(\mathcal L_n|_Y)\bigr)\delta_{x_Y}$, a signed measure of total mass
 zero.
\end{enumerate}
\end{proposition}

\begin{proof}
\cref{itm:ch19-model:mass} Summing \eqref{eq:ch19-model-mixed-measure} over
$Y$ computes the mixed intersection number of the $\mathcal L_i$ against the
fundamental cycle $\sum_Yb_Y[Y]=[\mathcal X_s]$, which by flatness and
\Cref{thm:model-projection-formula} equals $(L_1\cdots L_n)_X$.

\cref{itm:ch19-model:pos} If each $\mathcal L_i$ is relatively nef, then each
restriction $\mathcal L_i|_Y$ is nef on the proper $\widetilde K$-variety
$Y$, and a mixed intersection of nef classes is nonnegative; the
multiplicities $b_Y$ are positive.  No converse is asserted: positivity of
the top mixed measure alone does not characterize relative nefness, and the
correct characterization is by restriction to residual curves, which is
\Cref{thm:common-model-core}.

\cref{itm:ch19-model:vert} Multilinearity of the intersection product gives
the displayed difference, and its total mass is
$(\cO(D)|_X\cdot L_2\cdots L_n)=0$ because $\cO(D)$ is trivial on the generic
fibre; this is \Cref{prop:special-fibre-degree-conservation}.
\end{proof}

\begin{example}[Migration of mass under a vertical blowup]
\label{ex:ch19-mass-migration}
Take $X=\mathbf P^2_K$ with $K$ discretely valued, $v(\pi)=1$,
$L=\cO(1)$, and compare two model metrics.  The standard model
$\mathbf P^2_{K^\circ}$ with $\cL=\cO(1)$ gives
$c_1(\overline L)^2=\delta_{\zeta_G}$, the Gauss point, of mass
$1=(L^2)$.  Blowing up a rational toric fixed point $p$ of the special fibre, as in
\Cref{ex:metrics-cl-toric-surface}, produces a model with two components
$\widetilde{X_s}$ and $E$, multiplicities both one, and the model bundle
$\cL'=\rho^*\cO(1)\otimes\cO(-E)$; its restrictions have
$(\cL'|_E)^2=1$ and $(\cL'|_{\widetilde{X_s}})^2=0$, so
\[
 c_1(\overline L')^2=\delta_{x_E},
\]
again of mass $1$.  The unit of mass has moved from the Gauss point
$u=(0,0)$ to the monomial point $u=(1,1)$.  This is
\eqref{eq:metrics-toric-measure}, and it is the two-dimensional analogue of
\eqref{eq:exceptional-model-laplacian}: a vertical modification cannot change
the total mass, only its distribution.
\end{example}

\begin{example}[A multiplicity-two fibre as a normalization check]
\label{ex:ch19-nonreduced}
The multiplicities $b_Y$ in \eqref{eq:ch19-model-mixed-measure} are not
cosmetic.  Let $K$ be discretely valued, and suppose that a proper curve
$C/K$ has a normal proper model $\mathcal C$ whose special fibre is the
single component $2Y$.  Let $L$ have degree $d$, and let
$\mathcal L$ be a $\mathbf Q$-model of $L$.  Put
$e=\deg(\mathcal L|_Y)$.  Then
\[
 c_1(\overline L)=2e\,\delta_{x_Y},
\]
while conservation of degree gives
\[
                         2e=d.
\]
If $\mathcal L$ is an honest model line bundle, then $e\in\mathbf Z$,
so $d$ must be even.  For odd $d$, a model in this one-component
configuration must therefore use rational coefficients, with
$e=d/2$.  Dropping either the multiplicity $b_Y=2$ or the denominator
would give the wrong mass.  The corresponding normalization of the
divisorial point is $v_{x_Y}=\tfrac12\ord_Y$, as in
\Cref{rem:divisorial-normalization}.  Notice that this is a conditional
normalization check; it does not assert that an arbitrary curve, and in
particular $\mathbf P^1$, admits such a model.
\end{example}

\begin{remark}[What the model case does not give]
\label{rem:ch19-model-scope}
\Cref{ex:ch19-good-reduction,ex:ch19-several-components} determine the
measure attached to a model metric completely, and by uniform approximation
they determine it for continuous semipositive metrics as well
(\Cref{thm:cl-continuous-extension}).  They say nothing about three
questions that occupy the rest of this book.  First, whether an arbitrary
positive measure of the right mass arises this way---the Monge--Amp\`ere
existence problem, taken up in \Cref{chap:bfj-ma}.  Second, whether an
arbitrary positive CLD current is the curvature of some metric---there is no
representation theorem, and none is claimed.  Third, what happens when the
potential is unbounded, where \eqref{eq:ch19-model-mixed-measure} alone
does not define a product; this is the first row of the table in
\Cref{sec:ch19-limitations}.
\end{remark}

\section{Raynaud uniformization of abelian varieties}
\label{sec:ch19-abelian}

Abelian varieties are the third and last family of this chapter, and the
one where the local calculus is genuinely needed: neither the curve
reduction nor the toric reduction applies, and the answer---a Haar measure
on a real torus of dimension the torus rank---is not visible from any single
model.

\subsection{The uniformization}
\label{sec:ch19-raynaud}

Let $A$ be an abelian variety over a complete nontrivially valued field.
After a finite separable extension, Raynaud uniformization gives
\begin{equation}
                  0\longrightarrow T\longrightarrow E
                    \stackrel{q}{\longrightarrow}B\longrightarrow0,
             \qquad A^{\mathrm{an}}=E^{\mathrm{an}}/\Lambda,
\label{eq:raynaud-uniformization}
\end{equation}
where $T$ is a split torus, $B$ has good reduction, and $\Lambda$ is
a discrete cocompact subgroup.  If $N$ is the cocharacter lattice of
$T$, the canonical tropicalization descends to
\[
                  \operatorname{trop}_A:A^{\mathrm{an}}\longrightarrow
                                      N_{\mathbf R}/\operatorname{trop}(\Lambda).
\]

We fix once and for all the notation that will be used for the rest of the
chapter.  Write
\[
 g=\dim A,\qquad r=\dim T,\qquad b=\dim B,\qquad g=b+r,
\]
and call $r$ the \emph{torus rank} of $A$.  We say the situation is
\emph{split} when $B$ has good reduction over $K$ itself and $T$ is a
split torus, so that $\Lambda\simeq\mathbf Z^r$; this is achieved after a
finite separable extension of $K$, and every statement below is made in
the split case.  The two extreme cases have names:
\[
 r=0:\ \text{good reduction};
 \qquad
 r=g\ (\text{i.e.\ }b=0):\ \text{totally degenerate reduction}.
\]
Every intermediate value $0<r<g$ is called \emph{mixed reduction}, and it
is exactly there that the delicate questions of
\Cref{rem:ch19-canonical-delta-scope} live.

\begin{theorem}
\label{thm:ch19-raynaud}
Let $K$ be complete and nontrivially valued and let $A$ be an abelian
variety over $K$.  There exist a finite separable extension $K'/K$, an
abelian variety $B$ over $K'$ with good reduction, a split torus $T$
over $K'$, an extension $E$ of $B$ by $T$ in the category of
algebraic groups, and a discrete subgroup $\Lambda\subseteq E(K')$ free of
rank $r=\dim T$, such that
\eqref{eq:raynaud-uniformization} holds as analytic groups over $K'$.
Moreover:
\begin{enumerate}
\item\label{itm:ch19-ray:neron} $E$ contains a canonical compact analytic
 subgroup $N_{\mathrm{er}}$, the \emph{N\'eron domain} of $A$, which is
 an extension of $B^{\An}$ by the affinoid torus $T_1=\{|\chi^m|=1\}$
 and is a formal analytic group with semiabelian reduction;
\item\label{itm:ch19-ray:trop} the map
 $\tau:E^{\An}\to N_{\mathbf R}$ obtained by trivializing
 $E\to B^{\An}$ over a formal affinoid cover and tropicalizing the torus
 factor is well defined, surjective, proper, piecewise linear, and
 $G$-smooth;
\item\label{itm:ch19-ray:lattice} $\tau(\Lambda)$ is a cocompact lattice
 in $N_{\mathbf R}$, and $\tau$ descends to a surjection
 $\trop_A:A^{\An}\to N_{\mathbf R}/\tau(\Lambda)$;
\item\label{itm:ch19-ray:skeleton} there is a canonical deformation
 retraction of $E^{\An}$ onto a closed subset $S(E)$ which $\tau$
 identifies with $N_{\mathbf R}$, and it descends to a deformation
 retraction of $A^{\An}$ onto the \emph{skeleton}
 $S(A)=S(E)/\Lambda\simeq N_{\mathbf R}/\tau(\Lambda)$, a real torus of
 dimension $r$.
\end{enumerate}
\end{theorem}

\begin{proof}
This is a theorem of Raynaud, in the form written out by Bosch and
L\"utkebohmert; we quote it as
\cite[Theorem~B.5.7, \S\S B.5.8--B.5.12 and \S19.2.1--19.2.2]{CLD25}, whose
sources are \cite{BoschLutkebohmert1993}, and refer to
\cite[Chapter~6]{Berk12} for the retraction in
\cref{itm:ch19-ray:skeleton}.  We indicate only where each assertion comes
from.

The existence of $(B,T,E,\Lambda)$ after a finite separable extension is the
uniformization theorem itself; the N\'eron domain of
\cref{itm:ch19-ray:neron} is the unique maximal compact analytic subgroup
that is formal analytic with semiabelian reduction, and its uniqueness makes
its formation commute with base change.

For \cref{itm:ch19-ray:trop}: let $\mathcal V$ be a formal affinoid cover of
$B$ trivializing the extension $N_{\mathrm{er}}\to B^{\An}$, with sections
$f_V$.  On $V\cap W$ the sections differ by a morphism
$f_{VW}:V\cap W\to T_1$ with values in the \emph{unit} torus, so the
composites $q^{-1}(V)\simeq V\times T\to T\to N_{\mathbf R}$ agree on
overlaps and glue to a global map $\tau$.  Each local expression is
$-\log|\chi^m|$ for a character $\chi^m$, hence $G$-smooth and piecewise
linear.  Restricted to a fibre of $q$, $\tau$ is the ordinary
tropicalization of a torus, which is proper; since $B^{\An}$ is compact,
$\tau$ is proper and surjective.

Assertion \cref{itm:ch19-ray:lattice} is the statement that the composite
$\Lambda\to E(K')\to N_{\mathbf R}$ has cocompact image, which is part of
the uniformization theorem; \cref{itm:ch19-ray:skeleton} is Berkovich's
retraction of $E^{\An}$ onto the skeleton of the torus fibre over the Gauss
point of $B$, which is $\Lambda$-equivariant because $\Lambda$ acts by
translations fixing that point.
\end{proof}

\begin{remark}[The skeleton and the splitting field]
\label{rem:ch19-skeleton-intrinsic}
On the chosen splitting field, $S(A)$ is the image of the canonical
deformation retraction, and it is identified with the real torus by
\Cref{thm:ch19-raynaud}\,\cref{itm:ch19-ray:skeleton}.  Over the original
field its skeleton is the Galois quotient of the split skeleton
\cite[\S8.4.6]{BGJKK25}; that quotient need not itself be a real torus.
The geometric torus rank $r$ is an invariant of $A$.  The identification
with an essential skeleton is used only under the discretely valued,
characteristic-zero hypotheses of \Cref{sec:essential-skeleton}.
For $g=1$ the potential reduction dichotomy is that of \Cref{ex:ch08-tate-curve}: $r=0$ is potentially good
reduction, $r=1$ is multiplicative reduction, and the two are separated by
the sign of $v(j)$.
\end{remark}

\subsection{The canonical tropicalization is harmonic}
\label{sec:ch19-abelian-harmonic}

\begin{theorem}\label{thm:abelian-harmonic-tropicalization}
The coordinate functions of the canonical tropicalization on
$E^{\mathrm{an}}$ are integral harmonic functions.  Hence
$\operatorname{trop}_A$ is a harmonic tropicalization in the sense of
\Cref{chap:delta-forms}.
\end{theorem}
\begin{proof}
Let $\chi$ be a character of the split torus $T$.  The pushout of the
Raynaud extension by $\chi:T\to\mathbf G_m$ is a
$\mathbf G_m$-torsor over $B$, hence a rigidified translation-invariant line
bundle $M_\chi$ on $B$.  Because $B$ has good reduction, $M_\chi$ extends to
a formal line bundle whose restriction to the special fibre is numerically
trivial.  Its canonical formal metric pulls back to the torsor, and the
logarithmic norm of the tautological section is
\[
                         -\log|\chi|.
\]
This is exactly the tropical coordinate associated to $\chi$.  The residual
line bundle of this PL function is therefore numerically trivial; by
definition, the coordinate is integral harmonic.

A lattice basis of the character group gives all coordinates of the
canonical map $E^{\mathrm{an}}\to N_{\mathbf R}$.  Each is integral harmonic,
so this map is a harmonic tropicalization.  Translation by an element of
$\Lambda$ adds the integral affine constant determined by its character
values.  Hence the map descends to
$N_{\mathbf R}/\operatorname{trop}(\Lambda)$ on the quotient
$A^{\mathrm{an}}$.  This is \cite[Proposition~16.2]{GJR26}.
\end{proof}

The proof deserves to be unwound once, because the same three-step pattern
recurs in \Cref{sec:ch19-canonical}: an analytic object on $E$ is
identified with a line bundle on $B$; the good reduction of $B$ turns
that line bundle into a formal one; and the numerical triviality of its
residual bundle is what produces harmonicity.

\begin{proposition}
\label{prop:ch19-tau-coordinates}
Keep the notation of \Cref{thm:ch19-raynaud} and let
$\chi^{m}$, $m\in M=N^\vee$, be a character of $T$.  Let $Q_m$ be the
line bundle on $B$ obtained by pushing out the extension
$N_{\mathrm{er}}\to B^{\An}$ along $\chi^m:T_1\to\mathbf G_{m,1}$, and
let $\mathcal Q_m$ be its canonical formal extension to the abelian formal
scheme $\mathcal B$ over $K^\circ$.  Then:
\begin{enumerate}
\item\label{itm:ch19-tau:torsor} the pushout gives a morphism
 $E\to Q_m^\times$ over $B$.  Pulling back the tautological section gives
 a nowhere-vanishing frame $e_m$ of $q^*Q_m$, and
 $\langle\tau(\cdot),m\rangle=-\log\|e_m\|$ for the pulled-back formal
 metric of $\mathcal Q_m$.  This morphism is not in general an open
 embedding when $r>1$;
\item\label{itm:ch19-tau:numtriv} $Q_{m,s}$ is algebraically equivalent to
 zero on the abelian variety $\mathcal B_s$, hence numerically trivial;
\item\label{itm:ch19-tau:harmonic}
 $\dprimeddprime\bigl[\langle\tau(\cdot),m\rangle\bigr]=0$ for every
 $m\in M$, that is, all coordinates of $\tau$ are harmonic in the sense
 of \Cref{def:ch16-harmonic-function} and integral harmonic in the sense of
 \Cref{sec:harmonic-trop}.
\end{enumerate}
\end{proposition}

\begin{proof}
\cref{itm:ch19-tau:torsor} An extension of $B$ by $\mathbf G_m$ is the
complement of the zero section in the total space of a line bundle
algebraically equivalent to zero; pushing out along $\chi^m$ realizes the
$T$-torsor $E\to B$ as the fibre product of the $\mathbf G_m$-torsors
attached to a basis of $M$.  Since $\mathcal B$ is a formal abelian scheme,
$Q_m$ extends to $\mathcal Q_m$ and the extension is unique up to a unit
\cite[Lemma~B.2.2 and Remark~19.2.3]{CLD25}.  The logarithmic norm of its
nonzero tautological section, pulled back along $E\to Q_m^\times$, gives exactly
$\langle\tau(\cdot),m\rangle$ by the construction of $\tau$ in
\Cref{thm:ch19-raynaud}\,\cref{itm:ch19-ray:trop}.

\cref{itm:ch19-tau:numtriv} The line bundle $Q_m$ classifies an extension,
hence lies in $\Pic^0(B)$; the same is true of its formal extension
restricted to the special fibre, because algebraic equivalence to zero is
detected by triviality of the associated Mumford bundle and specializes.  A
line bundle algebraically equivalent to zero on an abelian variety is
numerically trivial.

\cref{itm:ch19-tau:harmonic} The formal model of $q^*Q_m$ has
numerically trivial residual bundle, and $e_m$ identifies its generic fibre
with the trivial bundle.  Thus its logarithmic norm is integral harmonic
by \cite[Lemma~7.4 and Proposition~16.2]{GJR26}.  Equivalently,
\cite[Theorem~18.8.4 and Remark~19.2.3]{CLD25} give
$c_1(q^*\overline{Q_m})=0$ as a first Chern \emph{current}.  Applying
\eqref{eq:cld-pl-section-sign} to $e_m$, whose divisor is empty, gives the
displayed zero-current equation.  A first Chern current in dimension
$b>1$ is not the atomic top-degree measure of
\Cref{thm:cld-recovers-cl}; that measure formula alone would not prove
the required vanishing.
\end{proof}

\subsection{When is the canonical tropicalization smooth?}
\label{sec:ch19-tau-smooth}

Harmonicity is strictly weaker than smoothness in the CLD sense, and the
gap is the whole subject of \Cref{sec:harmonic-trop}.  On abelian varieties
the gap can be described exactly.

\begin{proposition}
\label{prop:ch19-tau-smooth}
Keep the notation above.
\begin{enumerate}
\item\label{itm:ch19-smooth:rank1} For $r=1$, the function $\tau$ is
 CLD-smooth if and only if the residual line bundle $Q_s$ on the abelian
 variety $\mathcal B_s$ is torsion in $\Pic(\mathcal B_s)$.
\item\label{itm:ch19-smooth:general} In general, all coordinates of $\tau$
 are CLD-smooth if and only if the residual extension of $\mathcal B_s$ by
 $T_{\widetilde K}$ is isotrivial.
\item\label{itm:ch19-smooth:finite} If the residue field $\widetilde K$ is
 algebraic over a finite field, condition
 \cref{itm:ch19-smooth:general} holds automatically.
\item\label{itm:ch19-smooth:degenerate} If $A$ is totally degenerate,
 that is $b=0$, then $\tau$ is the tropicalization of a split torus and
 is CLD-smooth; if $A$ has good reduction, that is $r=0$, then $\tau$
 is constant and the statement is vacuous.
\end{enumerate}
\end{proposition}

\begin{proof}
\cref{itm:ch19-smooth:rank1} By
\Cref{prop:ch19-tau-coordinates}\,\cref{itm:ch19-tau:torsor}, $\tau$ is the
logarithmic norm of the pulled-back tautological section for the formal
metric of $\mathcal Q$.  The exact formal-metric smoothness criterion
\cite[Proposition~17.4.13]{CLD25}, as applied in
\cite[Remark~19.2.3]{CLD25}, says that this coordinate is smooth precisely
when $\mathcal Q_s$ is torsion on every integral subvariety of
$\mathcal B_s$.  The condition on $\mathcal B_s$ itself is necessary and
implies all the restrictions, giving the stated torsion criterion.  We use
this source input rather than infer ordinary smoothness from numerical
triviality.

\cref{itm:ch19-smooth:general} Identify $T\simeq\mathbf G_m^r$ so that
$\tau=(\tau_1,\ldots,\tau_r)$.  By the previous item each $\tau_i$ is
CLD-smooth exactly when the corresponding residual bundle is torsion, and
the family of residual bundles $(Q_{i,s})$ classifies the residual extension
of $\mathcal B_s$ by $T_{\widetilde K}$; simultaneous torsion of all of them
is precisely isotriviality of that extension, meaning that the extension is
split after a finite isogeny.  In positive characteristic the isogeny
need not be \'etale.  This is
\cite[Remark~19.2.3]{CLD25}.

\cref{itm:ch19-smooth:finite} If $\widetilde K$ is algebraic over a finite
field, then $\Pic^0$ of an abelian variety over $\widetilde K$ is a torsion
group, because it is the group of $\widetilde K$-points of the dual abelian
variety and every point of an abelian variety over $\overline{\mathbf F}_p$
is torsion, the group of points over a finite field being finite.  Hence
every $Q_{i,s}$ is torsion and \cref{itm:ch19-smooth:general} applies.

\cref{itm:ch19-smooth:degenerate} If $b=0$ then $B$ is a point, $E=T$, and
$\tau$ is the ordinary tropicalization of a split torus, which is a smooth
tropicalization by \Cref{lem:ch19-toric-chart}\,(2).  If $r=0$ then
$N_{\mathbf R}=0$.
\end{proof}

\subsection{Invariant potentials and the polarization form}
\label{sec:ch19-abelian-invariant}

For a polarization, the quotient real torus carries a positive translation
invariant Lagerberg $(1,1)$-form $\omega$.  In the totally degenerate
split case, pullback by $\operatorname{trop}_A$ gives the canonical form on
$A^{\mathrm{an}}$, and invariant $\omega$-psh potentials correspond to
invariant psh potentials on the abelian variety
\cite[Proposition~7.4.5]{BGJKK25}.  Their Monge--Amp\`ere measures correspond
under tropicalization.

For an arbitrary polarized abelian variety, the good-reduction part in
\eqref{eq:raynaud-uniformization} contributes an algebraic degree factor,
while the torus rank determines the dimension of the tropical real torus.
If $L$ is ample and $\mu$ is a positive measure of mass $(L^g)$
supported on the canonical skeleton, there is a continuous semipositive
metric solving
\[
                              c_1(L,\|\cdot\|)^g=\mu,
\]
unique up to scaling; over an algebraically closed nontrivially valued field
this is \cite[Theorem~8.3.3]{BGJKK25}, and descent to an arbitrary
nontrivially valued ground field is carried out in its Section~8.4.

\begin{warning}{Canonical metrics and ordinary delta-forms}
An earlier claim that the canonical metric on every line bundle of an
abelian variety is locally a tensor product of a smooth metric and a formal
metric used the incorrect assumption that the canonical tropicalization is
always a smooth tropicalization.  Gubler--Jell--Rabinoff exhibit a mixed
reduction abelian threefold for which it is not smooth and explicitly leave
ordinary delta-metricity of general canonical metrics open
\cite[Remark~16.4 and \S17]{GJR26}.  The proved source statement is harmonic
tropicalizability.  None of the abelian comparison results above uses the
retracted ordinary-delta claim.  The new product--Gauss criterion
\Cref{thm:p09p2-product-gauss-candidate} resolves the remaining ample case
without reinstating that argument: an ample cubical metric is an ordinary
delta-metric exactly when the residual Raynaud extension is isotrivial.
Thus non-isotrivial mixed reduction gives a genuine negative answer.  The
separate good-reduction, totally degenerate, algebraically trivial, and
isotrivial cases are recorded in
\cref{rem:ch19-canonical-delta-scope}.
\end{warning}

The unrestricted question whether the canonical map from invariant Lagerberg
forms to $H^{p,q}(A)$ is always an isomorphism is answered negatively by
\Cref{prop:ch18-gjr-cohomology-counterexample}.  The counterexample has good
reduction and does not settle the totally degenerate refinement.  Of the two
other source questions formerly grouped here, ample-symmetric ordinary
delta-metricity in non-isotrivial mixed reduction is answered negatively by
\Cref{thm:p09p2-product-gauss-candidate}; the earlier
\Cref{thm:p09p2-reduction} explains why this is the exact universal range
that formerly remained.  The semipositive integral-decomposition question is
settled in Chapter~18:
\cref{thm:ch18-semipositive-integral-residual-curve} gives the
positive-cone result in residual dimension at most one,
\cref{thm:p09-signed-semipositive-span} gives signed generation in every
dimension, and
\cref{cor:p09-good-reduction-positive-cone-counterexample} disproves the
unrestricted positive-cone assertion in residual dimension two.

\begin{remark}[Exact scope of the delta-metricity question]
\label{rem:ch19-canonical-delta-scope}
\index{abelian variety!canonical metric!delta-metricity}
It is worth recording the exact status in the different reduction and
line-bundle ranges.  Let $A$ be a
polarized abelian variety with Raynaud data $(B,T,E,\Lambda)$ as above,
let $L$ be a line bundle on $A$, and give $L$ its canonical
(cubical) metric of \Cref{sec:ch19-canonical}.  By
\Cref{thm:ch19-canonical-structure} below, the pullback of that metric to
$E$ is the product of the formal metric of a line bundle $M$ on $B$
and the metric with weight $f\circ\tau$, where $f$ is a quadratic
function on $N_{\mathbf R}$.  Hence:
\begin{enumerate}
\item\label{itm:ch19-scope:good} If $A$ has \emph{good reduction}
 ($r=0$), then $N_{\mathbf R}=0$, the second factor is trivial, and the
 canonical metric is an honest formal metric, hence a delta-metric in the
 sense of \Cref{chap:delta-forms} by \cite[Proposition~8.11]{GK17}.
\item\label{itm:ch19-scope:degenerate} If $A$ is \emph{totally
 degenerate} ($b=0$), then $B$ is a point, the first factor is trivial,
 and by
 \Cref{prop:ch19-tau-smooth}\,\cref{itm:ch19-smooth:degenerate} the map
 $\tau$ is a smooth tropicalization; so $f\circ\tau$ is a CLD-smooth
 function on the skeleton and the canonical metric is smooth there.  It is a
 delta-metric, the delta-form being an ordinary smooth form.
\item\label{itm:ch19-scope:isotrivial} More generally, if the residual
 extension of $\mathcal B_s$ by $T_{\widetilde K}$ is isotrivial, then
 \Cref{prop:ch19-tau-smooth}\,\cref{itm:ch19-smooth:general} shows that all
 coordinates of $\tau$ are CLD-smooth.  Thus
 \eqref{eq:ch19-canonical-weight} expresses the canonical metric locally as
 the tensor product of a formal metric and a CLD-smooth metric, hence as an
 ordinary delta-metric.
\item\label{itm:ch19-scope:finite} If the residue field $\widetilde K$ is
 algebraic over a finite field, then by
 \Cref{prop:ch19-tau-smooth}\,\cref{itm:ch19-smooth:finite} the residual
 extension is automatically isotrivial.  This is the special case of
 \cref{itm:ch19-scope:isotrivial} that is often easiest to recognize.
\item\label{itm:ch19-scope:piczero} Under the algebraically closed,
 complete, nontrivially valued hypotheses of
 \Cref{thm:p09p2-reduction}, if $L\in\Pic^0(A)$, then every cubical metric
 on $L$ is an ordinary delta-metric.  This conclusion is independent of
 the reduction type and of isotriviality.
\item\label{itm:ch19-scope:mixed} Under the field hypotheses of
 \Cref{thm:p09p2-product-gauss-candidate}, let $L$ be ample.  Then its
 cubical metric is an ordinary delta-metric if and only if the residual
 Raynaud extension is isotrivial.  Consequently, in \emph{non-isotrivial
 mixed reduction} no cubical metric on an ample line bundle is an ordinary
 delta-metric.  In particular, this answers negatively the ample-symmetric
 question left open in \cite[Remark~16.4 and \S17]{GJR26}.  The map
 $\tau$ is still harmonic, so the Bedford--Taylor theory below remains
 available even though ordinary delta-metricity fails.
\end{enumerate}
The negative assertion in \cref{itm:ch19-scope:mixed} concerns the ordinary
algebraic-chart delta-form category.  It does not contradict harmonic
tropicalizability, which suffices for
\Cref{thm:bt-harmonic-compatibility} and hence for all the measure
computations of \Cref{sec:ch19-canonical}.
\end{remark}

\section{Canonical metrics and their currents}
\label{sec:ch19-canonical}

The canonical metric on a line bundle over an abelian variety is the
rigidity statement that makes abelian varieties computable: it is unique up
to a scalar, its curvature is translation invariant, and its top power is a
Haar measure.  This section states the construction, computes the currents,
and carries two examples to numbers.

\subsection{Cubical metrics}
\label{sec:ch19-cubical}

Let $A$ be an abelian variety over $K$, and for $I\subseteq\{1,2,3\}$
let $q_I:A^3\to A$ be the sum of the coordinates indexed by $I$.  The
theorem of the cube asserts that
\begin{equation}
\label{eq:ch19-cube-bundle}
 D_3(L)=\bigotimes_{I\subseteq\{1,2,3\}}q_I^*L^{\otimes(-1)^{\#I+1}}
\end{equation}
is trivial on $A^3$, and canonically so: there is a distinguished
trivialization $\tau_L$, compatible with pullback, tensor product and
duality, and restricting to the evident trivialization on
$A\times A\times\{0\}$.

\begin{definition}[Cubical metric]
\label{def:ch19-cubical}
\index{cubical metric}\index{canonical metric!abelian variety}
A continuous metric on $L$ is \emph{cubical}, or \emph{canonical}, if the
induced metric on $D_3(L)$ gives the canonical trivialization
$\tau_L$ constant norm one.
\end{definition}

\begin{theorem}
\label{thm:ch19-cubical-exists}
Let $A$ be an abelian variety over a complete non-Archimedean field and
let $L$ be a line bundle on $A$.
\begin{enumerate}
\item\label{itm:ch19-cub:exist} $L$ carries a cubical metric, and two
 cubical metrics on $L$ differ by a positive real scalar; fixing a
 rigidification $\varepsilon$ of $L$ at the origin, there is a unique
 cubical metric with $\|\varepsilon\|=1$.
\item\label{itm:ch19-cub:tensor} Cubical metrics are stable under tensor
 product and duality, and pull back to cubical metrics under homomorphisms
 of abelian varieties.
\item\label{itm:ch19-cub:isom} If $L$ is symmetric, respectively
 antisymmetric, and $w=2$, respectively $w=1$, fix a rigidification
 $\varepsilon$ and normalize the cubical metric by
 $\|\varepsilon\|=1$.  Then, for $n\geq2$, the canonical isomorphism
 $[n]^*L\simeq L^{n^w}$ of rigidified line bundles is an isometry.  For
 an arbitrarily scaled cubical metric the same isomorphism merely has
 constant norm.  Conversely, a metric for which one such isomorphism has
 constant norm is cubical.
\item\label{itm:ch19-cub:goodred} If $A$ extends to a formal abelian
 scheme $\mathcal A$ over $K^\circ$, then $L$ extends to a formal line
 bundle $\mathcal L$ on $\mathcal A$ and the associated formal metric is
 cubical.
\end{enumerate}
\end{theorem}

\begin{proof}
All four assertions are \cite[Lemmas~19.1.4, 19.1.6 and
Proposition~19.1.8]{CLD25}, whose proofs we summarize.

For \cref{itm:ch19-cub:exist}, a contraction argument does the work.  If
$\iota:[n]^*L\simeq L^{n^w}$ is one of the isomorphisms of
\cref{itm:ch19-cub:isom}, then the assignment
$\|\cdot\|\mapsto\|\cdot\|'$, where $\iota$ transports the pullback of the
\emph{input} metric to $L^{n^w}$ and $\|\cdot\|'$ is its $n^w$-th root,
is a contraction of the space of continuous metrics for
the uniform distance $d_\infty$ of \eqref{eq:metrics-dinfty-sup}, with
contraction factor $n^{-w}$; by
\Cref{lem:metrics-dinfty-complete} it has a unique fixed point, which is the
sought metric, and iteration from any starting metric converges uniformly to
it.  The general $L$ is handled by writing
$L^2=(L\otimes[-1]^*L)\otimes(L\otimes[-1]^*L^{-1})$ as a product of a
symmetric and an antisymmetric bundle and taking a square root.  Uniqueness
reduces, by taking the quotient of two cubical metrics, to showing that a
cubical metric on $\cO_A$ with $\|1\|=1$ is trivial, which is the same
contraction argument.

Assertion \cref{itm:ch19-cub:tensor} follows from the natural isometries
$D_3(L)\otimes D_3(M)\simeq D_3(L\otimes M)$ and
$D_3(L^{-1})\simeq D_3(L)^{-1}$, both compatible with the canonical
trivializations.

For \cref{itm:ch19-cub:goodred}, the formal line bundle $\mathcal L$ exists
and is unique up to a unit; $D_3(\mathcal L)$ is then a formal model of the
trivial bundle $D_3(L)$, its restriction to $\mathcal A^2\times\{0\}$ carries
the evident trivialization, and the unique extension of that trivialization
to $D_3(\mathcal L)$ restricts on the generic fibre to $\tau_L$ and has norm
one.
\end{proof}

The ordinary delta-metricity problem has a useful exact reduction.  Besides
settling every algebraically trivial line bundle, it shows that neither the
antisymmetric part nor an arbitrary symmetric nonample bundle is the
essential obstruction.

% Standalone exact packet for P09.2.  Do not insert into a live chapter
% before production verification.  No novelty or priority claim is intended.

For a rigidified line bundle \(L\) on an abelian variety \(A\), write

\begin{equation}
\label{eq:p09p2-d3}
 D_3(L)=
 \bigotimes_{I\subseteq\{1,2,3\}}
 q_I^*L^{\otimes(-1)^{\# I+1}},
\end{equation}

where \(q_I:A^3\to A\) adds the coordinates indexed by \(I\), and \(q_\varnothing\)
is the zero map.  The rigidification determines the distinguished
trivialization of \eqref{eq:p09p2-d3}.  A continuous metric on \(L^{\mathrm{an}}\)
is called \emph{cubical} when that trivialization has constant norm one.
In this packet, an \emph{ordinary delta-metric} always means a delta-metric
in the algebraic tropical-chart sense of
\cite[Definition~9.9 and \S9.13]{GK17}; no harmonic or weakly smooth chart is
allowed in this terminology.

\begin{theorem}
\label{thm:p09p2-reduction}
Let \(K\) be an algebraically closed, complete, nontrivially valued
non-Archimedean field, let \(A\) be an abelian variety over \(K\), and let
\(\overline L=(L,\|\cdot\|_L)\) be a cubically metrized rigidified line
bundle on \(A\).
Then the following statements hold.
\begin{enumerate}
\item\label{itm:p09p2-pic0}
 If \(L\in\operatorname{Pic}^0(A)\), then \(\overline L\) is an ordinary
 delta-metrized line bundle.

\item\label{itm:p09p2-split}
 Put \(\iota=[-1]:A\to A\) and define, with the induced rigidifications
 and the directly induced tensor, pullback, and dual metrics,
 \begin{equation}
 \label{eq:p09p2-plus-minus}
  \overline L_+
    :=\overline L\otimes\iota^*\overline L,
  \qquad
  \overline L_-
    :=\overline L\otimes(\iota^*\overline L)^\vee.
 \end{equation}
 Then both metrics in \eqref{eq:p09p2-plus-minus} are cubical,
 \begin{equation}
 \label{eq:p09p2-symmetry}
  \iota^*L_+\simeq L_+,
  \qquad
  \iota^*L_-\simeq L_-^\vee,
  \qquad
  L_-\in\operatorname{Pic}^0(A),
 \end{equation}
 and evaluation gives an isometry
 \begin{equation}
 \label{eq:p09p2-square-split}
       \overline L^{\otimes2}
       \simeq
       \overline L_+\otimes\overline L_-.
 \end{equation}
 Consequently,
 \begin{equation}
 \label{eq:p09p2-delta-equivalence}
  \overline L\text{ is an ordinary delta-metric}
  \quad\Longleftrightarrow\quad
  \overline L_+\text{ is an ordinary delta-metric}.
 \end{equation}

\item\label{itm:p09p2-ample-reduction}
 For this fixed abelian variety \(A\), the following are equivalent:
 \begin{enumerate}
 \item\label{itm:p09p2-all}
 every cubical metric on every rigidified line bundle on \(A\) is an ordinary
 delta-metric;
 \item\label{itm:p09p2-ample-symmetric}
 every cubical metric on every ample symmetric rigidified line bundle on
 \(A\) is an ordinary delta-metric.
 \end{enumerate}
\end{enumerate}
\end{theorem}


The reduction proof is given in \Cref{sec:appN-cubical-reduction}.

\begin{remark}[Two things ``canonical'' can mean]
\label{rem:ch19-canonical-two-meanings}
On an abelian variety with good reduction, the cubical metric is the formal
metric of the N\'eron model, by
\Cref{thm:ch19-cubical-exists}\,\cref{itm:ch19-cub:goodred}; this is what
``canonical'' means in Arakelov theory at a place of good reduction.  On a
general abelian variety the two notions separate: there is no distinguished
model, and the cubical metric is characterized by the cube relation, not by
any model.  \Cref{thm:ch19-canonical-structure} makes the difference
explicit by writing the cubical metric as a formal metric \emph{times} a
quadratic tropical correction.
\end{remark}

\subsection{The structure of the canonical metric}
\label{sec:ch19-canonical-structure}

We now use \eqref{eq:raynaud-uniformization}.  Keep
$p:E^{\An}\to A^{\An}$ and $q:E\to B$, and fix a line bundle $L$ on
$A$ with cubical metric normalized at the origin.  The pullback
$p^*L$ descends from $B$: there are a line bundle $M$ on $B$ and an
isomorphism $\epsilon:p^*L\simeq q^*M$, normalized to carry the
tautological trivializations at the identity to one another, and a descent
datum $c=(c_\lambda)_{\lambda\in\Lambda}$.  We write $q_L$ for the
quadratic form on $N_{\mathbf R}$ attached to $L$ by the bilinear
pairing
\begin{equation}
\label{eq:ch19-qL}
 q_L\bigl(\tau(\lambda)\bigr)=\log\bigl\|\langle\lambda,u(\lambda)\rangle\bigr\| ,
 \qquad \lambda\in\Lambda,
\end{equation}
where $u:\Lambda\to\Lambda'$ is the polarization-type map into the lattice
of the dual uniformization and $\langle\cdot,\cdot\rangle$ is the
canonical section of the Poincar\'e bundle with its formal metric; this is
the notation of \cite[\S B.6.3]{CLD25}.

\begin{theorem}
\label{thm:ch19-canonical-structure}
With the notation above, there is a unique real-valued function $f$ on
$N_{\mathbf R}$ such that
\begin{equation}
\label{eq:ch19-canonical-weight}
 -\log\|s\|+\log\|\epsilon(s)\|=f\circ\tau
\end{equation}
for every local invertible section $s$ of $p^*L$, where the first norm
is the cubical metric on $L$ pulled back and the second is the formal
metric of a formal model $\mathcal M$ of $M$ on $\mathcal B$.  The
function $f$ is quadratic, and its homogeneous part of degree two is
$\tfrac12q_L$.  Consequently
\begin{equation}
\label{eq:ch19-canonical-curvature}
 c_1(p^*\overline L)=c_1(q^*\overline M)
   +\tau^*\dprimeddprime\bigl(\tfrac12q_L\bigr),
\end{equation}
the cubical metric is dpsh, and $c_1(\overline L)$ is translation
invariant on $A^{\An}$ and vanishes if $L$ is translation invariant.
\end{theorem}

\begin{proof}
Existence and uniqueness of $f$, and quadraticity, are
\cite[Proposition~19.2.5]{CLD25}; we recall the mechanism.  The descent
datum satisfies
$\log\|c_\lambda\|=\log\|r(\lambda)\otimes r(0)^{-1}\|
 +\langle\tau(t),u(\lambda)\rangle$ at a point $t\in E$, where $r$ is the
trivialization datum of the descent, and the cocycle relation
$r(\lambda+\mu)\otimes r(\lambda)^{-1}\otimes r(\mu)^{-1}\otimes r(0)
 =\langle\lambda,u(\mu)\rangle$ shows that
$\lambda\mapsto\log\|r(\lambda)\|$ is a quadratic function on $\Lambda$ with
degree-two part $\lambda\mapsto\tfrac12\log\|\langle\lambda,u(\lambda)\rangle\|$.
Since $\tau(\Lambda)$ is a lattice in $N_{\mathbf R}$, there is a unique
quadratic $f$ on $N_{\mathbf R}$ with $\log\|r(\lambda)\|=f(\tau(\lambda))$,
and multiplying the formal metric of $q^*\mathcal M$ by $\rme^{-f\circ\tau}$
makes the descent datum have norm one; the resulting metric on $L$ is
cubical, hence equals the given one up to the normalizing scalar.
Uniqueness of $f$ follows from surjectivity of $\tau$.

For \eqref{eq:ch19-canonical-curvature}, apply
\eqref{eq:cld-pl-section-sign} to a local invertible section: the divisor
terms cancel and
$c_1(p^*\overline L)=c_1(q^*\overline M)+\dprimeddprime[f\circ\tau]$.  Since
$\tau$ is harmonic by
\Cref{prop:ch19-tau-coordinates}\,\cref{itm:ch19-tau:harmonic} and $f$ is
smooth, harmonic-chart Stokes \cite[Corollary~7.8.8]{CLD25} gives
$\dprimeddprime[f\circ\tau]=\tau^*\dprimeddprime f$, and
$\dprimeddprime f=\dprimeddprime(\tfrac12q_L)$ because the affine part of a
quadratic function is annihilated by $\dprimeddprime$.

The cubical metric is dpsh because \eqref{eq:ch19-canonical-weight} writes
it as the product of a formal metric and a $G$-smooth one, and dpsh
functions are closed under sums (\Cref{def:ch16-dpsh}); here the
formal and $G$-smooth factors are dpsh by the special abelian construction
\cite[Corollary~19.2.6]{CLD25}.  The pulled-back quadratic form is
interpreted as a $G$-form current, without asserting ordinary smoothness.
Translation
invariance: for $a\in A(K)$, the bundle $t_a^*L\otimes L^{-1}$ is
algebraically equivalent to zero, so by \cite[Proposition~B.6.7]{CLD25} its
quadratic form vanishes and its bundle $M$ on $B$ is translation invariant,
whence $\mathcal M_s$ is numerically trivial and
$c_1(\overline M)=0$ by \cite[Corollary~18.8.8]{CLD25}, applied to the
formal metric and its dual; therefore
$c_1(t_a^*\overline L)=c_1(\overline L)$.  The same computation with
$L$ itself translation invariant gives $c_1(\overline L)=0$.  This is
\cite[Corollary~19.2.6]{CLD25}.
\end{proof}

\begin{theorem}
\label{thm:ch19-canonical-positivity}
Assume the absolute value of $K$ is nontrivial and keep the notation
above.  The following are equivalent:
\begin{enumerate}
\item\label{itm:ch19-pos:psh} the cubical metric of $L$ is psh;
\item\label{itm:ch19-pos:current} $c_1(\overline L)\geq0$;
\item\label{itm:ch19-pos:data} the quadratic form $q_L$ is positive
 semidefinite and the residual line bundle $\mathcal M_s$ is nef on the
 abelian variety $\mathcal B_s$;
\item\label{itm:ch19-pos:nef} $L$ is nef.
\end{enumerate}
Moreover $c_1(\overline L)=0$ if and only if $L$ is translation
invariant, if and only if $q_L=0$ and $\mathcal M_s$ is numerically
trivial; and $L$ is ample if and only if $\mathcal M_s$ is ample and
$q_L$ is positive definite.
\end{theorem}

\begin{proof}
This is \cite[Proposition~19.3.1, Corollary~19.3.2 and
Theorem~B.6.4]{CLD25}.  The implication
\cref{itm:ch19-pos:psh}$\Rightarrow$\cref{itm:ch19-pos:current} is
\Cref{prop:cld-psh-positive-current}.  For
\cref{itm:ch19-pos:data}$\Rightarrow$\cref{itm:ch19-pos:psh}: nefness of
$\mathcal M_s$ makes $q^*\overline M$ psh by
\Cref{thm:ch17-psh-residual-nef}, since its residual bundles are pullbacks
of the nef bundle $\mathcal M_s$.  We justify separately that
$\tau^*(\tfrac12q_L)$ is psh; harmonicity by itself is not an ordinary
smooth-approximation argument.

Choose a lattice basis and write $h_1,\ldots,h_r$ for the coordinates of
$\tau$.  By \Cref{prop:ch19-tau-coordinates}, each $h_j$ is integral harmonic:
its residual bundle is numerically trivial, hence both that bundle and its
dual are nef.  Applying \Cref{thm:ch17-psh-residual-nef} to the trivial
line bundles with weights $h_j$ and $-h_j$ therefore makes both signs
continuous psh.  Consequently every real affine
function of $(h_1,\ldots,h_r)$ is psh, by the convex-cone property of the
CLD sheaf and the inclusion of constants.

Put $F=\tfrac12q_L$.  Its Hessian $H$ is a constant positive semidefinite
matrix.  For $a\in\mathbf R^r$ let
$\ell_a(t)=F(a)+\nabla F(a)\cdot(t-a)$.  Then
\[
 0\leq F(t)-\ell_a(t)
       =\tfrac12(t-a)^{\mathsf T}H(t-a)
       \leq\tfrac12\|H\|\,\|t-a\|^2.
\]
On a fixed compact subset of $\mathbf R^r$, choose a finite
$\varepsilon$-net $(a_j)$.  The convex function $\max_j\ell_{a_j}$ differs
from $F$ there by at most $\tfrac12\|H\|\varepsilon^2$.  Its pullback by
$\tau$ is a finite maximum of continuous psh functions, hence psh.
As $\varepsilon\downarrow0$, these pullbacks converge uniformly on each
fixed compact analytic neighbourhood under consideration.  Stability under
local uniform limits gives $F\circ\tau$ psh.  If $r=0$, this assertion
just says that a constant is psh.  This is the quadratic instance of the
analytic result \cite[Corollary~18.8.9]{CLD25}; it does not use the
general harmonic-chart assertion discussed in
\Cref{cor:ch16-harmonic-convex-positive}.

The quadratic function $f$ in the metric decomposition differs from $F$
by a real affine function, whose pullback is psh by the preceding
argument.  Thus $f\circ\tau$ is psh as well.  Now $p^*\overline L$ is the
product of these two psh metrics by
\eqref{eq:ch19-canonical-weight}.  Since $p$ has local sections,
$\overline L$ is psh.

For \cref{itm:ch19-pos:current}$\Rightarrow$\cref{itm:ch19-pos:data}, two
separate arguments are needed.  Over a rigid point $b\in B$ the residual
bundle of $M$ is trivial, so on a neighbourhood $U$ of $b$ the metrized
bundle $q^*\overline M$ is trivial and $E$ is trivialized, whence
$c_1(p^*\overline L)=\tau^*(\tfrac12\dprimeddprime q_L)$ there; positivity
of a form pulled back by a surjection onto $N_{\mathbf R}$ forces
$\dprimeddprime q_L\geq0$, that is, $q_L$ positive semidefinite.  For the
residual nefness, let $\omega$ be the standard invariant $(1,1)$-form on
$N_{\mathbf R}$; then $\tau^*\omega^r$ kills the second term of
\eqref{eq:ch19-canonical-curvature} for dimension reasons, so
$c_1(q^*\overline M)\wedge\tau^*\omega^r\geq0$, and the fibre-integration
projection formula \Cref{thm:cld-fibre-projection} converts this into
$\langle c_1(\overline M),\alpha\rangle\geq0$ for every positive smooth
$\alpha$ on $B$, hence $c_1(\overline M)\geq0$ and $\mathcal M_s$ is nef by
\Cref{thm:ch17-psh-residual-nef}.

Equivalence with \cref{itm:ch19-pos:nef} in the algebraic case: if $L$ is nef,
split $L^2=L_+\otimes L_-$ with $L_+$ symmetric and $L_-$ antisymmetric, so
$c_1(\overline{L_-})=0$; for an ample symmetric $N$ the bundle
$L_+^{t}\otimes N$ is ample for all $t\geq0$ by Nakai--Moishezon and its
cubical metric is psh, whence $tc_1(\overline{L_+})+c_1(\overline N)\geq0$
and, dividing by $t\to\infty$, $c_1(\overline{L_+})\geq0$.  Conversely if
the metric is psh then for every irreducible curve $C\subset A$ one has
$(L\cdot C)=\int_{C^{\An}}c_1(\overline L|_{C^{\An}})\geq0$ by
\Cref{thm:ch17-degree}, so $L$ is nef.  The statements about
triviality and ampleness are the cited corollary and theorem.
\end{proof}

\subsection{Ordinary delta-metricity}
\label{sec:ch19-canonical-delta-metricity}

We can now decide exactly when the quadratic factor in
\eqref{eq:ch19-canonical-weight} belongs to the ordinary
Gubler--K\"unnemann delta-form category.  The point is not merely that the
coordinates of $\tau$ fail to be smooth.  One must recover their first
derivatives from the single ambient codimension-zero coefficient tensor
required in every ordinary delta-chart.  The torus action supplies the
missing transverse directions, while relative Gauss points specialize the
result back to the original analytic space.  This yields both a general
recognition criterion and the exact answer for ample cubical metrics.

Three chart bridges should be kept distinct.  A finite integral-PL map
factors locally through one analytic moment map; after shrinking, Gubler's
algebraization theorem replaces that finite analytic moment by an algebraic
one with the same tropicalization; and a canonical very affine moment then
dominates the finitely many algebraic charts under consideration.  Equality
of tropicalizations, not equality of the analytic functions, is the input
needed below.  In the proof, for a cell $\sigma$ we write
$L_\sigma=\operatorname{span}_{\mathbf R}(\sigma-\sigma)$,
$\operatorname{Star}(C)$ means the incident cells containing $C$, and
$D_\sigma H$ is the linear part of $H|_\sigma$.  An ordinary
neighbourhood is a neighbourhood for the Berkovich topology.  A relative
Gauss point over $x$ is formed after base change to $\mathcal H(x)$; it
is not a naive ordered pair of Berkovich points.

The proof proceeds by algebraic action pullback, a relative product--Gauss
path, a common-cell lineality argument, contraction of the single ambient
Hessian tensor, and specialization back to the original point.  The converse
uses algebraization to pass from CLD analytic-moment smoothness to ordinary
algebraic-chart smoothness.  Finally, the restricted analytic-local
invariance permits passage through the Raynaud covering; ordinary algebraic
pullback alone would not suffice.

% Exact candidate statement for source and production-Rethlas audit.
% This file is isolated from the live book. It makes no novelty or priority
% claim and must not be promoted before a final/correct production receipt.

\begin{theorem}
\label{thm:p09p2-product-gauss-candidate}
Let \(K\) be an algebraically closed complete non-Archimedean field with
nontrivial absolute value.  Let \(T\) be a split algebraic torus, put

\[
 M=X^*(T),\qquad N=\operatorname{Hom}(M,\mathbf Z),
\]

and let \(X\) be an algebraic \(K\)-variety endowed with an algebraic
\(T\)-action

\[
                     \mu:T\times X\longrightarrow X.
\]

Suppose that \(h:X^{\mathrm{an}}\to N_{\mathbf R}\) is locally integral
piecewise affine and that, as maps on \((T\times X)^{\mathrm{an}}\),

\begin{equation}
\label{eq:p09p2-product-equivariance}
 h\circ\mu^{\mathrm{an}}
 =\operatorname{trop}_T\circ\operatorname{pr}_T
  +h\circ\operatorname{pr}_X.
\end{equation}

Let \(b\) be a nondegenerate symmetric bilinear form on
\(N_{\mathbf R}\), let \(\ell\in M_{\mathbf R}=N_{\mathbf R}^{\vee}\), and
let \(c\in\mathbf R\).  Set

\[
                  F(w)=\frac12 b(w,w)+\ell(w)+c.
\]

Endow the trivial line bundle on \(X\) with the metric defined by

\[
                         -\log\lVert1\rVert=F\circ h.
\]

This metric is an ordinary Gubler--K\"unnemann delta-metric in the
algebraic tropical-chart sense of \cite[Definition~9.9]{GK17} if and only
if every coordinate \(\langle m,h\rangle\), \(m\in M\), is smooth in the
sense of Chambert-Loir--Ducros.

Now let \(A\) be an abelian variety over \(K\), with split Raynaud
uniformization

\[
  0\longrightarrow T\longrightarrow E\xrightarrow{q}B\longrightarrow0,
  \qquad
  A^{\mathrm{an}}=E^{\mathrm{an}}/\Lambda,
\]

and let \(L\) be an ample rigidified line bundle on \(A\), endowed with its
normalized cubical metric.  This metric is an ordinary
Gubler--K\"unnemann delta-metric if and only if the residual Raynaud
extension of \(B_s\) by \(T_{\widetilde K}\) is isotrivial.  Every cubical
metric on \(L\) has the same ordinary delta-metricity status.  Consequently,
on an abelian variety with non-isotrivial mixed reduction, no cubical metric
on an ample line bundle---in particular, on an ample symmetric rigidified
line bundle---is an ordinary delta-metric.
\end{theorem}


The product--Gauss proof is given in \Cref{sec:appN-product-gauss}.

\subsection{The canonical measure}
\label{sec:ch19-canonical-measure}

\begin{theorem}
\label{thm:ch19-canonical-ma}
Let $A$ be a $g$-dimensional abelian variety over a complete
nontrivially valued field with split Raynaud data, and let $L$ be an ample
line bundle with its cubical metric.  Then:
\begin{enumerate}
\item\label{itm:ch19-ma:E} $c_1(p^*\overline L)^g$ is a Haar measure on
 $S(E)\simeq N_{\mathbf R}$;
\item\label{itm:ch19-ma:A} $c_1(\overline L)^g$ is the unique Haar
 measure on the skeleton $S(A)\simeq N_{\mathbf R}/\tau(\Lambda)$ of total
 mass $(L^g)$;
\item\label{itm:ch19-ma:formula} explicitly, writing
 $\omega_L$ for the unique $G$-form on $A^{\An}$ with
 $p^*\omega_L=\tau^*\dprimeddprime(\tfrac12q_L)$,
 \[
   c_1(\overline L)^g=\binom{g}{r}\,S_b\wedge\omega_L^{\,r},
 \]
 where $S_b$ is the descent of $c_1(q^*\overline M)^b$ and
 $b=g-r$.
\end{enumerate}
\end{theorem}

\begin{proof}
This is \cite[Corollary~19.2.8]{CLD25}; we give the computation, since the
bookkeeping of which terms survive is exactly what the reader needs in the
examples below.

\emph{Step 1: expanding the binomial.}  By
\eqref{eq:ch19-canonical-curvature}, $c_1(p^*\overline L)$ is the sum of the
closed current $c_1(q^*\overline M)$ and the $G$-form
$\widetilde\omega_L:=\tau^*\dprimeddprime(\tfrac12q_L)$, which is locally
harmonically tropicalizable because $\tau$ is harmonic
(\Cref{prop:ch19-tau-coordinates}).  Both are obtained by applying
$\dprimeddprime$ to dpsh weights, so the Bedford--Taylor product is
multilinear by \Cref{thm:ch16-dpsh-product}, and
\[
 c_1(p^*\overline L)^g
 =\sum_{d=0}^{g}\binom{g}{d}
   c_1(q^*\overline M)^{g-d}\wedge\widetilde\omega_L^{\,d}.
\]
Each formal curvature power is a parlinear current and a difference of
skeletal Bedford--Taylor currents.  Thus the hypotheses of
\Cref{thm:bt-harmonic-compatibility} apply to its product with
$\widetilde\omega_L^{\,d}$, giving the direct $G$-form pairing used here.
This is the special product justified in \cite[Corollary~19.2.8]{CLD25},
not a multiplication rule for an arbitrary current and an arbitrary
weakly smooth form.

\emph{Step 2: only one term survives.}  The form
$\widetilde\omega_L^{\,d}$ vanishes for $d>r$, being pulled back from
$N_{\mathbf R}$, which has dimension $r$.  The current
$c_1(q^*\overline M)^{g-d}$ vanishes for $g-d>b$, because $q^*\overline M$
has dimension of definition at most $b=\dim B$ at every point
(\Cref{thm:ch14-support-dimension}).  Since $g=b+r$, both constraints are
equalities only for $d=r$, and
\[
 c_1(p^*\overline L)^g
 =\binom{g}{r}c_1(q^*\overline M)^{b}\wedge\widetilde\omega_L^{\,r}.
\]

\emph{Step 3: the fibre computation.}  Let $\beta$ be the Gauss point of
$B$, that is, the divisorial point over the generic point of
$\mathcal B_s$.  Since $\mathcal B_s$ is irreducible of dimension $b$,
the formal-intersection and relative fibre-integration calculation of
\cite[Proposition~18.6.1 and Corollary~19.2.8]{CLD25} gives a current
$c$ times integration along $q^{-1}(\beta)$, where
$c=(\mathcal M_s^b)>0$.  The fibre is an analytic torus over
$\mathcal H(\beta)$; it is not an ordinary closed analytic $K$-cycle.
This integration current is the source's relative construction, not an
arbitrary pullback of a Dirac measure.  On the fibre $\tau$ is the usual
tropicalization, and for a test function $u$,
\[
 \bigl\langle c_1(p^*\overline L)^g,u\bigr\rangle
 =c\binom{g}{r}\int_{q^{-1}(\beta)}u\,\tau^*\bigl(\dprimeddprime\tfrac12q_L\bigr)^r ,
\]
and $\bigl(\dprimeddprime\tfrac12q_L\bigr)^r$ is a translation-invariant
$(r,r)$-form on $N_{\mathbf R}$ because $q_L$ is quadratic.  Contracting
against the canonical calibration of $S(E)$ turns this into a
translation-invariant $r$-form, that is, a multiple of Lebesgue measure.
This proves \cref{itm:ch19-ma:E}.

\emph{Step 4: descent and mass.}  The map $p$ is a covering, so currents
pull back along it and $p^*c_1(\overline L)^g=c_1(p^*\overline L)^g$.  A
Haar measure on $S(E)$ invariant under $\tau(\Lambda)$ descends to a Haar
measure on $S(A)$; equality of the two pullbacks gives equality of the
descended currents.  The total mass is $(L^g)$ by
\Cref{thm:ch17-degree}, and a Haar measure on a compact real torus
is determined by its mass.  This is \cref{itm:ch19-ma:A}, and Step~2 with the
descent of $\widetilde\omega_L$ to $\omega_L$ is
\cref{itm:ch19-ma:formula}.
\end{proof}

\begin{proposition}
\label{prop:ch19-canonical-subvariety}
Let $U\subseteq A^{\An}$ be open, let $X\subseteq U$ be a closed
irreducible analytic subspace of dimension $d$, and let $L$ be a line
bundle on $A$ with its cubical metric.  For $j\in\mathbf N$, the
current $c_1(q^*\overline M|_{p^{-1}(X)})^j$ descends to a current
$S_j$ on $X$ of bidimension $(d-j,d-j)$, and
\begin{equation}
\label{eq:ch19-canonical-subvariety}
 c_1\bigl(\overline L|_X\bigr)^{e}
 =\sum_{j=\max(e-b,0)}^{\min(r,e)}\binom{e}{j}S_{e-j}\wedge\omega_L^{\,j},
 \qquad 0\leq e\leq d .
\end{equation}
In particular $c_1(\overline L|_X)^d$ is a finite sum of measures
$\mu_j$, each carried by a calibrated polyhedral subspace $P_j$ of
dimension $j$ on which $\tau$ is a piecewise immersion, and each locally
a constant multiple of Lebesgue measure on $P_j$.
\end{proposition}

\begin{proof}
Expand $c_1(p^*\overline L|_{p^{-1}X})^{e}$ as in Step~1 above and use the
two vanishing statements of Step~2, now with $\dim X=d$ in place of $g$: the
term of index $j$ vanishes if $e-j>b$ or if $e-j>d$ or if $j>r$.  Each
$c_1(q^*\overline M|_{p^{-1}X})^{i}$ is a difference of Bedford--Taylor
currents and descends because $c_1(p^*\overline L|_{p^{-1}X})$ and
$\widetilde\omega_L$ do.  The calibrated spaces $P_j$ are supplied by the parlinear-current
pairing in \cite[Corollary~19.2.9]{CLD25}, with its special formal curvature
powers and harmonic-chart $G$-forms.  The measures
$\omega_L^{\,j}|_{P_j}$ are locally constant multiples of Lebesgue measure
because $\omega_L$ is pulled back by piecewise linear $\tau$ from a
translation-invariant form \cite[Remark~19.3.4]{CLD25}.  No analogous
polyhedral decomposition is inferred for arbitrary closed currents.
\end{proof}

\begin{remark}[All the pieces can be nonzero]
\label{rem:ch19-canonical-pieces}
The range of $j$ in \eqref{eq:ch19-canonical-subvariety} is not merely a
bookkeeping device.  Gubler has constructed a subvariety of an abelian
variety for which none of the measures $\mu_j$ vanishes, for any
$j$ in the admissible range \cite[Remark~19.3.4\,c)]{CLD25}.  Moreover, if
$L$ and $L'$ are both ample, the supports of $\mu_j(L)$ and
$\mu_j(L')$ coincide for every $j$: choosing $p\geq1$ with
$L^p\otimes(L')^{-1}$ ample gives
$p\,c_1(\overline L|_X)\geq c_1(\overline{L'}|_X)$, so each measure
dominates a multiple of the other, and the polyhedral decomposition into
pieces of distinct dimensions forces the supports to agree piece by piece.
Thus the polyhedral decomposition is an invariant of $X$, not of the
polarization.
\end{remark}

\subsection{The Tate curve, completely}
\label{sec:ch19-tate-canonical}

\begin{example}[Canonical metric and canonical measure on a Tate curve]
\label{ex:ch19-tate-canonical}
\index{Tate curve!canonical metric}
Assume here that $K$ is algebraically closed, complete, and nontrivially
valued.  Let $q\in K^\times$ with $\ell=v(q)>0$, let
$E_q=\mathbf G_{m,K}^{\An}/q^{\mathbf Z}$, and keep $\theta$,
the real Gauss profile $V(t)=v_{\eta_t}(\theta)$, and $\tau(z)=v(z)$ from
\Cref{ex:ch19-tate-slope}.  Here $g=r=1$, $b=0$, $B$ is a point,
$T=\mathbf G_m$, $E=\mathbf G_m$, $\Lambda=q^{\mathbf Z}$, and
$\tau(\Lambda)=\ell\mathbf Z$; the skeleton is
$S(E_q)=\mathbf R/\ell\mathbf Z$, a circle of circumference $\ell$.

\emph{The line bundle.}  Let $L=\cO_{E_q}(d\cdot(O))$, of degree $d>0$,
where $O$ is the image of $z=1$.  Then
$p^*L=\cO_{\mathbf G_m}(D)$ with $D=d\sum_{n\in\mathbf Z}(q^n)$, and the
tautological section $s_D$ has $\divis(s_D)=D$.  A metric on
$p^*L$ is determined by the weight $\phi=-\log\|s_D\|$.

\emph{The functional equation.}  Because $p\circ t_q=p$, the canonical
identification $t_q^*p^*L=p^*L$ preserves the tautological divisor section,
and a metric descends to $E_q$ exactly when its weight is $q$-periodic:
$\phi(qz)=\phi(z)$.  Set $W(z)=-\log|\theta(z)|$, so that
$W(\eta_t)=V(t)$, and construct the candidate weight
\[
 \phi(z)=dW(z)+g(\tau(z)).
\]
Here $W$ is an analytic logarithmic function, whereas $V$ is only its
real skeletal profile.  This ansatz uses a tropical correction to the
specified divisor metric; the polar behaviour alone would not force an
arbitrary metric ratio to factor through $\tau$.
By \eqref{eq:ch19-theta-functional},
$W(qz)=W(z)-\tau(z)$, so periodicity becomes
\begin{equation}
\label{eq:ch19-tate-cocycle}
 g(t+\ell)-g(t)=d\,t,\qquad t\in\mathbf R .
\end{equation}
Its solutions differ by $\ell$-periodic functions; the \emph{quadratic}
solution is unique up to a constant, namely
\begin{equation}
\label{eq:ch19-tate-g}
 g(t)=\frac{d}{2\ell}t^2-\frac{d}{2}t
     =\frac{d}{2\ell}\,t\,(t-\ell),
\end{equation}
as one checks:
$g(t+\ell)-g(t)=\frac{d}{2\ell}(2t\ell+\ell^2)-\frac d2\ell=dt$.
By \Cref{thm:ch19-canonical-structure} this is the canonical metric:
$f=g$, and its degree-two part is $\tfrac12q_L$ with
\begin{equation}
\label{eq:ch19-tate-qL}
 q_L(t)=\frac d\ell t^2,
 \qquad
 q_L(\tau(q))=q_L(\ell)=d\ell=\log\bigl\|\langle q,u(q)\rangle\bigr\| ,
\end{equation}
the last equality because the polarization attached to $d\cdot(O)$ is
$d$ times the principal one, whose Poincar\'e pairing on
$\Lambda=q^{\mathbf Z}$ is $v(q)=\ell$.  This is the check of
\eqref{eq:ch19-qL}.

\emph{The weight on the skeleton.}  Since $V\equiv0$ on $[0,\ell]$, the
canonical weight restricted to the fundamental domain is simply
\begin{equation}
\label{eq:ch19-tate-weight}
 \phi(\eta_t)=\frac{d}{2\ell}\,t(t-\ell),
 \qquad 0\leq t\leq\ell,
\end{equation}
a convex quadratic function of $t$, with minimum $-d\ell/8$ at
$t=\ell/2$ and
$\phi(0)=\phi(\ell)=0$.  Up to the additive constant $d\ell/12$ this is
$\tfrac{d\ell}{2}B_2(t/\ell)$ with $B_2(x)=x^2-x+\tfrac16$ the second
Bernoulli polynomial; it is the non-Archimedean N\'eron function of
$d\cdot(O)$, and it is \eqref{eq:ch08-green-circle} integrated twice.

\emph{The curvature.}  On $\mathbf G_m^{\An}$,
\[
 c_1(p^*\overline L)
 =\dprimeddprime[\phi]+\delta_D
 =\underbrace{-\delta_D}_{\dprimeddprime[dW]}
   +\tau^*\dprimeddprime g+\delta_D
 =\tau^*\dprimeddprime g,
\]
using \Cref{thm:cld-poincare-lelong} for the first term and harmonicity of
$\tau$ for the second.  Since $g''=d/\ell$,
\begin{equation}
\label{eq:ch19-tate-measure}
 c_1(\overline L)=\frac d\ell\,\rd t
 \quad\text{on }S(E_q)=\mathbf R/\ell\mathbf Z,
 \qquad\text{total mass }\frac d\ell\cdot\ell=d=\deg L .
\end{equation}
This is \Cref{thm:ch19-canonical-ma} in the case $g=r=1$: a Haar measure
on the circle of the correct mass.  It can also be read off directly, since
on the circle
\[
 \Delta_\Gamma\bigl(\phi|_\Gamma\bigr)
 =\frac d\ell\,\rd t-d\,\delta_{t=0},
\]
the atom at $t=0$ coming from the slope jump
$g'(0^+)-g'(\ell^-)=-\tfrac d2-\tfrac d2=-d$, and adding
$\tau_*\delta_{d\cdot(O)}=d\delta_{t=0}$ restores
\eqref{eq:ch19-tate-measure}.

\emph{Comparison with the model metrics of \Cref{chap:model-calculations}.}
Subdivide the circle into $N$ arcs of equal length $h=\ell/N$ with
endpoints $t_i=ih$, and let $\phi_N$ be the piecewise affine
interpolation of $\phi|_\Gamma$ at the $t_i$, extended periodically to
$\mathbf R$.  To define its full analytic metric, put
$g_N(t)=\phi_N(t)-dV(t)$ and use the weight
$dW(z)+g_N(\tau(z))$.  Its cocycle is the same as above.  Its rational
slopes and value-group break points give a rational semistable model metric;
our algebraically closed field supplies the required subdivisions.
The symbol $\overline L_{\phi_N}$ below denotes this metric, not a
metric obtained by treating $V$ as an analytic function.
For $i\neq0$ the slope
jump of $\phi_N$ at $t_i$ is
\[
 \frac{\phi(t_{i+1})-2\phi(t_i)+\phi(t_{i-1})}{h}
 =\frac{g''h^2}{h}=\frac d\ell\,h=\frac dN ,
\]
exactly, since $g$ is a quadratic polynomial; and at $t_0=0$ it is
\[
 \frac{g(h)+g(\ell-h)}{h}
 =\frac{d}{2\ell}\bigl[(h-\ell)-(\ell-h)\bigr]
 =\frac dN-d .
\]
Adding $d\delta_{t=0}$ gives
\begin{equation}
\label{eq:ch19-tate-model-approx}
 c_1\bigl(\overline L_{\phi_N}\bigr)
 =\frac dN\sum_{i=0}^{N-1}\delta_{t_i},
 \qquad\text{mass }d .
\end{equation}
For $\ell=N=d=3$ this is
$\delta_{v_1}+\delta_{v_2}+\delta_{v_3}$, which is exactly the balanced
model metric of \Cref{ex:ch08-tate-curve}.  As $N\to\infty$ the measures
\eqref{eq:ch19-tate-model-approx} converge weakly to
\eqref{eq:ch19-tate-measure}; and since $\phi_N\geq\phi$ with
$\|\phi_N-\phi\|_\infty\leq\tfrac{d}{8\ell}h^2\to0$, the convergence is
uniform.  To obtain monotonicity, choose a nested sequence of subdivisions,
for example $N_k=N_0 2^k$.  Convexity then gives
$\phi_{N_{k+1}}\leq\phi_{N_k}$, so their full analytic weights decrease
uniformly to the canonical weight (the metric norms increase).  Thus the
canonical metric is a uniform limit of semipositive model metrics, hence a
continuous semipositive metric in the sense of
\Cref{sec:metrics-semipositive}.  Every step is consistent with
\Cref{thm:cld-bt-continuity}.

\emph{Three checks recorded.}  The mass $d$; the identification
$q_L(\ell)=d\ell$ with the Poincar\'e pairing; and the exact agreement of
\eqref{eq:ch19-tate-model-approx} with the model computation of
\Cref{chap:model-calculations}.
\end{example}

\subsection{An abelian surface}
\label{sec:ch19-abelian-surface}

\begin{example}[A totally degenerate abelian surface]
\label{ex:ch19-abelian-surface-degenerate}
\index{abelian variety!abelian surface}
Let $K$ be algebraically closed, complete, nontrivially valued, and let
$q\in K^\times$ with $v(q)=1$.  Take $T=\mathbf G_m^2$,
$N=\mathbf Z^2$, $B=\Spec K$, so $E=T$, $b=0$, $r=g=2$: totally
degenerate reduction.  Let
\[
 \Lambda=\langle\lambda_1,\lambda_2\rangle,
 \qquad \lambda_1=(q,1),\quad \lambda_2=(1,q),
\]
so that $\tau(\lambda_1)=(1,0)$, $\tau(\lambda_2)=(0,1)$, and
$\tau(\Lambda)=\mathbf Z^2=N$.  The skeleton is the real torus
$S(A)=\mathbf R^2/\mathbf Z^2$, of lattice-normalized volume $1$.

\emph{The polarization data.}  Choose the symmetric bilinear form
\[
 b(\lambda_i,\lambda_j)=B_{ij},
 \qquad
 B=\begin{pmatrix}2&1\\1&2\end{pmatrix},
\]
positive definite with $\det B=3$.  By \cite[Theorem~B.6.4]{CLD25} the
data $(\Lambda,u)$ with this pairing define an ample line bundle $L$ on
the abelian surface $A=E^{\An}/\Lambda$, and by \eqref{eq:ch19-qL} the
quadratic form on $N_{\mathbf R}=\mathbf R^2$ is
\[
 q_L(x)=x^{\mathsf T}Bx,
 \qquad
 \tfrac12q_L(x)=\tfrac12x^{\mathsf T}Bx .
\]

\emph{The canonical measure.}  Here $M$ is the trivial bundle on a point,
so $S_0=1$ and \Cref{thm:ch19-canonical-ma}\,\cref{itm:ch19-ma:formula}
reads $c_1(\overline L)^2=\binom22\omega_L^2=\omega_L^2$.  By
\eqref{eq:real-ma-density},
\[
 \Bigl(\dprimeddprime\tfrac12q_L\Bigr)^2
 =2!\det\bigl(D^2(\tfrac12q_L)\bigr)\,\rd x
 =2\det(B)\,\rd x
 =6\,\rd x ,
\]
so
\begin{equation}
\label{eq:ch19-abelian-surface-measure}
 c_1(\overline L)^2=6\,\rd x
 \quad\text{on }\mathbf R^2/\mathbf Z^2,
 \qquad\text{total mass }6\cdot1=6 .
\end{equation}

\emph{The check against the polarization type.}  For a polarization of type
$(d_1,d_2)$ one has $\chi(L)=d_1d_2$ and $(L^2)=2!\,d_1d_2$.  The
elementary divisors of $B$ over $\mathbf Z$ are $1$ and $3$, since
$\gcd$ of the entries is $1$ and $\det B=3$; hence $L$ has type
$(1,3)$, $\chi(L)=3$, and $(L^2)=2\cdot3=6$, in agreement with
\eqref{eq:ch19-abelian-surface-measure}.  Had we taken $B=I$, the
principal polarization, the measure would be $2\,\rd x$ of mass $2$,
which is $(L^2)=2!\cdot1$ for a principally polarized abelian surface.

\emph{A subvariety.}  Let $X\subset A$ be a curve, so $d=1$.  In
\eqref{eq:ch19-canonical-subvariety} with $b=0$ only $j=e=1$ survives,
so $c_1(\overline L|_X)$ is carried by the calibrated analytic polyhedral
space in \Cref{prop:ch19-canonical-subvariety}.  Its tropical pushforward
is the trace of $\omega_L$ on $\trop_A(X)$ with its tropical
multiplicities, and its total mass is $(L\cdot X)$.  Suppose
$\trop_A(X)$ is the image of the line $\mathbf Rw$ with
$w\in\mathbf Z^2$ primitive, carrying tropical multiplicity $m$.
Parametrize that line by lattice arclength $s$, so the point is $sw$ and
the closed geodesic has lattice length $1$.  Then
\[
 \tfrac12q_L(sw)=\tfrac12s^2\,w^{\mathsf T}Bw,
 \qquad
 \frac{\rd^2}{\rd s^2}\Bigl(\tfrac12q_L(sw)\Bigr)=w^{\mathsf T}Bw,
\]
so the density of $\omega_L$ along the geodesic is the constant
$w^{\mathsf T}Bw$ and
\[
 (L\cdot X)=m\,w^{\mathsf T}Bw .
\]
For $w=(1,0)$ and $m=1$ this gives $B_{11}=2$; for $w=(1,1)$ and
$m=1$ it gives $2+2\cdot1+2=6$.
\end{example}

\begin{example}[An abelian surface of mixed reduction]
\label{ex:ch19-abelian-surface-mixed}
Now take $g=2$ with $b=r=1$: mixed reduction, the case singled out in
\Cref{rem:ch19-canonical-delta-scope}.  Let $B$ be an elliptic curve over
$K$ with good reduction and formal model $\mathcal B$, let
$T=\mathbf G_m$, and let $E$ be an extension of $B$ by $T$
corresponding to a point $Q\in\Pic^0(B)=B^\vee(K)$.  Let
$\Lambda=\lambda^{\mathbf Z}$ with $\tau(\lambda)=\ell>0$, so that the
skeleton $S(A)=\mathbf R/\ell\mathbf Z$ is a circle of circumference
$\ell$.  Assume the polarization data are chosen so that
$\mathcal M_s$ has degree $m>0$ on the elliptic curve $\mathcal B_s$
and $q_L(t)=ct^2$ with $c>0$.  These are assumed to be genuine compatible
Raynaud descent data defining $L$, not arbitrary independently chosen
numbers and extension classes.  The ampleness criterion is
\Cref{thm:ch19-canonical-positivity}, from
\cite[Theorem~B.6.4 and Remark~B.6.6]{CLD25}.

\emph{The measure.}  Here $b=r=1$, so
\Cref{thm:ch19-canonical-ma}\,\cref{itm:ch19-ma:formula} gives
\[
 c_1(\overline L)^2=\binom21S_1\wedge\omega_L=2\,S_1\wedge\omega_L,
\]
where $S_1$ is the descent of $c_1(q^*\overline M)$.  By
\Cref{thm:cld-formal-intersection}, $c_1(\overline M)=m\,\delta_\beta$
with $\beta$ the Gauss point of $B$.  The relative fibre-integration
calculation in the proof of \Cref{thm:ch19-canonical-ma} then identifies
$S_1$ with the descent of $m$ times that relative current over
$q^{-1}(\beta)$; its torus skeleton descends to the whole circle.  Since
$\dprimeddprime(\tfrac12q_L)=c\,\dprime t\wedge\ddprime t$, we get
\[
 c_1(\overline L)^2=2mc\,\rd t
 \quad\text{on }\mathbf R/\ell\mathbf Z,
 \qquad\text{total mass }(L^2)=2mc\ell .
\]
With $\ell=1$, $c=1$, $m=3$ this is $(L^2)=6$, so $\chi(L)=3$ and
$L$ has type $(1,3)$, the same numerical polarization as in
\Cref{ex:ch19-abelian-surface-degenerate}; the two surfaces have the same
degree and totally different skeleta, of dimensions $2$ and $1$.  The
canonical measure remembers the reduction type through the \emph{dimension}
of its support, not through its mass.

\emph{Where ordinary delta-metricity fails.}  The metric here is, by
\eqref{eq:ch19-canonical-weight}, the product of the formal metric of
$\mathcal M$ and the weight $f\circ\tau$.  By
\Cref{prop:ch19-tau-smooth}\,\cref{itm:ch19-smooth:rank1}, $\tau$ is
CLD-smooth if and only if $Q_s\in\Pic^0(\mathcal B_s)$ is torsion.  If
$\widetilde K$ is algebraic over a finite field this is automatic; for a
general residue field it is a genuine condition on the extension class, and
when it fails \Cref{thm:p09p2-product-gauss-candidate} proves that the ample
cubical metric is not an ordinary delta-metric.  Thus a non-torsion residual
class gives an explicit mixed-reduction example that is harmonically
tropicalizable and psh but not ordinary delta.  All the measure computations
above remain valid, because their formal curvature powers and
harmonic-chart $G$-forms satisfy the hypotheses of
\Cref{thm:bt-harmonic-compatibility}.  This example therefore separates the
harmonic Bedford--Taylor category from the ordinary algebraic-chart
delta-form category in the smallest mixed dimension.
\end{example}

\section{The local--global common core}
\label{sec:ch19-common-core}

Curves, toric varieties and abelian varieties have now each been computed in
several languages, and each time the languages agreed.  The remaining two
sections say precisely why, and precisely how far the agreement extends.
This one packages the domain on which the local approaches of
\Cref{part:cld} coincide; \Cref{sec:ch19-cld-bfj} then compares that domain
with the global theory of \Cref{part:bfj}.

The following theorem packages the domain on which the principal
approaches genuinely coincide.

\begin{theorem}\label{thm:common-model-core}
Let $K$ be algebraically closed, complete, and nontrivially valued, let
$X$ be a proper algebraic variety, and let $\overline L$ be a formal
metric.  Then the following are equivalent:
\begin{enumerate}
\item the residual formal line bundle is nef;
\item the metric is semipositive in the local model sense;
\item its Chern delta-form is positive;
\item its local weights are functorially psh;
\item its restrictions to closed algebraic curves are psh.
\end{enumerate}
If the local weights are viewed as $\mathbf R$-PL functions, these
conditions are also equivalent to classical psh under the hypotheses of
\Cref{thm:gr-classical-pl}.  For $n$ such formal metrics, the CLD
curvature product, the delta-form wedge product, and the Chambert--Loir
special-fibre measure agree.
\end{theorem}
\begin{proof}
The first five equivalences are
\Cref{thm:gk-formal-positivity-comparison}.  The equivalence with
classical psh is \Cref{thm:gr-classical-pl}, applied in local
trivializations.  Equality of measures follows respectively from
the general formal intersection formula in
\Cref{thm:cld-formal-intersection} and from
\Cref{thm:gk-delta-cl-measure}.
\end{proof}

Over a discretely valued field of residue characteristic zero, on a smooth
projective variety, this model core is also the starting locus of the BFJ
global theory: semipositivity means relative nefness of the model line
bundle, and the BFJ model Monge--Amp\`ere measure is the same
special-fibre measure.  The assertion is deliberately confined to model
metrics.  The later BFJ extension to singular psh functions uses global
compactness, envelopes, and energy, not the local CLD definition alone.

\section{A first CLD--BFJ comparison}
\label{sec:ch19-cld-bfj}

\Cref{thm:common-model-core} compares notions that are all local and all
defined inside \Cref{part:cld}.  We now make the first comparison with a
theory that is not: the global model theory of Boucksom--Favre--Jonsson,
constructed in \Cref{part:bfj}.  This is the theorem that
\Cref{part:comparisons} will cite whenever a statement has to be transported
across the two frameworks, so its scope is stated with more care than its
proof requires.

\subsection{The two frameworks side by side}
\label{sec:ch19-cld-bfj-setup}

Throughout this section, and \emph{only} in this section, we adopt the
original BFJ hypotheses of \Cref{chap:bfj-model-forms}:
\begin{equation}
\label{eq:ch19-bfj-hypotheses}
\left\{
\begin{aligned}
 &K\ \text{complete discretely valued, residue characteristic }0,\
   v(\pi)=1;\\
 &X\ \text{smooth, connected, projective over }K,\ \dim X=n;\\
 &\text{models normal, flat and projective over }R=K^\circ .
\end{aligned}
\right.
\end{equation}
These are needed on the BFJ side for resolution of singularities and for the
model-envelope machinery; nothing in \Cref{part:cld} requires them.  Where a
statement below survives without \eqref{eq:ch19-bfj-hypotheses}, we say so.
Set $F=\widehat{K^a}$.  The ordinary Gubler--K\"unnemann delta-form and
functorial curve tests below are made on $X_F$, component by component.
The BFJ model data and CLD measures remain over the discretely valued $K$;
these two field settings must not be identified.

The two frameworks assign different names to the same underlying data, and
the first task is to align them.

\begin{center}
\small
\begin{tabular}{p{.25\textwidth}|p{.29\textwidth}|p{.30\textwidth}}
\textbf{Datum}&\textbf{CLD name (\Cref{part:cld})}
 &\textbf{BFJ name (\Cref{part:bfj})}\\ \hline
a model $(\cX,\cL)$ of $(X,L)$
 &formal metric $\|\cdot\|_{\cL}$ on $L^{\An}$
 &closed model form $\theta\in\mathcal Z^{1,1}(X)$
  determined on $\cX$\\
a vertical $\mathbf Q$-Cartier divisor $D$
 &model function $\varphi_D=-\log|f_D|$, a PL weight
 &model function, with
  $\varphi_D(x_E)=b_E^{-1}\ord_E(D)$\\
change of metric
 &$\|\cdot\|_\varphi=\rme^{-\varphi}\|\cdot\|_0$
 &$\theta+\ddc\varphi$\\
positivity
 &$c_1(\overline L)\geq0$ as a CLD current
 &$\theta+\ddc\varphi$ semipositive, i.e.\ relatively nef\\
the measure
 &$c_1(\overline L_1)\wedge\cdots\wedge c_1(\overline L_n)$,
  \Cref{thm:cld-formal-intersection}
 &$\theta_1\wedge\cdots\wedge\theta_n$,
  \eqref{eq:model-mixed-measure}\\
continuous common locus
 &uniform limits of semipositive model metrics on $L$
 &uniform limits of $\theta$-psh model functions\\
singular potentials
 &the sheaf $\operatorname{Psh}$ of
  \Cref{def:ch15-hierarchy}, local
 &$\PSH(X,\theta)$ of \Cref{chap:bfj-psh}, global\\
\end{tabular}
\end{center}

The last row is where the two theories part company, and it is the reason
the comparison theorem has the shape it has.  The CLD class is a
\emph{sheaf}: membership is a local condition, and there is no reference to
a global class $\theta$.  The BFJ class is defined by a condition on
\emph{all models} of a fixed projective $X$ and is not local in any
obvious sense.  A theorem comparing them must therefore either restrict to
potentials for which both descriptions are computable---model and
continuous---or prove a genuinely new statement.  Only the first is
available, and that is the content of what follows.

\subsection{The comparison theorem}
\label{sec:ch19-cld-bfj-theorem}

\begin{theorem}
\label{thm:ch19-cld-bfj-comparison}
Assume \eqref{eq:ch19-bfj-hypotheses}.  Let $L$ be an ample line bundle on
$X$, let $(\cX_0,\cL_0)$ be a model of $(X,L)$ with associated formal
metric $\|\cdot\|_0$ and closed model form $\theta$, and let
$V=(\{\theta\}^n)=(L^n)>0$.

\emph{\textbf{(A) Model locus.}}  For a model function $\varphi$ on
$X^{\An}$, the following are equivalent:
\begin{enumerate}
\item\label{itm:ch19-cb:bfj} $\varphi$ is $\theta$-psh in the BFJ sense,
 that is, some determination of $\theta+\ddc\varphi$ on a model is
 relatively nef;
\item\label{itm:ch19-cb:zhang} the model metric
 $\|\cdot\|_\varphi=\rme^{-\varphi}\|\cdot\|_0$ on $L$ is semipositive
 in the sense of \Cref{sec:metrics-semipositive};
\item\label{itm:ch19-cb:cld} the CLD Chern current
 $c_1(L,\|\cdot\|_\varphi)$ is a positive current;
\item\label{itm:ch19-cb:delta} on $X_F$ the Chern delta-form
 $c_1^\delta(L_F,\|\cdot\|_{\varphi,F})$ is positive;
\item\label{itm:ch19-cb:funct} on $X_F$ the local metric weights
 are functorially psh in the sense of \cite{GK19};
\item\label{itm:ch19-cb:curve} for every closed integral curve
 $C\subseteq X_F$, the restricted metric on $C^{\An}$ is psh in the
 current-positive sense of \cite{GK19}.
\end{enumerate}
When these hold, and more generally for any $n$-tuple of model functions,
the BFJ mixed model measure and the CLD curvature product agree.  For
semipositive factors this is the positive Bedford--Taylor product; arbitrary
model factors use its dpsh multilinear extension, yielding a signed measure:
\begin{equation}
\label{eq:ch19-cld-bfj-model}
 \begin{aligned}
  \MA_\theta(\varphi_1,\ldots,\varphi_n)
   &=c_1\bigl(L,\|\cdot\|_{\varphi_1}\bigr)\wedge\cdots\wedge
     c_1\bigl(L,\|\cdot\|_{\varphi_n}\bigr)\\
   &=\sum_Eb_E
     \bigl(\cL_{\varphi_1}|_E\cdots\cL_{\varphi_n}|_E\bigr)
     \delta_{x_E},
 \end{aligned}
\end{equation}
and this common measure is the Chambert--Loir measure of the same models.
Its total mass is $V$.

\emph{\textbf{(B) Continuous locus.}}  Let $\varphi\in C^0(X^{\An})$ be a
uniform limit of $\theta$-psh model functions.  Then
$\|\cdot\|_\varphi$ is a continuous semipositive metric on $L$ in the
sense of Zhang, $c_1(L,\|\cdot\|_\varphi)$ is a positive CLD current, and
\begin{equation}
\label{eq:ch19-cld-bfj-continuous}
 \MA_\theta(\varphi)=c_1\bigl(L,\|\cdot\|_\varphi\bigr)^n
\end{equation}
as measures of mass $V$ on $X^{\An}$.  The same holds for mixed products
of $n$ such potentials.

\emph{\textbf{(C) Scope.}}  Neither \eqref{eq:ch19-cld-bfj-model} nor
\eqref{eq:ch19-cld-bfj-continuous} is asserted outside the model and
continuous loci.  Specifically, this book proves no statement of the
following four forms:
\begin{enumerate}[resume]
\item\label{itm:ch19-cb:scope1} that a bounded $\theta$-psh function is
 locally CLD-psh;
\item\label{itm:ch19-cb:scope2} that a locally CLD-psh function on
 $X^{\An}$ is $\theta$-psh for some closed model form $\theta$;
\item\label{itm:ch19-cb:scope3} that the BFJ monotone Monge--Amp\`ere
 measure of an unbounded $\theta$-psh function equals a CLD
 Bedford--Taylor product;
\item\label{itm:ch19-cb:scope4} that the two constructions assign the same
 mass to a pluripolar set.
\end{enumerate}
\end{theorem}

\begin{proof}
\emph{Step 1: the equivalences in (A).}  For integral parlinear
metric weights over $K$, positivity of the original CLD Chern current is
equivalent to residual nefness by \Cref{thm:ch17-psh-residual-nef}.
Taking a tensor power handles rational model metrics.  Hence
\cref{itm:ch19-cb:zhang}$\Leftrightarrow$\cref{itm:ch19-cb:cld} holds
in the original field setting.  Residue-nef semipositivity is preserved and
reflected after complete scalar extension
\cite[Proposition~7.6\,(2')]{GJR26}.  On $X_F$, the equivalences between
formal semipositivity and \cref{itm:ch19-cb:delta}--\cref{itm:ch19-cb:curve}
are exactly \Cref{thm:gk-formal-positivity-comparison}.  This does not
descend an ordinary delta-form to the discretely valued field, nor does
it infer scalar-extension invariance of arbitrary positive currents.

It remains to identify \cref{itm:ch19-cb:bfj} with
\cref{itm:ch19-cb:zhang}.  Choose a model $\cX$ dominating $\cX_0$ on which
$\varphi=\varphi_D$ for a vertical $\mathbf Q$-Cartier divisor $D$, and let
$\cL_\varphi=\rho^*\cL_0\otimes\cO_{\cX}(D)$.  By
\Cref{prop:metrics-divisor-model-function} and the sign convention
\eqref{eq:appJ-model-divisor}, the model metric of $\cL_\varphi$ is exactly
$\|\cdot\|_\varphi$, and by \eqref{eq:ch20-sign} the closed model form
determined by $\cL_\varphi$ is $\theta+\ddc\varphi$.  Semipositivity in
\Cref{sec:metrics-semipositive} is by definition relative nefness of
$\cL_\varphi$, which is the BFJ definition of $\theta$-psh; and both are
independent of the chosen dominating model by
\Cref{lem:metrics-nef-independence} and \Cref{prop:nef-pullback}.  Hence
\cref{itm:ch19-cb:bfj}$\Leftrightarrow$\cref{itm:ch19-cb:zhang}.

\emph{Step 2: the measures in (A).}  The right-hand equality of
\eqref{eq:ch19-cld-bfj-model} is the definition
\eqref{eq:model-mixed-measure} of the BFJ mixed model measure, once one
observes that the divisorial-point normalization used there,
$\varphi_D(x_E)=b_E^{-1}\ord_E(D)$, is the same as the CLD one; this is
\Cref{rem:divisorial-normalization}.  The middle equality is
\Cref{thm:cld-recovers-cl}, whose hypotheses---$K$ discretely valued, $X$
proper and reduced, $\cX$ relatively normal---are implied by
\eqref{eq:ch19-bfj-hypotheses}.  That the same measure is the
Chambert--Loir measure is the last sentence of that theorem, and that it is
also, after extension to $F$ and pushforward of the resulting measure,
the delta-form product is \Cref{thm:gk-delta-cl-measure}.  The total
mass is $(L^n)=V$ by \Cref{prop:ch19-model-basic}\,\cref{itm:ch19-model:mass}.

\emph{Step 3: reduction of (B) to a decreasing sequence.}  Choose
$\theta$-psh model functions $\varphi_m$ with
$\|\varphi_m-\varphi\|_\infty\leq2^{-m-3}$, and put
\[
 c_m=2^{-m},\qquad \psi_m=\varphi_m+c_m.
\]
Because $v(\pi)=1$ and $c_m$ is rational, $c_m$ is a rational model
constant, obtained from a rational multiple of the full special fibre.
Thus $\psi_m$ is again a $\theta$-psh model function.  Moreover,
\[
 \psi_m-\psi_{m+1}
 \geq(c_m-c_{m+1})-(2^{-m-3}+2^{-m-4})>0.
\]
Also $\psi_m\geq\varphi$ and
$\|\psi_m-\varphi\|_\infty\leq2^{-m}+2^{-m-3}\to0$.
Hence $\psi_m\searrow\varphi$ uniformly.  Arbitrary real shifts preserve
pshness, but need not be rational model constants over a DVR; that is why
the shifts here are chosen rationally.

\emph{Step 4: passing to the limit on both sides.}  On the BFJ side,
$\MA_\theta$ is continuous along decreasing sequences of bounded
$\theta$-psh functions by \Cref{prop:e1-monotone-continuity}, so
$\MA_\theta(\psi_m)\to\MA_\theta(\varphi)$ weakly.  On the CLD side, the
metrics $\|\cdot\|_{\psi_m}$ converge uniformly to $\|\cdot\|_\varphi$, so
$c_1(L,\|\cdot\|_{\psi_m})^n\to c_1(L,\|\cdot\|_\varphi)^n$ weakly by
\Cref{thm:cld-bt-continuity}; the hypotheses of that theorem are met because
each $\psi_m$ is continuous psh, and the products are taken with all $n$
factors equal.  By Step~2 the two sequences of measures are equal term by
term, hence so are their limits.  This is
\eqref{eq:ch19-cld-bfj-continuous}.  That $\|\cdot\|_\varphi$ is a
continuous semipositive metric in Zhang's sense is the definition, given
Step~1; that $c_1(L,\|\cdot\|_\varphi)$ is positive follows from
\Cref{prop:bt-current-stability}, positivity being closed under weak limits.
The mixed case is identical, using the multilinear version of both
continuity statements.  This common-locus comparison will be used again
in \Cref{thm:bfj-cld-ma-common-locus}; that later formulation is not an
input to the present proof.

\emph{Step 5: (C).}  Nothing is to be proved; the assertion is that the
four displayed statements are not proved.  We record why each is out of
reach of the argument above.  Items \cref{itm:ch19-cb:scope1} and
\cref{itm:ch19-cb:scope2} concern classes that are not both local: Steps~1--2
identify \emph{model} objects, and the identification propagates in Step~4
only along uniformly convergent sequences, whereas a bounded $\theta$-psh
function is in general obtained from a \emph{decreasing pointwise net}
of model functions and need not be continuous.  Item
\cref{itm:ch19-cb:scope3} fails to follow because the ordinary CLD
recursion used here has continuous admissible factors; its products are not
supplied for arbitrary unbounded limits.  The decreasing-net consequence
\Cref{cor:ch16-decreasing-nets} still requires a continuous finite limit.
The BFJ monotone extension of
\Cref{thm:bounded-ma-extension} is constructed by a different mechanism.
Item \cref{itm:ch19-cb:scope4} is not a gap in the proof but a genuine
difference: \Cref{ex:ch16-polar-failure} exhibits a canonical extension of a
CLD product carrying mass one on a polar set, whereas the non-pluripolar
products of \Cref{def:nonpluripolar-ma} discard that mass by construction.
\end{proof}

\begin{remark}[What the hypotheses are doing]
\label{rem:ch19-cld-bfj-hypotheses}
It is worth separating the roles of the assumptions
\eqref{eq:ch19-bfj-hypotheses}.
\begin{itemize}
\item \emph{Discreteness} is used in Step~2, through
\Cref{thm:cld-recovers-cl}, to make the defect factor of
\Cref{thm:cld-formal-intersection} equal to the scheme-theoretic
multiplicity $b_E$.  Over a general nontrivially valued field the CLD
formula still holds, with the defect factor written out; what changes is the
BFJ side, which has no theory of models over a non-discrete valuation ring
in the form used in \Cref{part:bfj}.
\item \emph{Residue characteristic zero} is not used anywhere in the proof
above.  It is used in \Cref{part:bfj} for continuity of envelopes, and
therefore for everything beyond bounded potentials; the removal of that
hypothesis by test ideals is discussed in
\Cref{chap:bfj-envelopes}.
\item \emph{Smoothness and projectivity of $X$} enter only through the
existence of the models used in \Cref{part:bfj}.  Statement (A) holds for
proper varieties in the formal scope of
\Cref{thm:common-model-core}, with ordinary delta tests over completed
algebraic closure.  Statement (B) assumes uniform approximation by
\emph{semipositive} model metrics.  Density of all model functions,
\Cref{prop:metrics-model-function-density}, does not preserve positivity
and cannot replace that assumption.
\item \emph{Ampleness of $L$} is used for the mass statement and for the
BFJ theory's standing hypothesis that $\{\theta\}$ be ample.  For a
general $L$, equality \eqref{eq:ch19-cld-bfj-model} of signed measures
still holds, since both sides are given by the same intersection numbers.
\end{itemize}
\end{remark}

\begin{corollary}
\label{cor:ch19-transfer}
Assume \eqref{eq:ch19-bfj-hypotheses}.
\begin{enumerate}
\item\label{itm:ch19-tr:curve} If $\dim X=1$, then for every continuous
 semipositive metric the BFJ measure, the Chambert--Loir measure and the CLD
 curvature all equal $\mu_0+\Delta_\Gamma F$ of
 \Cref{prop:ch19-curve-model-curvature} and its uniform limits.
\item\label{itm:ch19-tr:toric} If $X$ is toric and the metric is
 invariant, all three equal $n!\MA_{\mathbf R}(\phi)$ after tropical
 pushforward, by \Cref{thm:ch19-toric-ma}.
\item\label{itm:ch19-tr:abelian} If $X=A$ is abelian with split Raynaud
 data and $L$ is ample with its cubical metric, all three equal the Haar
 measure of \Cref{thm:ch19-canonical-ma}, of mass $(L^g)$.
\end{enumerate}
\end{corollary}

\begin{proof}
In each case the metric is a uniform limit of semipositive model metrics:
for \cref{itm:ch19-tr:curve} by the definition of a continuous
semipositive metric, with the approximating curvatures computed by
\Cref{prop:ch19-curve-model-curvature}; for \cref{itm:ch19-tr:toric} by
Step~3 of the proof of \Cref{thm:ch19-toric-ma}; for
\cref{itm:ch19-tr:abelian} by the explicit approximation
\eqref{eq:ch19-tate-model-approx} in dimension one, and in general by the
semipositive rational polyhedral metrics attached to Mumford models.  On
increasingly fine rational periodic polytopal decompositions, convex
quasi-periodic polyhedral weights converge uniformly to the quadratic weight
$f$ of \eqref{eq:ch19-canonical-weight}.  Convexity is part of the
Mumford-model construction, rather than a consequence of arbitrary affine
interpolation; see \cite{Gub10CM}.  Now apply
\Cref{thm:ch19-cld-bfj-comparison}\,(B).
\end{proof}

\begin{warning}{The comparison is about measures, not about classes}
\Cref{thm:ch19-cld-bfj-comparison} identifies two \emph{operators} on a
common domain.  It does not identify the two \emph{classes} of psh objects,
and the difference is not a technicality.  On the CLD side the class is a
sheaf with no global positivity reference; on the BFJ side it is a global
class attached to a fixed $\theta$.  Already for $X=\mathbf P^1$ the
constant function $0$ is CLD-psh on every open set but is
$\theta$-psh only when $\theta$ is semipositive, and a CLD-psh function
on a small open set need not extend to $X^{\An}$ at all.  The local versus
global metric question is \Cref{prob:frontier-cld-global-metric}; the
present theorem does not touch it.
\end{warning}

\section{What is not known for singular psh functions}
\label{sec:ch19-limitations}

The ordinary curvature-product computations of this chapter use model
metrics, continuous metrics, or the specifically admissible potentials.
The Thuillier statements and invariant toric existence theorem also treat
singular functions, through their own constructions.  For each of the
assertions below, either the statement is false for arbitrary singular psh
functions, or it is not known, or it has no meaning at all until an extra
construction is supplied.  The purpose of this section is to say which of
the three, and where the book resolves the matter.

\Cref{part:bfj,part:comparisons} use this table as a scope-discipline
device: a hypothesis in a later chapter that looks unnecessarily strong
should be checked against it before being weakened.

\subsection{The table}
\label{sec:ch19-limitations-table}

In the third column, \emph{false} means a counterexample is available in
this book; \emph{undefined} means the expression on which the assertion
depends is not defined; and \emph{open} means the assertion is meaningful
and neither proved nor disproved here.  In the second column, ``model'' and
``continuous'' refer to the loci of
\Cref{thm:ch19-cld-bfj-comparison}\,(A) and~(B).

\begingroup
\footnotesize
\begin{longtable}{p{.195\textwidth}|p{.135\textwidth}|p{.275\textwidth}|p{.205\textwidth}}
\toprule
\textbf{Assertion}&\textbf{Model / cont.}
 &\textbf{Status and precise reason for arbitrary singular $u$}
 &\textbf{Resolved or left open at}\\
\midrule
\endfirsthead
\toprule
\textbf{Assertion}&\textbf{Model / cont.}
 &\textbf{Status and precise reason for arbitrary singular $u$}
 &\textbf{Resolved or left open at}\\
\midrule
\endhead

(L1) The Bedford--Taylor product $(\dprimeddprime u)^n$ exists as a
 current.
 &Yes,
  \Cref{thm:bt-symmetry-locality}
 &\emph{Undefined.}  The recursion \eqref{eq:ch16-recursion-schema} needs
  the product $uT$ of \eqref{eq:ch16-uT-definition}, which is constructed
  in \Cref{lem:continuous-times-positive-current} for $u$ continuous and
  $T$ positive.  For arbitrary unbounded $u$, multiplication may fail
  the required integrability and no unrestricted canonical extension of
  this recursion is supplied.
 &Global substitutes:
  \Cref{thm:bounded-ma-extension},
  \Cref{def:nonpluripolar-ma}.  A local product is
  \Cref{prob:frontier-local-ma}.\\

(L2) The total mass is $(L^n)$.
 &Yes, \eqref{eq:ch19-cld-bfj-model},
  \Cref{thm:ch17-degree}
 &\emph{False in general.}  Even where a product is defined by monotone or
  non-pluripolar extension, mass can drop: the non-pluripolar product of
  \Cref{def:nonpluripolar-ma} has mass $\leq(L^n)$, with equality
  characterizing the full-mass class.
 &\Cref{chap:bfj-energy};
  \Cref{def:e1}.\\

(L3) The measure puts no mass on the polar locus $\{u=-\infty\}$.
 &Yes; vacuous for continuous $u$
 &\emph{Depends on the construction, and the two answers differ.}
  \Cref{thm:cld-polar-negligibility} proves that a product of continuous
  psh factors does not charge the polar locus of an \emph{auxiliary}
  locally approachable, smoothly integrable function, under the stated
  separation and compact-sublevel barrier hypotheses.  It does not define
  that auxiliary function's own product.  A polar extension can carry mass
  there, whereas non-pluripolar products discard it by definition.
 &\Cref{ex:ch16-polar-failure};
  \Cref{prop:bounded-ma-no-polar-mass},
  \Cref{def:nonpluripolar-ma}.\\

(L4) Positivity of $\dprimeddprime[u]$ implies that $u$ belongs to the
 approximation-generated CLD hierarchy.
 &Yes for continuous functions on smooth proper algebraic curves over an
  algebraically closed complete nontrivially valued field,
  \Cref{thm:p063-current-psh-ordinary-regularization}
 &\emph{Open beyond the continuous smooth-proper algebraic-curve and
  $\mathbf R$-PL cases.}
  Positivity of the current is a condition from above, whereas the CLD
  class is generated from below by decreasing limits of smooth psh
  functions.  The tree regularization above supplies the missing implication
  only for continuous functions in its stated curve scope; arbitrary
  singular or $-\infty$-valued functions remain unresolved even in
  dimension one.  For $\mathbf R$-PL functions on a boundaryless space
  the implication is available in every dimension.
 &\Cref{thm:verified-pl-psh-equivalence,%
thm:p01-real-pl-positive-current-candidate}.\\

(L5) $u$ is a local decreasing limit of \emph{smooth} psh functions.
 &Yes for continuous current-psh functions on smooth proper algebraic curves
  over an algebraically closed complete nontrivially valued field
 &\emph{Open in general.}  Eventual stabilization is a theorem
  (\Cref{prop:ch15-hierarchy-stabilizes}).  A separate finite-dimensional
  result proves full first-stage collapse for the generated sheaf on finite
  ambient-induced real-affine stars, including extended-real functions with
  polar branches (\Cref{thm:ch16-collapse-finite-tropical-curves}).  For
  every nonempty branch set $B$, the exact CLD star $X_B$ of
  \cite[\S6.7.3]{CLD25} also has full first-stage collapse and
  fixed-neighbourhood decreasing smooth-psh approximation
  (\Cref{thm:p071-infinite-star-collapse-candidate}).  These theorems do not
  assert that every analytic, sedentary, or other nonambient tropical germ
  has one of the stated smooth structures.
  \Cref{thm:p063-current-psh-ordinary-regularization} settles the continuous
  current-positive subcase in its smooth-proper algebraic-curve scope;
  arbitrary singular analytic-curve functions outside that scope,
  arbitrary higher-dimensional spaces, and other nonambient local models
  remain open.
 &\Cref{prob:frontier-cld-hierarchy,%
  thm:ch16-collapse-finite-tropical-curves,%
  thm:p071-infinite-star-collapse-candidate};
  \cite[Corollary~5.4]{Wanner19}.\\

(L6) An interior local maximum forces local constancy.
 &Yes on analytic spaces,
  \Cref{cor:ch15-cld-curve-max}
 &\emph{Proved for analytic CLD-psh functions in every dimension.}
  CLD-to-classical psh and the analytic curve-connectivity theorem give
  the local maximum principle, also for singular CLD-psh functions.
  This supersedes the historical uncertainty in the cited CLD source.
  It supplies no unrestricted maximum principle on abstract tropical
  spaces; the finite one-dimensional case has its separate proof.
 &\Cref{cor:ch15-cld-curve-max};
  \cite[\S0.3.5]{CLD25}.\\

(L7) The comparison principle
 $\int_{\{u<w\}}\MA(w)\leq$\newline
 $\int_{\{u<w\}}\MA(u)$.
 &Chartwise for bounded $u,w$,
  \Cref{thm:ch16-comparison-chart}
 &\emph{Undefined, then unavailable.}  The statement presupposes (L1); and
  the chartwise proof uses the subgradient inclusion
  \Cref{lem:ch16-subgradient-inclusion}, which is a statement about finite
  convex functions.  Globally the comparison principle needs finite energy,
  not local positivity.
 &\eqref{eq:model-maximum-principle};
  \eqref{eq:bj-comparison-principle}.\\

(L8) Restriction to a closed subvariety $Z$ is psh and its curvature is
 the restriction.
 &Yes,
  \Cref{cor:ch15-pullback-cases};\newline
  \cref{itm:ch15-pb-closed}
 &\emph{Meaningless when $Z\subseteq\{u=-\infty\}$.}  Then
  $u|_Z\equiv-\infty$ is not a psh function in the sense used here, the
  current $[u|_Z]$ is not defined, and no restriction formula holds.  The
  same failure affects pullback along a morphism whose image lies in the
  polar locus.
 &\Cref{prop:ch15-pl-pullback}; see also the warning after that result.\\

(L9) The metric is a delta-metric, so that $c_1^\delta$ is an ordinary
 delta-form.
 &Yes for formal and PL metrics,
  \cite[Proposition~8.11]{GK17}
 &\emph{False for ample cubical metrics in non-isotrivial mixed
  reduction.}  The canonical tropicalization is harmonic but not smooth;
  the full ambient tensor condition forces all its coordinates to be smooth.
  Conversely, isotriviality makes those coordinates smooth and the metric
  ordinary delta.  Algebraically trivial cubical metrics remain ordinary
  delta in every reduction type.
 &\Cref{rem:ch19-canonical-delta-scope,thm:p09p2-product-gauss-candidate}.\\

(L10) The CLD-psh class coincides with a global $\theta$-psh class.
 &Compared only on model and continuous loci
 &\emph{Open, and not merely technical.}  The CLD class is a sheaf with no
  reference to a global positivity datum, and a CLD-psh function on a small
  open set need not extend.  No unrestricted singular BFJ--CLD equivalence
  is proved anywhere.
 &\Cref{thm:ch19-cld-bfj-comparison}\,(C);
  \Cref{prob:frontier-cld-global-metric}.\\

(L11) The Monge--Amp\`ere equation $\MA(u)=\mu$ is solvable for the
 relevant $\mu$.
 &Yes for invariant toric and abelian data,
  \cite[Theorems~6.3.13, 8.3.3]{BGJKK25}
 &\emph{Not a local statement at all.}  The invariant toric and abelian results have their own global data and
  product constructions; the toric theorem can even treat measures of
  infinite energy.  The original BFJ variational theory has different
  hypotheses, including an ample class, discreteness, and residue
  characteristic zero.  None supplies an unrestricted local CLD existence
  theorem for arbitrary singular potentials.
 &\Cref{chap:bfj-ma};
  \Cref{prob:frontier-general-ma}.\\

(L12) The curvature current determines the potential up to an additive
 constant.
 &Up to harmonic ambiguity, \Cref{sec:ch16-ambiguity-levels}
 &\emph{False without global hypotheses.}  Locally, any harmonic function
  may be added.  On a compact connected boundaryless analytic space, if both signs of
  the difference are CLD-psh, the analytic maximum principle makes it
  constant.  A zero curvature current alone does not establish that
  two-sided CLD membership for an arbitrary continuous function.
 &\Cref{thm:ch16-pluriharmonic-criterion};
  \Cref{prob:frontier-pluriharmonic}.\\

(L13) Smooth decreasing regularization can be arranged so that the
 curvature measures converge weakly.
 &Yes under uniform convergence,
  \Cref{thm:cld-bt-continuity}
 &\emph{Undefined, then unproved.}  Weak continuity of the CLD product is
  proved for continuous psh functions converging locally uniformly, and
  for decreasing nets of continuous psh functions with a finite continuous
  limit.  For arbitrary decreasing nets with unbounded limits the present
  recursion does not supply the limiting products, so this theorem gives
  no unrestricted convergence assertion.
 &\Cref{cor:ch16-decreasing-nets},
  \Cref{rem:ch16-not-claimed}.\\

(L14) The Chern--Levine--Nirenberg estimate bounds the local mass.
 &Yes,
  \Cref{thm:cld-cln}
 &\emph{Fails as stated.}  The constant in \Cref{thm:cld-cln} is a product
  of sup-norms, which is infinite for unbounded potentials.  The weighted
  form \Cref{thm:ch16-weighted-cln} replaces sup-norms by weights and is what
  survives; it is the tool used to prove (L3) in its stated scope.
 &\eqref{eq:ch16-weighted-cln-main};
  \eqref{eq:cld-weighted-cln}.\\
\bottomrule
\end{longtable}
\endgroup

\subsection{How to read the table}
\label{sec:ch19-limitations-reading}

Three structural remarks make the fourteen rows less arbitrary than they
look.

\emph{First, most failures come from one place.}  Rows (L1), (L7), (L13) and
(L14) are all consequences of the single fact that the recursion defining
Bedford--Taylor products multiplies a function by a current, and that
operation needs an admissible integrable multiplier, supplied here by
continuous bounded factors.  Everything downstream of the recursion
inherits the hypothesis.  The global theories evade it not by improving the
local product but by replacing it: BFJ by monotone extension from bounded
potentials (\Cref{thm:bounded-ma-extension}), and the non-pluripolar
formalism by truncation (\Cref{def:nonpluripolar-ma}).  Consequently rows
(L2) and (L3) record genuine \emph{differences} between the resulting
operators, not gaps.

\emph{Second, two failures come from the definition of the class.}  Rows
(L4) and (L5) are the two halves of the same question: the CLD class is
generated from below, so membership is proved by exhibiting approximants,
and any statement asserting membership from a property of the current is a
converse that must be proved.  Wanner's curve theorem compares current
positivity with Thuillier's comparison class; the separate
\Cref{thm:p063-current-psh-ordinary-regularization} supplies the required
ordinary CLD-smooth regularization for continuous functions on smooth proper
curves.  In higher dimension the only converse recorded here is for
$\mathbf R$-PL functions.

\emph{Third, one failure is geometric.}  Row (L9) is not about singular
functions at all: it concerns a perfectly continuous, even canonical, metric
on an abelian variety, and the obstruction is that a harmonic
tropicalization need not be smooth.  It is included because it is the row
most likely to be misremembered as settled, and because it delimits exactly
what \Cref{sec:ch19-canonical} proves.

\noindent\begin{minipage}{\textwidth}
\begin{warning}{Three phrases that are not interchangeable}
The table distinguishes them and so does the rest of the book.
\begin{itemize}
\item \emph{``Not known''} means the assertion is meaningful, and neither a
 proof nor a counterexample is given here.  The unresolved ranges of (L4), (L5), and (L10) are examples;
 the analytic maximum principle in (L6) and the ample criterion in (L9)
 are proved.
\item \emph{``Not defined''} means the expression itself has no referent
 until a further construction is supplied.  Rows (L1), (L7), (L8), (L13).
\item \emph{``Different''} means two specified constructions exist and give different answers, as in
 (L2) and (L3).  Local harmonic ambiguity in (L12) and the missing finite
 sup-norm bound in (L14) are separate limitations.
\end{itemize}
Conflating the second with the first produces false open problems;
conflating the third with either produces false theorems.
\end{warning}
\end{minipage}\par

\section{Comparison and limitations ledger}
\label{sec:ch19-ledger}

\Cref{sec:ch19-limitations} lists what fails.  This section lists, in the
complementary direction, what has been proved and against which precise
hypotheses, one entry per class of object met in this chapter.  Each entry
has two halves: what is established, with its source, and the boundary of
that establishment.  The list is meant to be quotable, and later chapters
quote it.

\begin{description}[style=nextline,leftmargin=0pt,labelindent=0pt,
  labelsep=0pt,itemsep=.8\baselineskip,font=\normalfont\bfseries]
\item[Model metric on a curve]
\emph{Proved.} The Chern current is $\mu_0+\Delta_\Gamma F$, an atomic
signed measure of mass $\deg_CL$, positive exactly when a finite list of
outgoing-slope inequalities holds
(\Cref{prop:ch19-curve-model-curvature}); the slope formula is the
retraction of Poincar\'e--Lelong
(\Cref{thm:ch19-curve-slope-formula}).
\par\emph{Boundary.} Both statements are for potentials that factor through
a skeleton; a continuous semipositive metric on a curve has a nonatomic
curvature in general (\Cref{ex:ch19-p1-nonpl}).

\item[Formal metric on a proper algebraic variety]
\emph{Proved.} Nef model, positive Chern delta-form, functorial psh, and curve
positivity are equivalent \cite[Theorem~6.10]{GK19}.
\par\emph{Boundary.} There is no representation theorem for arbitrary
positive currents by formal metrics.

\item[Continuous function on a smooth proper curve]
\emph{Proved.} Positivity of $\dprimeddprime[u]$ is equivalent to
Thuillier subharmonicity \cite[Def.~3.1.5]{Thu05} over an algebraically
closed complete nontrivially valued field
\cite[Theorem~4.7 and Corollary~4.8]{Wanner19}; continuous CLD-psh functions
belong to both classes.  Conversely, every continuous member of either class
lies locally in the first ordinary CLD-smooth decreasing-limit stage by
\Cref{thm:p063-current-psh-ordinary-regularization}.
\par\emph{Boundary.} The ordinary-smooth theorem does not treat arbitrary
upper semicontinuous functions that may equal $-\infty$, singular or
nonproper analytic curves outside its stated algebraic scope, or higher
dimension.  Decreasing regularization by Thuillier-smooth, hence skeletal
piecewise affine, subharmonic functions is unconditional and holds for
arbitrary upper semicontinuous subharmonic functions
(\Cref{thm:ch19-thuillier-approximation}).

\item[Subharmonic function on a smooth strictly analytic curve]
\emph{Proved.} Subharmonic functions are exactly the degree-zero currents
with positive $\ddc$, positive degree-one currents are exactly the positive
Radon measures, and a subharmonic function is determined by its values at the
points with residue field finite over the base
(\Cref{thm:ch19-thuillier-sh-currents}, from
\cite[Thm.~3.4.12, Cor.~3.4.13, Thm.~3.6.26]{Thu05}); no algebraic closedness
and no properness is required.
\par\emph{Boundary.} The statement is one-dimensional.  No corresponding
description of the singular psh class by a positivity property of a current
is available in higher dimension; that is the first row of
\Cref{sec:ch19-limitations}.

\item[$\mathbf R$-PL function on a good strict analytic space]
\emph{Proved.} Residue-nef semipositivity equals classically psh
\cite[Theorem~4.3]{GR25Classical}; on the relative interior it also equals
CLD-v2 psh by \Cref{thm:verified-pl-psh-equivalence}.
\par\emph{Boundary.} The CLD-v2 reverse implication at a relative-boundary
point is not proved, and no Monge--Amp\`ere operator is asserted for every
classically psh function.

\item[Toric continuous invariant metric]
\emph{Proved.} Tropical psh, local psh approximation, and, on a complete fan,
Zhang semipositivity agree under \cite[Theorem~5.3.11]{BGJKK25}.
\par\emph{Boundary.} No corresponding equivalence is asserted for arbitrary
noninvariant metrics.

\item[Toric Bedford--Taylor product]
\emph{Proved.} Pushforward equals the tropical product
\cite[Theorem~6.1.1]{BGJKK25}.
\par\emph{Boundary.} The formula requires the locally psh-approximable class
stated there.

\item[Canonical abelian tropicalization]
\emph{Proved.} The canonical tropicalization is a harmonic tropicalization
\cite[Proposition~16.2]{GJR26}.
\par\emph{Boundary.} It need not be a smooth tropicalization.

\item[Canonical abelian metrics]
\emph{Proved.} Invariant semipositive metrics and skeletal Monge--Amp\`ere
equations are handled in \cite[\S\S7--8]{BGJKK25}.
\par\emph{Boundary.} Ordinary delta-metricity was left open in
\cite[Remark~16.4]{GJR26}.  For ample cubical metrics it is characterized by
residual isotriviality in
\Cref{thm:p09p2-product-gauss-candidate}; no classification of every
individual nonample cubical line bundle is asserted here.

\item[Continuous Zhang-semipositive limit]
\emph{Proved.} The first Chern current is positive under the local formal
approximation hypotheses of \cite[Theorem~1.3]{GJR26}.
\par\emph{Boundary.} A positive first Chern current alone does not
automatically define all higher Bedford--Taylor products.

\item[BFJ singular psh metric]
\emph{Proved.} Global Monge--Amp\`ere and energy are available under the BFJ
hypotheses developed in \cref{part:bfj}.
\par\emph{Boundary.} No equality is asserted between the entire BFJ singular
class and the entire local CLD or classical-psh class.

\item[CLD versus BFJ on the model and continuous loci]
\emph{Proved.} Under \eqref{eq:ch19-bfj-hypotheses}, the BFJ notion of
$\theta$-psh model function, Zhang semipositivity of the associated model
metric, positivity of the CLD Chern current over $K$, and the ordinary Chern
delta-form, functorial psh, and closed-curve tests over $F$ all agree;
and the BFJ model
measure, the Chambert--Loir measure, the delta-form product and the CLD
Bedford--Taylor product all agree, in the model case and after uniform
approximation in the continuous case
(\Cref{thm:ch19-cld-bfj-comparison}).
\par\emph{Boundary.} Four statements are explicitly not asserted:
that bounded $\theta$-psh functions are locally CLD-psh; that locally
CLD-psh functions are $\theta$-psh for some $\theta$; that the two
operators agree on unbounded potentials; and that they assign the same mass
to a pluripolar set.  See \Cref{thm:ch19-cld-bfj-comparison}\,(C) and
\Cref{sec:ch19-limitations}.
\end{description}

\section{Chapter dictionary}
\label{sec:ch19-dictionary}

\subsection{From this chapter to the rest of the book}
\label{sec:ch19-dict-forward}

\begin{center}
\small
\begin{tabular}{p{.28\textwidth}|p{.26\textwidth}|p{.28\textwidth}}
\textbf{Computed here}&\textbf{Used as}&\textbf{Where}\\ \hline
graph Laplacian of a model function
 &curvature of a model metric on a curve
 &\Cref{prop:skeletal-curvature-curve},
  \eqref{eq:ch19-curve-formal-measure}\\
slope formula from Poincar\'e--Lelong
 &consistency of \Cref{part:cld} with
  \Cref{chap:model-calculations}
 &\Cref{thm:ch19-curve-slope-formula},
  \Cref{ex:ch19-tate-slope}\\
Wanner's comparison
 &identification of current positivity with Thuillier subharmonicity on
  curves
 &\Cref{thm:wanner-curve-comparison},
  \Cref{cor:ch19-curve-three-classes}\\
ordinary-smooth curve regularization
 &identification of the continuous current-positive class with the local
  first decreasing-limit stage
 &\Cref{thm:p063-current-psh-ordinary-regularization}\\
convexity is toric positivity
 &normalization check for every later positivity notion
 &\Cref{thm:toric-positivity},
  \Cref{rem:ch19-toric-no-content}\\
tropical Monge--Amp\`ere
 &$n!\MA_{\mathbf R}$ with mass $(L^n)$
 &\Cref{thm:ch19-toric-ma},
  \Cref{thm:toric-bfj-alexandrov}\\
atomic, segment, and smooth toric measures
 &the three shapes a Monge--Amp\`ere measure can take
 &\Cref{ex:ch19-toric-two-atoms,ex:ch19-toric-segment,%
        ex:ch19-toric-p2-smooth}\\
model mixed measure in dimension $n$
 &the BFJ model measure and the Chambert--Loir measure
 &\eqref{eq:ch19-model-mixed-measure},
  \eqref{eq:model-mixed-measure}\\
harmonic canonical tropicalization
 &chart input for the special formal-current products
 &\Cref{thm:abelian-harmonic-tropicalization},
  \Cref{thm:bt-harmonic-compatibility}\\
smoothness of $\tau$
 &characterizes ordinary delta-metricity of ample cubical metrics
 &\Cref{prop:ch19-tau-smooth},
  \Cref{thm:p09p2-product-gauss-candidate}\\
cubical metric and its curvature
 &translation-invariant currents and Haar measures
 &\Cref{thm:ch19-canonical-structure,thm:ch19-canonical-ma}\\
Tate curve canonical metric
 &the smallest complete example of the abelian theory
 &\Cref{ex:ch19-tate-canonical}\\
CLD--BFJ comparison
 &the transfer used throughout \Cref{part:comparisons}
 &\Cref{thm:ch19-cld-bfj-comparison},
  \Cref{thm:bfj-cld-ma-common-locus}\\
the limitations table
 &scope discipline for hypotheses in later parts
 &\Cref{sec:ch19-limitations},
  \Cref{sec:ch16-delimitation-table}
\end{tabular}
\end{center}

\subsection{From the general theory back to this chapter}
\label{sec:ch19-dict-backward}

\begin{center}
\small
\begin{tabular}{p{.28\textwidth}|p{.28\textwidth}|p{.26\textwidth}}
\textbf{Later object}&\textbf{Smallest instance here}
 &\textbf{Computed in}\\ \hline
$\MA_\theta(\varphi)$ for a model function
 &$\mu_0+\Delta_\Gamma F$ on a metric graph
 &\Cref{prop:ch19-curve-model-curvature},
  \Cref{ex:two-vertex-ma}\\
uniform approximation by model metrics
 &interpolation of a quadratic on a circle
 &\eqref{eq:ch19-tate-model-approx}\\
the factor $n!$
 &$2!\cdot\tfrac12=1=(\cO_{\mathbf P^2}(1)^2)$
 &\Cref{ex:ch19-toric-p2-smooth}\\
a Monge--Amp\`ere measure of intermediate dimension
 &uniform measure on a segment in $\mathbf R^2$
 &\Cref{ex:ch19-toric-segment}\\
translation-invariant curvature
 &$\tfrac d\ell\,\rd t$ on a circle
 &\eqref{eq:ch19-tate-measure}\\
a polarization of type $(1,3)$
 &$B=\begin{psmallmatrix}2&1\\1&2\end{psmallmatrix}$,
  $\det B=3$
 &\Cref{ex:ch19-abelian-surface-degenerate}\\
mixed reduction
 &$b=r=1$, skeleton a circle, mass $2mc\ell$
 &\Cref{ex:ch19-abelian-surface-mixed}\\
a harmonic but nonsmooth coordinate
 &$-\log|\chi|$ on a nonisotrivial extension
 &\Cref{prop:ch19-tau-smooth}\,\cref{itm:ch19-smooth:rank1}\\
the model versus continuous distinction
 &an atom versus a smooth density on $\mathbf P^1$
 &\Cref{ex:ch19-p1-standard,ex:ch19-p1-nonpl}
\end{tabular}
\end{center}

\subsection{The mass ledger}
\label{sec:ch19-mass-ledger}

Every Monge--Amp\`ere computation of this chapter has had its total mass
checked against a degree, in the manner of
\Cref{sec:ch08-mass-ledger}.  The checks are collected here.

\begin{center}
\small
\begin{tabular}{p{.36\textwidth}|p{.24\textwidth}|p{.24\textwidth}}
\textbf{Computation}&\textbf{Mass}&\textbf{Degree}\\ \hline
model metric on a curve,
 \eqref{eq:ch19-curve-formal-measure}
 &$\sum_Yb_Y\deg(\cL|_Y)$&$\deg_CL$\\
slope formula on the Tate curve,
 \Cref{ex:ch19-tate-slope}
 &$2-1-1=0$&$\deg\divis(f)=0$\\
model intersection check,
 \Cref{ex:ch19-tate-slope-model}
 &$(H\cdot Y_i)+(V\cdot Y_i)=0$&$0$\\
$\cO(1)$ on $\mathbf P^1$, standard,
 \Cref{ex:ch19-p1-standard}
 &$1$&$1$\\
$\cO(1)$ on $\mathbf P^1$, smoothed,
 \Cref{ex:ch19-p1-nonpl}
 &$0-(-1)=1$&$1$\\
$\cO(1,1)$ on $\mathbf P^1\times\mathbf P^1$, two atoms,
 \Cref{ex:ch19-toric-two-atoms}
 &$2(\tfrac12+\tfrac12)=2$&$(H_1+H_2)^2=2$\\
$\cO(1,1)$, segment measure,
 \Cref{ex:ch19-toric-segment}
 &$2\cdot1\cdot1=2$&$2$\\
$\cO(1)$ on $\mathbf P^2$, log-sum-exp,
 \Cref{ex:ch19-toric-p2-smooth}
 &$2\cdot\tfrac12=1$&$(\cO(1)^2)=1$\\
$\mathbf F_1$, $2H-E$,
 \Cref{ex:ch19-toric-hirzebruch}
 &$2+1=3$&$(2H-E)^2=3$\\
good reduction in dimension $n$,
 \Cref{ex:ch19-good-reduction}
 &$(L^n)$&$(L^n)$\\
nonreduced fibre,
 \Cref{ex:ch19-nonreduced}
 &$2e=d$&$\deg L=d$\\
Tate curve, canonical metric,
 \eqref{eq:ch19-tate-measure}
 &$\tfrac d\ell\cdot\ell=d$&$\deg L=d$\\
Tate curve, $N$-fold interpolation,
 \eqref{eq:ch19-tate-model-approx}
 &$N\cdot\tfrac dN=d$&$d$\\
totally degenerate abelian surface,
 \Cref{ex:ch19-abelian-surface-degenerate}
 &$2!\det B\cdot1=6$&$(L^2)=2!\cdot3=6$\\
mixed abelian surface,
 \Cref{ex:ch19-abelian-surface-mixed}
 &$2mc\ell=6$&$(L^2)=6$
\end{tabular}
\end{center}

Two entries deserve a comment.  The two abelian surfaces have the same
degree $6$ and the same polarization type $(1,3)$ but skeleta of
dimensions $2$ and $1$; the mass therefore does \emph{not} detect the
reduction type, whereas the dimension of the support does.  And the two
Tate-curve entries are the same measure computed twice, once as a Haar
measure and once as a limit of atomic model measures; their agreement is the
one-dimensional instance of
\Cref{thm:ch19-cld-bfj-comparison}\,(B).

\subsection{Four checks to run on any new computation}
\label{sec:ch19-checks}

\begin{enumerate}
\item \emph{Sign.}  Is the potential $\log|f|$ or $v(f)=-\log|f|$?  The
 two have opposite Laplacians, \eqref{eq:appJ-cld-pl} against
 \eqref{eq:appJ-valuation-pl}.  On a curve, a convex kink must carry
 positive mass.
\item \emph{Factorial.}  In dimension $n$, is the answer the Alexandrov
 measure or $n!$ times it?  The geometric mass is $n!\vol_M(P)=(L^n)$,
 \eqref{eq:appJ-toric-degree}.
\item \emph{Lattice.}  Are edge lengths paired with their
 corresponding weights, and are volumes lattice-normalized?
 See \Cref{rem:lattice-length} and
 \Cref{def:ch11-lattice-measure}.
\item \emph{Mass against degree.}  Does the total mass equal the
 intersection number it must?  Every row of
 \Cref{sec:ch19-mass-ledger} is such a check, and any discrepancy is a sign
 error in one of the first three.
\end{enumerate}

\section*{Source notes}

\paragraph{Attributions.}
The potential theory of Berkovich curves underlying
\Cref{sec:ch19-curves,sec:ch19-wanner,sec:ch19-thuillier-currents}---the
sheaves of harmonic and subharmonic germs, the operator $\ddc$, the
current spaces, the description of subharmonicity by positivity of $\ddc$,
and the approximation theorem---is Thuillier's \cite{Thu05}.
The curve comparison and the previously known restricted regularization
theorem are from \cite{Wanner19}.  The unrestricted ordinary-smooth curve
regularization is proved in
\Cref{thm:p063-current-psh-ordinary-regularization}.  Delta-form positivity,
toric convexity, and formal curve
tests are from \cite{GK19}.  The invariant toric and abelian
pluripotential theory, including the exact tropical pushforward formula, is
from \cite{BGJKK25}.  Harmonic canonical tropicalizations and the correction
that exposed the ordinary delta-metric question are from \cite{GJR26}.  The
product--Gauss criterion of
\Cref{thm:p09p2-product-gauss-candidate} is a book-level deduction from the
ordinary delta-chart formalism of \cite{GK17}, the algebraization theorem of
\cite{Gub16}, and the Raynaud metric and smoothness criteria of
\cite{CLD25}; no novelty or priority claim is made.  The common model measure
is traced independently to \cite{CL06,CLD25,GK17}.  These examples
will be reused in later parts as normalization checks for every global
Monge--Amp\`ere construction.

In detail: \Cref{thm:wanner-curve-comparison} is
\cite[Theorem~4.7 and Corollary~4.8]{Wanner19}, with the functoriality
statement \cite[Corollary~5.1]{Wanner19} and the restricted regularization
\cite[Corollary~5.4]{Wanner19}.  The comparison theorem is imported at a
source boundary because its proof uses Wanner's current-positive meaning of
``psh'', not the approximation-generated CLD hierarchy of this book.
\Cref{thm:toric-positivity} is
\cite[Proposition~8.13 and Corollary~8.14]{GK19}, and
\eqref{eq:cld-toric-ma-comparison} is \cite[Theorem~6.1.1]{BGJKK25}, whose
invariant psh comparison is \cite[Theorem~5.3.11]{BGJKK25} and whose
existence theorems are \cite[Theorems~6.3.13 and~8.3.3]{BGJKK25}.
\Cref{thm:ch19-raynaud} is the uniformization theorem of Raynaud in the
form of \cite[Appendix~B.5]{CLD25}, resting on
\cite{BoschLutkebohmert1993}, with the deformation retraction of
\cite[Chapter~6]{Berk12}.  \Cref{thm:ch19-cubical-exists},
\Cref{thm:ch19-canonical-structure},
\Cref{thm:ch19-canonical-positivity},
\Cref{thm:ch19-canonical-ma} and
\Cref{prop:ch19-canonical-subvariety} are, respectively,
\cite[Lemmas~19.1.4 and~19.1.6, Proposition~19.1.8; Proposition~19.2.5;
Proposition~19.3.1 with Corollary~19.3.2 and Theorem~B.6.4;
Corollary~19.2.8; Corollary~19.2.9 with Remark~19.3.4]{CLD25}.
\Cref{prop:ch19-tau-smooth} is \cite[Remarks~19.2.3 and
Proposition~17.4.13]{CLD25}.  \Cref{thm:abelian-harmonic-tropicalization} is
\cite[Proposition~16.2]{GJR26}, and the mixed-reduction example behind
\Cref{rem:ch19-canonical-delta-scope} is \cite[Remark~16.4 and \S17]{GJR26}.
\Cref{thm:ch19-cld-bfj-comparison} assembles
\Cref{thm:gk-formal-positivity-comparison},
\Cref{thm:cld-recovers-cl}, \Cref{thm:gk-delta-cl-measure} and the BFJ
monotone continuity of \cite[\S\S2.7 and~3]{BFJ15}; its part~(B) is the common-locus measure comparison used again in
\Cref{thm:bfj-cld-ma-common-locus}.  Its proof is given here independently,
so that the scope in part~(C) can be read with the argument in view.

\paragraph{What is new in the presentation.}
Five things are arranged differently from the sources.

First, the slope formula is derived twice.  \Cref{chap:model-calculations}
proves it by intersection theory on a semistable model;
\Cref{thm:ch19-curve-slope-formula} rederives it as the retraction of
Poincar\'e--Lelong, and \Cref{rem:ch19-slope-two-proofs} records that the
two arguments have genuinely different hypotheses.  This is done because
the identity is used in both directions later.

Second, \Cref{ex:ch19-tate-slope} and \Cref{ex:ch19-tate-slope-model} carry
one rational function on a Tate curve of circumference three through both
routes and check them atom by atom.  The numerical data are those of
\Cref{ex:ch08-tate-curve}, so the reader can compare the three computations
directly.

Third, the toric section is organized around the three shapes a
Monge--Amp\`ere measure can take.  \Cref{ex:ch19-toric-two-atoms} is
atomic, \Cref{ex:ch19-toric-segment} is carried by a segment, and
\Cref{ex:ch19-toric-p2-smooth} is absolutely continuous with the explicit
density $\varepsilon^{-2}p_0p_1p_2$.  The last of these is written out
because it makes the factor $n!$ and the total mass $\vol_M(P)$ visible
in a single change of variables, and because the same family reappears as
the toric Fubini--Study degeneration in
\Cref{sec:metrics-fs}.

Fourth, the canonical metric on a Tate curve is computed from scratch in
\Cref{ex:ch19-tate-canonical}, including the functional equation
\eqref{eq:ch19-tate-cocycle}, its quadratic solution
\eqref{eq:ch19-tate-g}, the identification
$q_L(\ell)=d\ell$ with the Poincar\'e pairing, and the exact
piecewise affine interpolation \eqref{eq:ch19-tate-model-approx} that
recovers the balanced model metric of
\Cref{chap:model-calculations}.  The general theory of
\Cref{sec:ch19-canonical} is stated so that every step of this example is an
instance of it.

Fifth, \Cref{sec:ch19-limitations} is a device rather than an exposition.
Its fourteen rows are the assertions that later chapters are most tempted to
apply outside their range, and the three-way classification into \emph{not
known}, \emph{not defined}, and \emph{different} is used consistently from
here on.

\paragraph{A note on the delta-metricity question.}
\index{abelian variety!canonical metric!scope of results}
The statement that a canonical metric on an abelian variety is locally the
tensor product of a formal metric and a smooth metric is correct for good
reduction, total degeneration, and every isotrivial residual Raynaud
extension; a residue field algebraic over a finite field is a useful special
case.  Each has a one-line reason in
\Cref{rem:ch19-canonical-delta-scope}.  In the good-reduction case the
canonical metric \emph{is} a formal metric
(\Cref{thm:ch19-cubical-exists}\,\cref{itm:ch19-cub:goodred}); in the totally
degenerate case the canonical tropicalization is the tropicalization of a
split torus and the metric is CLD-smooth on the skeleton; and over a residue
field algebraic over a finite field, $\Pic^0$ of the residual abelian
variety is torsion, so the residual extension is automatically isotrivial
and the canonical tropicalization is smooth.  There is also a logically
independent settled case: by \Cref{thm:p09p2-reduction}, every cubical metric
on a line bundle in $\Pic^0(A)$ is an ordinary delta-metric.  The same
theorem reduces the universal question on a fixed abelian variety to ample
symmetric cubical metrics.  The product--Gauss criterion
\Cref{thm:p09p2-product-gauss-candidate} closes precisely this reduced
problem: in the non-isotrivial mixed-reduction range the quadratic part is
nondegenerate, and ordinary delta-metricity would force the nonsmooth
canonical tropicalization to be smooth, a contradiction.  Hence the answer
is negative.  The present chapter nevertheless proves the abelian measure computations
using formal curvature powers and harmonic-chart $G$-forms in the precise
scope of \Cref{thm:bt-harmonic-compatibility}; ordinary delta-metricity is
not needed.

\paragraph{Historical note.}
Potential theory on Berkovich curves, in the form used in
\Cref{sec:ch19-curves,sec:ch19-wanner}, was built by Thuillier
\cite{Thu05}; for the projective line the parallel development is Baker and
Rumely's, whose monograph is named here rather than cited by key because no
proof of this chapter takes an input from it.  The two statements of
\Cref{sec:ch19-thuillier-currents} are quoted from \cite{Thu05} with his
numbering, and they are the only inputs of this chapter taken from him that
are not also proved in \Cref{chap:model-calculations}.

Thuillier's own account of how his theory sits relative to its neighbours is
\cite[\S\S5.1--5.4]{Thu05}, and it is worth summarizing because the four
comparisons are of four different kinds.
\begin{itemize}
\item \emph{Rumely.}  Rumely's capacity theory on algebraic curves works with
the rigid points $X_1=X(k)$ and with lemniscate domains
$\{|f|\leq1\}$; \cite[\S5.1]{Thu05} shows that these are the traces on
$X_1$ of strictly affinoid domains, that Rumely's algebraic capacity is
Thuillier's Choquet capacity of the Berkovich closure
\cite[Thm.~5.1.2]{Thu05}, and---the point relevant to us---that Rumely's
smooth functions and Laplacian on a reduction graph differ from Thuillier's
by a single constant, $\ddc=(\log p)^{-1}\ddc_{\mathrm R}$, because the
two normalize edge lengths differently.  This is the factor discussed in
\Cref{rem:ch08-thuillier-normalization}; it never appears in this book.
\item \emph{Baker--Rumely.}  Their Hsia kernel $\delta_\infty$ on
$\mathbf P^{1,\An}$ is recovered as $\rme^{-g}$, $g$ being the
symmetric continuous Green function attached by
\cite[Thm.~5.1.5]{Thu05} to the point $\infty$; see
\cite[\S5.2]{Thu05}.  Their theory is thus the projective-line specialization
of his, with a kernel constructed by nested discs rather than by solving
$\ddc T=\delta_x-\delta_\infty$.
\item \emph{Kani.}  Kani's potential-theoretic reading of intersection theory
on an arithmetic surface is what \cite{Thu05} identifies as the origin of the
whole subject; \cite[\S5.3]{Thu05} explains that Kani's ``local potential
theories'' are the linear algebra of harmonic functions restricted to
skeleta, and that his admissible measures are the Dirac measures at points of
$\mathrm I(X)$.
\item \emph{Favre--Jonsson.}  The valuative tree of Favre and Jonsson carries
a Laplacian $\Delta$ defined by a choice of root $t_0$;
\cite[Prop.~5.4.2]{Thu05} shows that $\ddc g=g(t_0)\delta_{t_0}-\Delta g$
for every complex potential $g$, so that the root dependence is exactly the
difference between a Laplacian and a Laplacian normalized to kill a marked
point, and \cite[\S5.4]{Thu05} then removes the root by sheafifying over
distinguished open subsets of a metrized real graph.  This is the same
bookkeeping that reappears in \Cref{part:bj} as the choice of base point on
the trivially valued cone.
\end{itemize}

Chambert-Loir and Ducros come after all of this and in a different spirit:
they construct a calculus of forms and currents in every dimension, at the
cost of a smooth sheaf that, already on curves, is too small to contain the
potentials with atomic curvature (\Cref{rem:ch19-thuillier-smooth-vs-cld}).
The identification of the two notions of subharmonicity on curves---
comparison against harmonic functions, and positivity of a
Chambert--Loir--Ducros current---is Wanner's \cite{Wanner19}, and it remains
the basic comparison between these two current-theoretic notions.  The
additional identification of the continuous current-positive class with the
first ordinary-smooth decreasing-limit stage is
\Cref{thm:p063-current-psh-ordinary-regularization}.
What
\Cref{thm:ch19-thuillier-sh-currents} adds is that on the Thuillier side the
identification of subharmonic functions with positive-$\ddc$ currents was
already unconditional, for arbitrary upper semicontinuous potentials and over
an arbitrary base field.

The toric picture has three independent origins.  Arithmetic toric geometry,
in which invariant metrics correspond to concave functions and heights to
integrals over polytopes, is due to Burgos Gil, Philippon and Sombra
\cite{BGPS14}, with the underlying toric geometry in \cite{Fultor,CLS11} and
the convexity in \cite{Roc70,Sch14}.  The model-theoretic measure is
Chambert-Loir's \cite{CL06}.  The systematic invariant tropical
pluripotential theory, including singular potentials and the solution of
invariant Monge--Amp\`ere equations, is
Burgos Gil--Gubler--Jell--K\"unnemann--Kramer \cite{BGJKK25}; Vilsmeier's
polystable-face comparison \cite{Vil20} is a different, local statement
carrying its own calibration factor and is used in this book only in that
form.

The abelian theory begins with Tate's uniformization of elliptic curves with
multiplicative reduction, extended by Raynaud to abeloid varieties and
written out in detail by Bosch and L\"utkebohmert
\cite{BoschLutkebohmert1993}; Berkovich's deformation retraction
\cite{Berk12} supplies the skeleton.  Cubical metrics are the
non-Archimedean form of the theorem of the cube, and their curvature
currents were computed by Chambert-Loir and Ducros
\cite[Chapter~19]{CLD25}; the structural description
\eqref{eq:ch19-canonical-weight}, splitting the canonical metric into a
formal factor and a quadratic tropical factor, is theirs, as is the
observation that the canonical tropicalization is smooth exactly when the
residual extension is isotrivial.  Gubler's earlier work on canonical
measures on abelian varieties supplied the polyhedral decomposition recorded
in \Cref{rem:ch19-canonical-pieces} and the example showing that every piece
can be nonzero.  Gubler, Jell and Rabinoff \cite{GJR26} then isolated
harmonic tropicalizations as the right enlargement of the chart category,
which is what makes the Bedford--Taylor computations of
\Cref{sec:ch19-canonical} legitimate without any smoothness assumption on
$\tau$.

Finally, the comparison of \Cref{sec:ch19-cld-bfj} has a longer prehistory
than its statement suggests.  Zhang's semipositive metrics \cite{Zhang1995}
and Chambert-Loir's measures \cite{CL06} already contained the model half;
what the Chambert-Loir--Ducros calculus \cite{CLD12,CLD25} adds is that the
same measure is a wedge product of currents, and what
Boucksom--Favre--Jonsson \cite{BFJ15,BFJ16} add is a global variational
theory that extends the operator far beyond the continuous locus.  The
theorem proved here is the precise statement of where those two extensions
overlap, and its part~(C) is the list of places where they do not.  Later
comparisons, in particular those of \cite{BE21} under semiample and CLD
hypotheses, enlarge the overlap but do not close it; see
\Cref{chap:compare-ma}.
